% Snapshot hashes (paper statements and Lean excerpts are embedded):
% STOC2027.tex SHA-256 6a2a92e1e2bf7c2b257ecf39f52a1181d490f36f8da947ab251a09cfec0e8a3b
% DPCut/AdaptiveComposition.lean SHA-256 41491f37471f802e18b616bfc8159d87da2ce814e1bd63f25ee62dfed96eccf7
% DPCut/Algorithms.lean SHA-256 3bbbc30ebcea06242debe684af99c95ba3b6dc2a5ed9c722af8deb908febb207
% DPCut/Applications.lean SHA-256 5c0a47c5f21fdd5832bd2ea312a91ab60561643f764ba0aa15d3d12fae78d086
% DPCut/ApproxSDPPrecision.lean SHA-256 ab05c8f093985f47a0aa2fbb27698e3ac7ee661792f45e7f5d25bf10ebcbd862
% DPCut/ApproxSDPRelease.lean SHA-256 9c71b8751308e7e491321e551d6f9336babb4dc0dc2417a9fc7804a298237744
% DPCut/ApproximationTransfer.lean SHA-256 d8fed844f7d28a5f50f8004d7d4ddd14cac128ff3c6063b0d75bcb33ce6c22ee
% DPCut/Arithmetic.lean SHA-256 4f01accd1d29b1d436045b70177b9bdb52e45eecc2822c0abbbfaa6663614875
% DPCut/CCApplication.lean SHA-256 a4c60738d433d40b7f2ae02a9eb399c693a67699e41aedf959a306f8dd2e6a9a
% DPCut/CappedLSWRuntime.lean SHA-256 11993c48d2d3d482375147e84aac96a4b6a4c391fb68416a27cdd89cdaa4428c
% DPCut/CostedProgram.lean SHA-256 84d0d160b697a3b74c3dd7d0cf26f901cdd8eee0ff9d5113f568c6d12b863a7b
% DPCut/Definitions.lean SHA-256 16bfd53531f6d75717c874db65c2d7b88c47dbcbe937ad0602a4bfa7d1f43735
% DPCut/ExternalApplicationContracts.lean SHA-256 09d973c5cc4c728959bfd0a85fa8d459a5c86721c4cee66e54920b7f48e0dcce
% DPCut/ImplementedApplications.lean SHA-256 3116047f2e04de82cd5810a876b28be55f43d16ce767603c6715a735dba016d2
% DPCut/LSWSolverFamily.lean SHA-256 bb450706c2cd5e637d6d95cbbe10a833243222e4e33bfc3031453d08bde829d9
% DPCut/LSWStatement.lean SHA-256 206ee470873f09d8cd66d3270eff3b47819ce5477bffa11aed9f345f645cc2e7
% DPCut/Lemma10.lean SHA-256 59ce2bb86a7e308b3e621dc340de92d62ecadab3268c07e366ee6b633d5ae7b2
% DPCut/Lemma11.lean SHA-256 09fc200e3bbf028ff90a6c1a77bc8212e5d7b34058a62f6b6a54f87fc49c1cee
% DPCut/Lemma12.lean SHA-256 1bde04dd674005eb2a75d0a5b7a2ed0c6e7f985d1b60ba4a2aa77a4f76bbb002
% DPCut/Lemma13.lean SHA-256 51737f672ce5ee34cf6b57721e7d36de297ab4dfc245adb918c61f26140d5a3d
% DPCut/Lemma14.lean SHA-256 9259cab1b694eb97195a54e500e0ca1b601be9e882c1dcc3d033474168c1e0cd
% DPCut/Lemma15.lean SHA-256 ccae3c648ec27a1e0caf9377acc0a9efa939731c8c9ae571e231107f9b52e06c
% DPCut/Lemma16.lean SHA-256 5973e067101b5e2dc15db4d6a41b972028fdbee5e4546f23828cb7e2da90ac07
% DPCut/Lemma17.lean SHA-256 12899d9d73db4f27fbb2ec2bb092e315f85d97249f5366ddbf07e6a2c9db4be5
% DPCut/Lemma18.lean SHA-256 908ac40d7695096b2f5e3a6cb0dfdab4d2b51fa233afa90be5e3e9fc59574a6c
% DPCut/Lemma19.lean SHA-256 4ff1943ee93645562ed0c9e6228bcf1655868698656fee93f1255305b4c4253e
% DPCut/Lemma20.lean SHA-256 9e56f8793a48cf9c9feb7a1e47e722f511d5b2fda5f160202ea85daa21eb817f
% DPCut/Lemma21.lean SHA-256 6ee223a1483e69518fdc2768a7507fcd04b18e1fafeee86bf4bf7c8acf56cfce
% DPCut/Lemma22.lean SHA-256 53b1ecf7bea247749a10e272723d01a44f406af5acc860c85b7c2810d9276007
% DPCut/Lemma23.lean SHA-256 c5db6d54b4574ed8309aab27768f16a191525b67e663179e1ba69ff9420e94ff
% DPCut/Lemma29.lean SHA-256 161dfc1f67014ff8e7136e4ef4a3bed80858222f78c4857604433f0201366b05
% DPCut/Lemma35.lean SHA-256 bae48c02eee55b6988e496881ea7637fe988434f022822178c305a5643793a66
% DPCut/OuterUtilityCalibration.lean SHA-256 c3621544661c7d8c2fca59c7a400684e85d938b85564a5f91bd757e7d1eca66d
% DPCut/Privacy.lean SHA-256 1c0a37f6b9b72c34539befeaa5c18fa54e9eb0178c59f978fa601a735eaaa5aa
% DPCut/PublicOracleCalibration.lean SHA-256 c9fcc955bdfe17e0716d476619ecb745b93bb7f790ba5321f6ae0b097d25a1d8
% DPCut/Remark33.lean SHA-256 ed6fbf67152bb47634bebb03c92d17706f573ae8f3c10aeef7904baaef7bd07c
% DPCut/SDPLinearOracle.lean SHA-256 7e1f4f0961b6273f1f2552af708a950f8d895d4007ed8df599438a2c14344c59
% DPCut/SDPRuntimeData.lean SHA-256 ea4501f344565071efc6168d4f3e9f10b52e339f9741f5b17742df62fefe2447
% DPCut/Statements.lean SHA-256 c8d02eb343a68f3493e2f0b17fcd94b5fecc3d5daaa3a56ea0b0a489bb1608b4
% DPCut/Theorem1.lean SHA-256 11b34f029aff0b3a0f77a43f2478c51f6d54c0d3d63a42adc21434e230558008
% DPCut/Theorem1LSW.lean SHA-256 6b82f50f1f84808a9d7868799ade993d13d6ed9cfa55f31cee4a97f304cabe6a
% DPCut/Theorem30.lean SHA-256 25ca5d8ca6a5fa639f7ca059752c9f4dee34a339ea56c06861293088813eb083
% DPCut/Theorem38.lean SHA-256 b9f4a4a044521169d2c5ead65c8931c21ce8823a8d0f2d34592004c58a9c1866
% DPCut/TimedAlgorithm.lean SHA-256 713fed499fbb5e98ba0bb3f8abffb3548cdbbc91a626bdffb544ecbcb49e031d
% DPCut/Utility.lean SHA-256 3b6ecb86b96d05aa8e3b0b04b898c44df1d1a83f23208e8ebd981ae108fd5dcd
% DPCut/UtilityCalibration.lean SHA-256 b85b60973063d27bb79c388de80a94c941ff5fa1804da1d54b841e7ee88a00a9
% DPCut/VerifiedLinearAlgebra.lean SHA-256 efe2d88c189b819a8f7fd3e49b4b78c5aceaf2d943db63762c88f7a07a7d1bab
% Generated by make_lemma_comparison.py; edit the sources, then regenerate.
% Compile from formalization/ with: pdflatex LemmaComparison.tex (twice).
\documentclass[a4paper,10pt]{article}
\usepackage[T1]{fontenc}
\usepackage[utf8]{inputenc}
\usepackage[margin=23mm]{geometry}
\usepackage{amsmath,amssymb,amsthm}
\usepackage{thm-restate}
\usepackage{xcolor,listings,textcomp}
\usepackage{xparse}
\usepackage[colorlinks=true,linkcolor=blue,urlcolor=blue]{hyperref}
\setcounter{secnumdepth}{0}
\setcounter{tocdepth}{1}
\setlength{\parindent}{0pt}
\setlength{\parskip}{5pt}
\setlength{\emergencystretch}{3em}
\allowdisplaybreaks
\newtheorem{theorem}{Theorem}
\newtheorem{lemma}[theorem]{Lemma}
\newtheorem{proposition}[theorem]{Proposition}
\newtheorem{corollary}[theorem]{Corollary}
\theoremstyle{definition}
\newtheorem{definition}[theorem]{Definition}
\theoremstyle{remark}
\newtheorem{remark}[theorem]{Remark}
\newcommand{\R}{\mathbb{R}}
\newcommand{\E}{\mathbb{E}}
\newcommand{\eps}{\varepsilon}
\newcommand{\del}{\delta}
\newcommand{\one}{\mathbf{1}}
\newcommand{\ind}{\mathbb{1}}
\newcommand{\Prb}{\mathbf{Pr}}
\newcommand{\cM}{\mathcal{M}}
\newcommand{\pdf}{\mathrm{pdf}}
\newcommand{\KL}{D_{\mathrm{KL}}}
\newcommand{\Lap}{\operatorname{Lap}}
\newcommand{\bottom}{\bot}
\newcommand{\norm}[1]{\left\Vert#1\right\Vert}
\newcommand{\ip}[2]{\left\langle#1,#2\right\rangle}
\newcommand{\odotc}{\mathbin{\odot}}
\DeclareMathOperator*{\argmin}{arg\,min}
\DeclareMathOperator*{\argmax}{arg\,max}
\renewcommand{\epsilon}{\varepsilon}
\renewcommand{\emptyset}{\varnothing}
\newcommand{\SourceId}[1]{\texttt{\detokenize{#1}}}
% Paper references retain paper numbering where known. Equation/section
% references and citations use stable source keys, avoiding external .aux/.bib files.
\newcommand{\PaperRef}[1]{\ifcsname PaperNumber@#1\endcsname\csname PaperNumber@#1\endcsname\else\SourceId{#1}\fi}
\newcommand{\PaperCref}[1]{\ifcsname PaperKind@#1\endcsname\csname PaperKind@#1\endcsname~\PaperRef{#1}\else\SourceId{#1}\fi}
\RenewDocumentCommand\ref{m}{\PaperRef{#1}}
\RenewDocumentCommand\eqref{m}{\textup{(\PaperRef{#1})}}
\NewDocumentCommand\Cref{m}{\PaperCref{#1}}
\NewDocumentCommand\cref{m}{\PaperCref{#1}}
\RenewDocumentCommand\cite{O{} m}{[\if\relax\detokenize{#1}\relax\else#1; \fi\SourceId{#2}]}
\lstdefinestyle{lean}{
  basicstyle=\ttfamily\footnotesize, columns=fullflexible, keepspaces=true,
  breaklines=true, breakatwhitespace=false, showstringspaces=false,
  numbers=left, numberstyle=\tiny\color{gray}, numbersep=7pt,
  xleftmargin=2em, frame=single, framerule=0.2pt, rulecolor=\color{gray!40},
  backgroundcolor=\color{gray!3}, aboveskip=7pt, belowskip=9pt,
  tabsize=2, upquote=true,
  morekeywords={def,abbrev,structure,inductive,where,Prop,Type,fun,let,if,then,else,match,with,theorem,lemma},
  keywordstyle=\color{blue!60!black},
  literate=
    {¹}{{${}^{1}$}}1
    {×}{{$\times$}}1
    {Δ}{{$\Delta$}}1
    {Θ}{{$\Theta$}}1
    {Ω}{{$\Omega$}}1
    {α}{{$\alpha$}}1
    {β}{{$\beta$}}1
    {γ}{{$\gamma$}}1
    {δ}{{$\delta$}}1
    {ε}{{$\varepsilon$}}1
    {η}{{$\eta$}}1
    {θ}{{$\theta$}}1
    {ι}{{$\iota$}}1
    {κ}{{$\kappa$}}1
    {μ}{{$\mu$}}1
    {ν}{{$\nu$}}1
    {ξ}{{$\xi$}}1
    {π}{{$\pi$}}1
    {ρ}{{$\rho$}}1
    {σ}{{$\sigma$}}1
    {τ}{{$\tau$}}1
    {χ}{{$\chi$}}1
    {ᵐ}{{${}^{m}$}}1
    {ᵥ}{{${}_{v}$}}1
    {‖}{{$\Vert$}}1
    {•}{{$\bullet$}}1
    {⁻}{{${}^{-}$}}1
    {₀}{{${}_{0}$}}1
    {₁}{{${}_{1}$}}1
    {₂}{{${}_{2}$}}1
    {₊}{{${}_{+}$}}1
    {ℕ}{{$\mathbb{N}$}}1
    {ℝ}{{$\mathbb{R}$}}1
    {→}{{$\to$}}1
    {↔}{{$\leftrightarrow$}}1
    {∀}{{$\forall$}}1
    {∂}{{$\partial$}}1
    {∃}{{$\exists$}}1
    {∅}{{$\varnothing$}}1
    {∈}{{$\in$}}1
    {∉}{{$\notin$}}1
    {∑}{{$\sum$}}1
    {∞}{{$\infty$}}1
    {∧}{{$\wedge$}}1
    {∨}{{$\vee$}}1
    {∪}{{$\cup$}}1
    {∫}{{$\int$}}1
    {≠}{{$\neq$}}1
    {≤}{{$\leq$}}1
    {≥}{{$\geq$}}1
    {≪}{{$\ll$}}1
    {⊆}{{$\subseteq$}}1
    {⊕}{{$\oplus$}}1
    {⌈}{{$\lceil$}}1
    {⌉}{{$\rceil$}}1
    {⟨}{{$\langle$}}1
    {⟩}{{$\rangle$}}1
    {⟪}{{$\langle\!\langle$}}1
    {⟫}{{$\rangle\!\rangle$}}1
    {⬝}{{$\cdot$}}1
}
\title{Paper--Lean statement comparison}
\author{}
\date{Source snapshot: 6a2a92e1e2bf7c2b}
\begin{document}
\maketitle
This document pairs every active theorem, corollary, lemma, proposition,
numbered definition, and remark in \SourceId{STOC2027.tex} with its existing Lean counterpart.
The paper statements are extracted from the current manuscript, excluding
comments and disabled material. The Lean definitions are copied from the current
\SourceId{DPCut/} files, not paraphrased. Source locations accompany both versions.

Each \texttt{def ...Statement : Prop} defines a proposition; it is not itself a
proof. Where the project has a corresponding numbered theorem, its declaration
is also shown, with the proof body omitted. Absence of such a declaration here
is not a claim that no related result has been proved. This document compares
statements; it does not re-run or certify Lean proofs.

Cited and standard background lemmas are marked below and included for statement
comparison. Proving their full generality within this project is not required
to validate this paper's contributions. External hypotheses used by the
formalized proofs remain explicit. Cost conclusions bound assigned counters and
reported times. The current interfaces do not derive these from execution steps:
in particular, \texttt{TimedAlgorithm} permits any measurable function at time
zero. The formal LSW and application contracts therefore do not certify actual
polynomial running time, even in the real-arithmetic model.

Paper statement and algorithm references retain the manuscript's numbering.
Equation and section references are shown by their source labels, and citations
by their bibliography keys. Thus this file needs neither the paper's auxiliary
files nor its bibliography. Lean listings retain Unicode in the source; the
typeset glyphs are supplied by \texttt{listings}. Line numbers in listings refer
to the original Lean files.

Unless a source annotation specifies otherwise, displayed Lean declarations
are in namespace \texttt{DPCut}, using
\texttt{MeasureTheory}, \texttt{ProbabilityTheory}, and the usual
\texttt{BigOperators}/\texttt{ENNReal}/\texttt{MatrixOrder} scopes.
The LSW excerpts also use \texttt{Metric} and the \texttt{RealInnerProductSpace}
scope. Excerpts retain surrounding section variables from their source files;
in particular, \texttt{TimedAlgorithm.expectedCost} has a parameter
\texttt{A : TimedAlgorithm} $\Theta\ \alpha$.
The notation appendix expands the main predicates, parameter records, and
algorithm laws used in the targets.
\tableofcontents

\expandafter\def\csname PaperNumber@thm:main\endcsname{1}
\expandafter\def\csname PaperKind@thm:main\endcsname{Theorem}
\expandafter\def\csname PaperNumber@cor:cc-intro\endcsname{2}
\expandafter\def\csname PaperKind@cor:cc-intro\endcsname{Corollary}
\expandafter\def\csname PaperNumber@cor:maxcut-intro\endcsname{3}
\expandafter\def\csname PaperKind@cor:maxcut-intro\endcsname{Corollary}
\expandafter\def\csname PaperNumber@def:DP\endcsname{4}
\expandafter\def\csname PaperKind@def:DP\endcsname{Definition}
\expandafter\def\csname PaperNumber@def:laplace-dist\endcsname{5}
\expandafter\def\csname PaperKind@def:laplace-dist\endcsname{Definition}
\expandafter\def\csname PaperNumber@lem:laplace-mech\endcsname{6}
\expandafter\def\csname PaperKind@lem:laplace-mech\endcsname{Lemma}
\expandafter\def\csname PaperNumber@lem:post-processing\endcsname{7}
\expandafter\def\csname PaperKind@lem:post-processing\endcsname{Lemma}
\expandafter\def\csname PaperNumber@lem:basic-composition\endcsname{8}
\expandafter\def\csname PaperKind@lem:basic-composition\endcsname{Lemma}
\expandafter\def\csname PaperNumber@lem:advanced-composition\endcsname{9}
\expandafter\def\csname PaperKind@lem:advanced-composition\endcsname{Lemma}
\expandafter\def\csname PaperNumber@lem:chi\endcsname{10}
\expandafter\def\csname PaperKind@lem:chi\endcsname{Lemma}
\expandafter\def\csname PaperNumber@lem:chaos\endcsname{11}
\expandafter\def\csname PaperKind@lem:chaos\endcsname{Lemma}
\expandafter\def\csname PaperNumber@lem:preprocess\endcsname{12}
\expandafter\def\csname PaperKind@lem:preprocess\endcsname{Lemma}
\expandafter\def\csname PaperNumber@lem:stability_of_optimizer\endcsname{13}
\expandafter\def\csname PaperKind@lem:stability_of_optimizer\endcsname{Lemma}
\expandafter\def\csname PaperNumber@lem:gaussian-privacy\endcsname{14}
\expandafter\def\csname PaperKind@lem:gaussian-privacy\endcsname{Lemma}
\expandafter\def\csname PaperNumber@lem:concentration_chi-squares\endcsname{15}
\expandafter\def\csname PaperKind@lem:concentration_chi-squares\endcsname{Lemma}
\expandafter\def\csname PaperNumber@lem:pay-what-you-get\endcsname{16}
\expandafter\def\csname PaperKind@lem:pay-what-you-get\endcsname{Lemma}
\expandafter\def\csname PaperNumber@lem:gradient_privacy\endcsname{17}
\expandafter\def\csname PaperKind@lem:gradient_privacy\endcsname{Lemma}
\expandafter\def\csname PaperNumber@lem:main-privacy\endcsname{18}
\expandafter\def\csname PaperKind@lem:main-privacy\endcsname{Lemma}
\expandafter\def\csname PaperNumber@lem:sdp_error\endcsname{19}
\expandafter\def\csname PaperKind@lem:sdp_error\endcsname{Lemma}
\expandafter\def\csname PaperNumber@lem:gradient_estimation\endcsname{20}
\expandafter\def\csname PaperKind@lem:gradient_estimation\endcsname{Lemma}
\expandafter\def\csname PaperNumber@lem:inner_round\endcsname{21}
\expandafter\def\csname PaperKind@lem:inner_round\endcsname{Lemma}
\expandafter\def\csname PaperNumber@lem:outer_round\endcsname{22}
\expandafter\def\csname PaperKind@lem:outer_round\endcsname{Lemma}
\expandafter\def\csname PaperNumber@lem:utility\endcsname{23}
\expandafter\def\csname PaperKind@lem:utility\endcsname{Lemma}
\expandafter\def\csname PaperNumber@lem:chi_square_tail\endcsname{24}
\expandafter\def\csname PaperKind@lem:chi_square_tail\endcsname{Lemma}
\expandafter\def\csname PaperNumber@lem:hypercontractivity\endcsname{25}
\expandafter\def\csname PaperKind@lem:hypercontractivity\endcsname{Lemma}
\expandafter\def\csname PaperNumber@lem:holder_interpolation\endcsname{26}
\expandafter\def\csname PaperKind@lem:holder_interpolation\endcsname{Lemma}
\expandafter\def\csname PaperNumber@prop:unbiased_estimator\endcsname{27}
\expandafter\def\csname PaperKind@prop:unbiased_estimator\endcsname{Proposition}
\expandafter\def\csname PaperNumber@lem:pythagorean\endcsname{28}
\expandafter\def\csname PaperKind@lem:pythagorean\endcsname{Lemma}
\expandafter\def\csname PaperNumber@lem:capped-unit-release\endcsname{29}
\expandafter\def\csname PaperKind@lem:capped-unit-release\endcsname{Lemma}
\expandafter\def\csname PaperNumber@thm:cc-black-box-reduction\endcsname{30}
\expandafter\def\csname PaperKind@thm:cc-black-box-reduction\endcsname{Theorem}
\expandafter\def\csname PaperNumber@cor:cc-fixed-k\endcsname{31}
\expandafter\def\csname PaperKind@cor:cc-fixed-k\endcsname{Corollary}
\expandafter\def\csname PaperNumber@thm:private-max-k-cut\endcsname{32}
\expandafter\def\csname PaperKind@thm:private-max-k-cut\endcsname{Theorem}
\expandafter\def\csname PaperNumber@rem:maxcut-tightness\endcsname{33}
\expandafter\def\csname PaperKind@rem:maxcut-tightness\endcsname{Remark}
\expandafter\def\csname PaperNumber@thm:private-constrained-cut\endcsname{34}
\expandafter\def\csname PaperKind@thm:private-constrained-cut\endcsname{Theorem}
\expandafter\def\csname PaperNumber@lem:hc-bicriteria-recursion\endcsname{35}
\expandafter\def\csname PaperKind@lem:hc-bicriteria-recursion\endcsname{Lemma}
\expandafter\def\csname PaperNumber@thm:private-hierarchical-clustering\endcsname{36}
\expandafter\def\csname PaperKind@thm:private-hierarchical-clustering\endcsname{Theorem}
\expandafter\def\csname PaperNumber@lem:ekkl-discrepancy\endcsname{37}
\expandafter\def\csname PaperKind@lem:ekkl-discrepancy\endcsname{Lemma}
\expandafter\def\csname PaperNumber@thm:lower_bound\endcsname{38}
\expandafter\def\csname PaperKind@thm:lower_bound\endcsname{Theorem}
\expandafter\def\csname PaperNumber@alg:main\endcsname{1}
\expandafter\def\csname PaperKind@alg:main\endcsname{Algorithm}
\expandafter\def\csname PaperNumber@alg:gradient_estimation\endcsname{2}
\expandafter\def\csname PaperKind@alg:gradient_estimation\endcsname{Algorithm}
\expandafter\def\csname PaperNumber@alg:gradient_aggregation\endcsname{3}
\expandafter\def\csname PaperKind@alg:gradient_aggregation\endcsname{Algorithm}
\expandafter\def\csname PaperNumber@alg:preprocess\endcsname{4}
\expandafter\def\csname PaperKind@alg:preprocess\endcsname{Algorithm}
\clearpage
\section{Theorem 1: Private relative cut release}
\SourceId{thm:main}\hfill\SourceId{STOC2027.tex}:336--350
\paragraph{Paper statement.}
\setcounter{theorem}{0}
\begin{theorem}\label{thm:main}
Let
$n\geq2$, $0<\gamma\leq 1/4$, $0<\varepsilon\leq 1$, and
$0<\delta,\beta<1/10$.  There is an $(\varepsilon, \delta)$ edge-DP algorithm that, on every
nonnegative weighted $n$-vertex graph $G$, runs in time polynomial in
$n$, $1/\gamma$, $1/\varepsilon$, $\log(1/\delta)$, and $\log(1/\beta)$, and outputs a
nonnegative weighted graph $\widetilde{G}$ satisfying, with probability at least $1-\beta$, simultaneously
for every $S\subseteq V$,
\begin{align*}
 \left|w_G(S,V\setminus S)-w_{\widetilde{G}}(S,V\setminus S)\right|
 \leq \gamma w_G(S,V\setminus S)  + O\left(\frac{n}{\gamma\varepsilon}
 \left(\log(n/(\delta\beta\gamma\eps))\right)^{7}\right) .
\end{align*}
\end{theorem}
\paragraph{Lean statement.}
\textit{Source:} \SourceId{DPCut/Statements.lean}, line 178.
\begin{lstlisting}[style=lean,firstnumber=178]
def MainStatisticalStatement : Prop :=
  ∃ C : ℝ, 0 < C ∧ ∀ (n : ℕ), 2 ≤ n → ∀ p : Parameters, p.Valid →
    ∃ M : Kernel (Weights n) (Weights n), IsMarkovKernel M ∧
      GraphValued M ∧ EdgeDP M p.epsilon p.delta ∧
      ∀ w : Weights n, Nonnegative w →
        HighProbability (M w) p.beta (fun h => AllCuts w h p.gamma (C * mainError n p))
\end{lstlisting}

\textit{Source:} \SourceId{DPCut/Statements.lean}, line 202.
\begin{lstlisting}[style=lean,firstnumber=202]
def MainTheoremStatement (model : ComputationalModel) : Prop :=
  ∃ C Ctime : ℝ, ∃ degree : ℕ, ∃ program : model.Program, 0 < C ∧ 0 < Ctime ∧
  ∀ (n : ℕ), 2 ≤ n → ∀ p : Parameters, p.Valid →
    ∃ M : Kernel (Weights n) (Weights n),
      model.realizes n program p M ∧ IsMarkovKernel M ∧ GraphValued M ∧
      EdgeDP M p.epsilon p.delta ∧ ∀ w : Weights n, Nonnegative w →
        model.cost n program p w ≤ ENNReal.ofReal (Ctime * runtimeScale n p ^ degree) ∧
        HighProbability (M w) p.beta (fun h => AllCuts w h p.gamma (C * mainError n p))
\end{lstlisting}

\paragraph{Corresponding theorem declaration (proof omitted).}
\textit{Source:} \SourceId{DPCut/Theorem1.lean}, line 17.
\begin{lstlisting}[style=lean,firstnumber=17]
theorem theorem1_statistical {n : ℕ} (hn : 2 ≤ n) {p : Parameters} (hp : p.Valid) :
    Measurable (fullReleaseLaw (n := n) p outerUtilityConstants) ∧
      (∀ w : Weights n, IsProbabilityMeasure (fullReleaseLaw p outerUtilityConstants w) ∧
        ∀ᵐ h ∂fullReleaseLaw p outerUtilityConstants w, Nonnegative h) ∧
      EdgeDP (fullReleaseLaw (n := n) p outerUtilityConstants) p.epsilon p.delta ∧
      ∀ w : Weights n, Nonnegative w →
        HighProbability (fullReleaseLaw p outerUtilityConstants w) p.beta
          (fun h => AllCuts w h p.gamma ((2 : ℝ) ^ 166 * mainError n p))
\end{lstlisting}

\paragraph{Corresponding theorem declaration (proof omitted).}
\textit{Source:} \SourceId{DPCut/Theorem1.lean}, line 30.
\begin{lstlisting}[style=lean,firstnumber=30]
theorem theorem1_statistical_exists : MainStatisticalStatement
\end{lstlisting}

\paragraph{Corresponding theorem declaration (proof omitted).}
\textit{Source:} \SourceId{DPCut/Theorem1LSW.lean}, line 104.
\begin{lstlisting}[style=lean,firstnumber=104]
theorem theorem1_of_lsw (hLSW : LSWStatement) :
    MainTheoremStatement (realArithmeticLSWModel hLSW)
\end{lstlisting}

\paragraph{Comparison note.}
The statistical theorem \texttt{theorem1\_statistical} explicitly names \texttt{fullReleaseLaw}, bundling measurability, normalization, nonnegative output, privacy and all-cut accuracy with constant $2^{166}$. The first target is its existential corollary \texttt{theorem1\_statistical\_exists}. The second target adds a polynomial bound on the supplied model's cost. The current \texttt{theorem1\_of\_lsw} uses \texttt{realArithmeticLSWModel} and the proposition parameter \texttt{LSWStatement}, intended to specialize Lee--Sidford--Wong, Theorem~42 (arXiv:1508.04874v2). Its \texttt{TimedAlgorithm} interface permits arbitrary measurable functions with reported time zero. Thus this is cost accounting, not an operational runtime theorem conditional only on the cited solver. The SDP geometry, public mass clipping, objective and precision bounds, reported-time truncation, restarts and privacy/utility transfer are proved internally. The constructed matrix routines, fallback decoder and covariance Gaussian sampler are connected to the program through semantic refinement and counter bounds. The value/count interfaces themselves do not enforce permitted instructions; scalar sampling and outer-block prices remain model conventions. The parameter ranges, nonnegative output and simultaneous cut guarantee with logarithmic exponent $7$ match the paper. The identity covariance fallback matches the implementation footnote and its failure probability is included in the total-variation budget. Neither actual polynomial runtime nor the footnote's sharper $\widetilde O(n^8)$ bound is certified by the current model. The notation appendix expands the interfaces and certificates. The older exact-SDP contract and the application contracts share the assigned-cost limitation.
\clearpage
\section{Corollary 2: Private correlation clustering}
\SourceId{cor:cc-intro}\hfill\SourceId{STOC2027.tex}:490--498
\paragraph{Paper statement.}
\setcounter{theorem}{1}
\begin{restatable}[Private correlation clustering]{corollary}{privatecorrelationclustering}
\label{cor:cc-intro}
For every constant $\varepsilon>0$ and $\delta=n^{-10}$, there is a
polynomial-time $(\varepsilon, \delta)$-edge-DP algorithm for unweighted complete-graph min-disagree
correlation clustering that, with high probability, outputs a partition with cost at most
\[
1.387\operatorname{OPT}_G+\widetilde O_{\varepsilon}(n).
\]
\end{restatable}
\paragraph{Lean statement.}
\textit{Source:} \SourceId{DPCut/Applications.lean}, line 154.
\begin{lstlisting}[style=lean,firstnumber=154]
def CorrelationClusteringStatement : Prop :=
  ∀ ε : ℝ, 0 < ε → ∃ C Cfail : ℝ, ∃ exponent : ℕ,
    0 < C ∧ 0 < Cfail ∧ PrivateCCGuarantee C Cfail ε (1387 / 1000) exponent id
\end{lstlisting}

\paragraph{Corresponding theorem declaration (proof omitted).}
\textit{Source:} \SourceId{DPCut/Theorem30.lean}, line 28.
\begin{lstlisting}[style=lean,firstnumber=28]
theorem private_correlation_clustering (baseline : CCBaselineFamily)
    (alg : CCApproximationContract (5547 / 4000) 1) : CorrelationClusteringStatement
\end{lstlisting}

\paragraph{Combined implementation theorem (proof omitted).}
\textit{Source:} \SourceId{DPCut/ImplementedApplications.lean}, line 142. \textit{Namespace:} \SourceId{DPCut.LSWApplications}.
\begin{lstlisting}[style=lean,firstnumber=142]
theorem private_correlation_clustering (hLSW : LSWStatement) (baseline : CCBaselineFamily)
    (alg : CCApproximationContract (5547/4000) 1) {ε : ℝ} (hε : 0 < ε) :
    ∃ C Cfail : ℝ, ∃ exponent : ℕ, 0 < C ∧ 0 < Cfail ∧
      ImplementedCCGuarantee hLSW (baseline ε hε).toPolynomialGraphAlgorithm
        alg.toPolynomialGraphAlgorithm C Cfail ε (1387/1000) exponent id
\end{lstlisting}

\paragraph{Comparison note.}
The original Lean target records the statistical guarantee. The combined implementation theorem also provides a polynomial assigned-cost bound for the specified approximate program, using the formal LSW parameter and cited external algorithm contracts. These interfaces allow arbitrary kernels with cost zero; actual running times of the external algorithms and full reduction are not certified. It makes $\widetilde O(n)$ explicit as $C n(\log n)^d$ and high probability as failure at most $C_{\rm fail}/n^2$, with constants chosen before $n$. Here \texttt{id} allows up to $n$ clusters, so the optimum is unrestricted.
\clearpage
\section{Corollary 3: Private Max-Cut}
\SourceId{cor:maxcut-intro}\hfill\SourceId{STOC2027.tex}:510--517
\paragraph{Paper statement.}
\setcounter{theorem}{2}
\begin{restatable}[Private Max-Cut]{corollary}{privatemaxcut}
\label{cor:maxcut-intro}
For every constant $\xi>0, \varepsilon>0$, and $\delta=n^{-10}$, there is a polynomial-time $(\varepsilon, \delta)$-edge-DP algorithm for Max-Cut that, for any non-negative weighted graph $G$, with high probability outputs a bipartition with cut value at least
\[
(0.878-\xi)\operatorname{OPT}_G-\widetilde O_{\xi,\varepsilon}(n).
\]
Here, $\operatorname{OPT}_G$ is the value of the real max-cut in $G$.
\end{restatable}
\paragraph{Lean statement.}
\textit{Source:} \SourceId{DPCut/Applications.lean}, line 194.
\begin{lstlisting}[style=lean,firstnumber=194]
def PrivateMaxCutStatement : Prop :=
  ∀ (ξ ε : ℝ), 0 < ξ → 0 < ε → ∃ C Cfail : ℝ, ∃ exponent : ℕ,
    0 < C ∧ 0 < Cfail ∧ PrivateMaxKGuarantee C Cfail ε (878 / 1000 - ξ) exponent 2
\end{lstlisting}

\paragraph{Combined implementation theorem (proof omitted).}
\textit{Source:} \SourceId{DPCut/ImplementedApplications.lean}, line 256. \textit{Namespace:} \SourceId{DPCut.LSWApplications}.
\begin{lstlisting}[style=lean,firstnumber=256]
theorem private_max_cut (hLSW : LSWStatement) (gw : MaxKCutContract 2 (878/1000))
    {ξ ε : ℝ} (hξ : 0 < ξ) (hε : 0 < ε) :
    ∃ C Cfail : ℝ, ∃ exponent : ℕ, 0 < C ∧ 0 < Cfail ∧
      ImplementedMaxKGuarantee hLSW gw.toPolynomialGraphAlgorithm C Cfail ε (878/1000-ξ) exponent 2
\end{lstlisting}

\paragraph{Comparison note.}
The original Lean target records the statistical guarantee. The combined implementation theorem also provides a polynomial assigned-cost bound for the specified approximate program, using the formal LSW parameter and cited external algorithm contracts. These interfaces allow arbitrary kernels with cost zero; actual running times of the external algorithms and full reduction are not certified. It makes $\widetilde O(n)$ explicit as $C n(\log n)^d$ and high probability as failure at most $C_{\rm fail}/n^2$, with constants chosen before $n$. The ratio is $878/1000-\xi$, with at most two clusters.
\clearpage
\section{Definition 4: Adjacency and edge differential privacy}
\SourceId{def:DP}\hfill\SourceId{STOC2027.tex}:1473--1485
\paragraph{Paper statement.}
\setcounter{theorem}{3}
\begin{definition}
\label{def:DP}
Consider histograms $w,w' \in \R_{\geq0}^{\binom{n}{2}}$.
We say that histograms $w,w'$ are adjacent, written $w\sim w'$, if
$\norm{w-w'}_1\leq1$.
A randomized algorithm $\cM$ is $(\varepsilon, \delta)$ edge-DP if for
all adjacent $w,w'$ and measurable output events $A$,
\[
 \Prb[\cM(w)\in A]\leq e^\varepsilon\Prb[\cM(w')\in A]+\delta,
\]
and symmetrically with $w,w'$ exchanged.
\end{definition}
\paragraph{Lean statement.}
\textit{Source:} \SourceId{DPCut/Definitions.lean}, line 47.
\begin{lstlisting}[style=lean,firstnumber=47]
def Adjacent {n : ℕ} (w w' : Weights n) : Prop := l1 (w - w') ≤ 1
\end{lstlisting}

\textit{Source:} \SourceId{DPCut/Definitions.lean}, line 63.
\begin{lstlisting}[style=lean,firstnumber=63]
def DP {α β : Type*} [MeasurableSpace β] (adj : α → α → Prop)
    (M : α → Measure β) (ε δ : ℝ) : Prop :=
  ∀ x x', adj x x' → ∀ s, MeasurableSet s →
    M x s ≤ ENNReal.ofReal (Real.exp ε) * M x' s + ENNReal.ofReal δ
\end{lstlisting}

\textit{Source:} \SourceId{DPCut/Definitions.lean}, line 67.
\begin{lstlisting}[style=lean,firstnumber=67]
def EdgeDP {n : ℕ} {α : Type*} [MeasurableSpace α]
    (M : Weights n → Measure α) (ε δ : ℝ) : Prop :=
  DP (fun w w' => Nonnegative w ∧ Nonnegative w' ∧ Adjacent w w') M ε δ
\end{lstlisting}

\paragraph{Comparison note.}
\texttt{Adjacent} defines the $\ell_1$ distance condition. \texttt{EdgeDP} restricts both inputs to nonnegative weights. Quantification over ordered adjacent pairs supplies both directions of the privacy inequality.
\clearpage
\section{Definition 5: Laplace distribution}
\SourceId{def:laplace-dist}\hfill\SourceId{STOC2027.tex}:1505--1512
\paragraph{Paper statement.}
\setcounter{theorem}{4}
\begin{definition}[Laplace Distribution]
\label{def:laplace-dist}
The centered Laplace distribution with scale parameter $b > 0$, denoted as $\text{Lap}(b)$, is a continuous probability distribution with the probability density function
\[
    h(z) = \frac{1}{2b} \exp\left(-\frac{|z|}{b}\right), \quad z \in \mathbb{R}.
\]
If $Z \sim \text{Lap}(b)$, then $\Prb[|Z| \ge t] = \exp\left(-\frac{t}{b}\right)$ for $t\geq 0$.
\end{definition}
\paragraph{Lean statement.}
\textit{Source:} \SourceId{DPCut/Definitions.lean}, line 155.
\begin{lstlisting}[style=lean,firstnumber=155]
def laplaceDensity (b x : ℝ) : ℝ := (2 * b)⁻¹ * Real.exp (-|x| / b)
\end{lstlisting}

\textit{Source:} \SourceId{DPCut/Definitions.lean}, line 156.
\begin{lstlisting}[style=lean,firstnumber=156]
def laplace (b : ℝ) : Measure ℝ := volume.withDensity (fun x => ENNReal.ofReal (laplaceDensity b x))
\end{lstlisting}

\textit{Source:} \SourceId{DPCut/Statements.lean}, line 131.
\begin{lstlisting}[style=lean,firstnumber=131]
def LaplaceStatement : Prop :=
  ∀ b : ℝ, 0 < b → IsProbabilityMeasure (laplace b) ∧
    ∀ t : ℝ, 0 ≤ t → prob (laplace b) {x | t ≤ |x|} = Real.exp (-t / b)
\end{lstlisting}

\paragraph{Comparison note.}
\texttt{laplaceDensity} is the real-valued density, and \texttt{laplace} is the measure obtained from it. \texttt{LaplaceStatement} records normalization and the tail identity for $b>0$.
\clearpage
\section{Lemma 6: Laplace mechanism}
\SourceId{lem:laplace-mech}\hfill\SourceId{STOC2027.tex}:1536--1543
\textit{Background result; not a contribution of this paper.}
\paragraph{Paper statement.}
\setcounter{theorem}{5}
\begin{lemma}[Laplace Mechanism]
\label{lem:laplace-mech}
Let $f \colon \mathcal{X} \to \mathbb{R}^k$ be a function with $\ell_1$-sensitivity $\Delta_1 f = \max_{x \sim x'} \|f(x) - f(x')\|_1$. For any $\varepsilon > 0$, the Laplace mechanism $\mathcal{M}_L$ defined by
\[
    \mathcal{M}_L(x) = f(x) + (Z_1, \dots, Z_k),
\]
where $Z_1, \dots, Z_k$ are drawn i.i.d. from $\text{Lap}(\Delta_1 f / \varepsilon)$, satisfies $\varepsilon$-differential privacy.
\end{lemma}
\paragraph{Lean statement.}
\textit{Source:} \SourceId{DPCut/Privacy.lean}, line 31.
\begin{lstlisting}[style=lean,firstnumber=31]
def LaplaceMechanismStatement : Prop :=
  ∀ (X : Type) (adj : X → X → Prop) (k : ℕ) (f : X → Fin k → ℝ) (Δ ε : ℝ),
    0 ≤ Δ → 0 < ε →
    (∀ x x', adj x x' → (∑ i, |f x i - f x' i|) ≤ Δ) →
    DP adj (laplaceMechanism k f Δ ε) ε 0
\end{lstlisting}

\paragraph{Comparison note.}
Lean uses a sensitivity upper bound $\Delta$ and explicitly handles $\Delta=0$ in \texttt{laplaceMechanism}.
\clearpage
\section{Lemma 7: Post-processing}
\SourceId{lem:post-processing}\hfill\SourceId{STOC2027.tex}:1545--1548
\textit{Background result; not a contribution of this paper.}
\paragraph{Paper statement.}
\setcounter{theorem}{6}
\begin{lemma}[Post-Processing]
\label{lem:post-processing}
Let $\mathcal{M} \colon \mathcal{X} \to \mathcal{Y}$ be a randomized mechanism that is $(\varepsilon, \delta)$-differentially private. Let $f \colon \mathcal{Y} \to \mathcal{Z}$ be an arbitrary randomized mapping whose internal randomness is independent of the private input data. Then the composite mechanism $f \circ \mathcal{M} \colon \mathcal{X} \to \mathcal{Z}$ remains $(\varepsilon, \delta)$-differentially private.
\end{lemma}
\paragraph{Lean statement.}
\textit{Source:} \SourceId{DPCut/Privacy.lean}, line 18.
\begin{lstlisting}[style=lean,firstnumber=18]
def PostProcessingStatement : Prop :=
  ∀ (X Y Z : Type) [MeasurableSpace Y] [MeasurableSpace Z],
  ∀ (adj : X → X → Prop) (M : X → Measure Y) (K : Kernel Y Z) (ε δ : ℝ),
    (∀ x, IsProbabilityMeasure (M x)) → IsMarkovKernel K → 0 ≤ ε → 0 ≤ δ →
    DP adj M ε δ → DP adj (fun x => (M x).bind K) ε δ
\end{lstlisting}

\clearpage
\section{Lemma 8: Basic adaptive composition}
\SourceId{lem:basic-composition}\hfill\SourceId{STOC2027.tex}:1550--1553
\textit{Background result; not a contribution of this paper.}
\paragraph{Paper statement.}
\setcounter{theorem}{7}
\begin{lemma}[Basic Adaptive Composition]
\label{lem:basic-composition}
Let $\mathcal{M}_1, \dots, \mathcal{M}_k$ be a sequence of adaptively chosen randomized mechanisms, where for each $i \in [k]$, the mechanism $\mathcal{M}_i$ is $(\varepsilon_i, \delta_i)$-differentially private. Then their adaptive composition is $\left( \sum_{i=1}^k \varepsilon_i, \sum_{i=1}^k \delta_i \right)$-differentially private.
\end{lemma}
\paragraph{Lean statement.}
\textit{Source:} \SourceId{DPCut/Privacy.lean}, line 43.
\begin{lstlisting}[style=lean,firstnumber=43]
def BasicCompositionStatement : Prop :=
  ∀ (X Y : Type) [MeasurableSpace Y] [StandardBorelSpace Y],
  ∀ (adj : X → X → Prop) (P : X → Sequential Y) (k : ℕ) (ε δ : Fin k → ℝ),
    (∀ x t, IsMarkovKernel (P x t)) →
    (∀ t, 0 ≤ ε t ∧ 0 ≤ δ t) →
    (∀ t : Fin k, ∀ h, DP adj (fun x => P x t h) (ε t) (δ t)) →
    DP adj (fun x => sequentialLaw (P x) k) (∑ t, ε t) (∑ t, δ t)
\end{lstlisting}

\clearpage
\section{Lemma 9: Advanced adaptive composition}
\SourceId{lem:advanced-composition}\hfill\SourceId{STOC2027.tex}:1555--1561
\textit{Background result; not a contribution of this paper.}
\paragraph{Paper statement.}
\setcounter{theorem}{8}
\begin{lemma}[Advanced Adaptive Composition~\cite{DRV10,DR14}]
\label{lem:advanced-composition}
Let $\mathcal{M}_1, \dots, \mathcal{M}_k$ be a sequence of adaptively chosen randomized mechanisms, where, conditional on every fixed preceding transcript, each $\mathcal{M}_i$ is $(\varepsilon_0, \delta_0)$-differentially private for some $\varepsilon_0 > 0$ and $\delta_0 \ge 0$. Then for any $\delta'\in(0,1)$, the adaptive composition of these $k$ mechanisms is $(\varepsilon, k\delta_0 + \delta')$-differentially private, where
$$
   \varepsilon = \varepsilon_0 \sqrt{2k \ln(1/\delta')} + k \varepsilon_0 (e^{\varepsilon_0} - 1).
$$
\end{lemma}
\paragraph{Lean statement.}
\textit{Source:} \SourceId{DPCut/Privacy.lean}, line 55.
\begin{lstlisting}[style=lean,firstnumber=55]
def AdvancedCompositionStatement : Prop :=
  ∀ (X Y : Type) [MeasurableSpace Y] [StandardBorelSpace Y],
  ∀ (adj : X → X → Prop) (P : X → Sequential Y) (k : ℕ) (ε₀ δ₀ δ' : ℝ),
    Symmetric adj →
    (∀ x t, IsMarkovKernel (P x t)) →
    0 < ε₀ → 0 ≤ δ₀ → 0 < δ' → δ' < 1 →
    (∀ t < k, ∀ h, DP adj (fun x => P x t h) ε₀ δ₀) →
    DP adj (fun x => sequentialLaw (P x) k) (advancedEpsilon k ε₀ δ') (k * δ₀ + δ')
\end{lstlisting}

\paragraph{Corresponding theorem declaration (proof omitted).}
\textit{Source:} \SourceId{DPCut/AdaptiveComposition.lean}, line 93.
\begin{lstlisting}[style=lean,firstnumber=93]
theorem advanced_composition_statement : AdvancedCompositionStatement
\end{lstlisting}

\paragraph{Comparison note.}
The hypothesis \texttt{Symmetric adj} ensures that each per-step DP guarantee holds in both directions. The paper's edge adjacency already has this property. The displayed theorem proves the full generic target, including $k=0$ and the trivial case $\delta_0\geq1$, without adding hypotheses.
\clearpage
\section{Lemma 10: Coordinatewise Gaussian witness estimate}
\SourceId{lem:chi}\hfill\SourceId{STOC2027.tex}:2018--2022
\paragraph{Paper statement.}
\setcounter{theorem}{9}
\begin{lemma}\label{lem:chi}
For every fixed $X \in \mathcal{X}$, choose independent random vectors $Z^{(j)}\sim \mathcal{N}(0,X)$ for each $j=1,\dots,k$. Set $\overline g(u,v)=\frac1k\sum_{j=1}^k\frac{(Z_u^{(j)}-Z_v^{(j)})^2}{4},$ and $\hat g(u,v)=\min\{\overline{g}(u,v)/2,1\}$. For some universal constant $C > 0$, if
$k\geq C\log(n/(\beta_\chi\del_1))$, then with probability at least $1-\min(\beta_\chi, \del_1)$, for every pair of nodes $e=\{u,v\}$ simultaneously,
\[\frac{q_{e}(X)}{4}\leq \hat{g}(u,v)\leq q_{e}(X).\]
\end{lemma}
\paragraph{Lean statement.}
\textit{Source:} \SourceId{DPCut/Statements.lean}, line 63.
\begin{lstlisting}[style=lean,firstnumber=63]
def ChiStatement : Prop :=
  ∃ C : ℝ, 0 < C ∧ ∀ (n : ℕ), 2 ≤ n → ∀ X ∈ feasible n,
  ∀ (βχ δ₁ : ℝ), 0 < βχ → βχ < 1 → 0 < δ₁ → δ₁ < 1 →
  ∀ (k : ℕ), 0 < k → C * Real.log (n / (βχ * δ₁)) ≤ k →
    HighProbability (gaussianBatch X k) (min βχ δ₁) (ChiGood X)
\end{lstlisting}

\paragraph{Corresponding theorem declaration (proof omitted).}
\textit{Source:} \SourceId{DPCut/Lemma10.lean}, line 134.
\begin{lstlisting}[style=lean,firstnumber=134]
theorem lemma10 : ChiStatement
\end{lstlisting}

\clearpage
\section{Lemma 11: A useful witness with constant probability}
\SourceId{lem:chaos}\hfill\SourceId{STOC2027.tex}:2026--2032
\paragraph{Paper statement.}
\setcounter{theorem}{10}
\begin{lemma}\label{lem:chaos}
For every fixed $D \in \mathbb{R}^{\binom{n}{2}}$ and $X \in \mathcal{X}$, choose independent random vectors $Z^{(j)}\sim \mathcal{N}(0,X)$ for each $j=1,\dots,k$. Set $\overline g(u,v)=\frac1k\sum_{j=1}^k\frac{(Z_u^{(j)}-Z_v^{(j)})^2}{4},$ and $\hat g(u,v)=\min\{\overline{g}(u,v)/2,1\}$.  For some universal constant $C > 0$, if
$k\geq C\log(n/(\beta_\chi\del_1))$, for a universal constant $c_0>0$, conditioning on the good event in Lemma~\ref{lem:chi},
\[
 \Prb\left[\langle D,\hat{g}\rangle\geq\frac{1}{2}\sum_{e=\{u,v\} \in \binom{n}{2}} D_{uv} q_{e}(X)\right]\geq c_0.
\]
\end{lemma}
\paragraph{Lean statement.}
\textit{Source:} \SourceId{DPCut/Statements.lean}, line 70.
\begin{lstlisting}[style=lean,firstnumber=70]
def ChaosStatement : Prop :=
  ∃ C c₀ : ℝ, 0 < C ∧ 0 < c₀ ∧ ∀ (n : ℕ), 2 ≤ n →
  ∀ (D : Weights n) X, X ∈ feasible n →
  ∀ (βχ δ₁ : ℝ), 0 < βχ → βχ < 1 → 0 < δ₁ → δ₁ < 1 →
  ∀ (k : ℕ), 0 < k → C * Real.log (n / (βχ * δ₁)) ≤ k →
    0 < prob (gaussianBatch X k) {z | ChiGood X z} ∧
    c₀ ≤ prob (gaussianBatch X k)
      {z | ChiGood X z ∧ pairing D (clippedWitness z) ≥ pairing D (fractionalQuery X) / 2} /
      prob (gaussianBatch X k) {z | ChiGood X z}
\end{lstlisting}

\paragraph{Corresponding theorem declaration (proof omitted).}
\textit{Source:} \SourceId{DPCut/Lemma11.lean}, line 107.
\begin{lstlisting}[style=lean,firstnumber=107]
theorem lemma11 : ChaosStatement
\end{lstlisting}

\paragraph{Comparison note.}
The conditional probability is written as a ratio of two probabilities; positivity of its denominator is part of the statement. The Lean theorem chooses $c_0=1/2048$, whereas the paper's proof obtains $1/648$; both satisfy the stated existential constant guarantee.
\clearpage
\section{Lemma 12: Preprocessing}
\SourceId{lem:preprocess}\hfill\SourceId{STOC2027.tex}:2204--2210
\paragraph{Paper statement.}
\setcounter{theorem}{11}
\begin{lemma}\label{lem:preprocess}
Fix any $\eps_0, \beta_0, \gamma \in (0, 1)$, and $n$-vertex input graph $G$ with edge-weight vector $w$. Decompose the weight vector $w=w'+v$, where $w'$ and $v$ are defined in \eqref{eq: def_split_weight}.
Then releasing $Y$ (as defined in \eqref{eq:noise_split_weight}) is $\eps_0$-DP, and, with probability at least $1-\beta_0$, for every cut $S$,
\[
 |Y(S,V\setminus S)+w'(S,V\setminus S)-w(S,V\setminus S)|\leq(\gamma/2)w(S,V\setminus S).
\]
\end{lemma}
\paragraph{Lean statement.}
\textit{Source:} \SourceId{DPCut/Statements.lean}, line 136.
\begin{lstlisting}[style=lean,firstnumber=136]
def PreprocessStatement : Prop :=
  ∀ (n : ℕ), 2 ≤ n → ∀ (ε₀ β₀ γ : ℝ),
  0 < ε₀ → ε₀ < 1 → 0 < β₀ → β₀ < 1 → 0 < γ → γ < 1 →
    EdgeDP (preprocessLaw (n := n) ε₀ β₀ γ) ε₀ 0 ∧
    ∀ w : Weights n, Nonnegative w →
      HighProbability (preprocessLaw ε₀ β₀ γ w) β₀
        (fun y => AllCuts w (y + clipWeights (clipThreshold n ε₀ β₀ γ) w) (γ / 2) 0)
\end{lstlisting}

\paragraph{Corresponding theorem declaration (proof omitted).}
\textit{Source:} \SourceId{DPCut/Lemma12.lean}, line 14.
\begin{lstlisting}[style=lean,firstnumber=14]
theorem lemma12 : PreprocessStatement
\end{lstlisting}

\clearpage
\section{Lemma 13: Stability of the SDP optimizer}
\SourceId{lem:stability_of_optimizer}\hfill\SourceId{STOC2027.tex}:2251--2261
\paragraph{Paper statement.}
\setcounter{theorem}{12}
\begin{lemma}\label{lem:stability_of_optimizer}
Let $A_1=A_0+\sum_e\xi_e b_eb_e^\top/4$, and let
\[
 X_i=\arg\max_{X\in\mathcal X}\{\ip{A_i}{X}+\lambda\log\det X\},
 \quad i\in\{0,1\}.
\]
Then
\[\norm{X_0^{-1/2}(X_1-X_0)X_1^{-1/2}}_F\leq \frac{\sum_e|\xi_e|\sqrt{q_e(X_0)q_e(X_1)}}{\lambda}.\]
Additionally, if $\norm{X_0^{-1/2}(X_1-X_0)X_1^{-1/2}}_F\leq1/2$, then for every $e$ it holds that
\[\frac{1}{2}q_e(X_1)\leq q_e(X_0)\leq 2q_e(X_1).\]
\end{lemma}
\paragraph{Lean statement.}
\textit{Source:} \SourceId{DPCut/Statements.lean}, line 31.
\begin{lstlisting}[style=lean,firstnumber=31]
def OptimizerStabilityStatement : Prop :=
  ∀ (n : ℕ), 2 ≤ n → ∀ (lam : ℝ), 0 < lam →
  ∀ (A₀ A₁ X₀ X₁ : Mat n) (ξ : Weights n),
    A₁ = A₀ + objectiveMatrix ξ →
    IsMatrixOptimizer lam A₀ X₀ → IsMatrixOptimizer lam A₁ X₁ →
    covarianceDistance X₀ X₁ ≤
      (∑ e, |ξ e| * Real.sqrt (fractionalQuery X₀ e * fractionalQuery X₁ e)) / lam ∧
    (covarianceDistance X₀ X₁ ≤ 1 / 2 → ∀ e,
      fractionalQuery X₁ e / 2 ≤ fractionalQuery X₀ e ∧
      fractionalQuery X₀ e ≤ 2 * fractionalQuery X₁ e)
\end{lstlisting}

\paragraph{Corresponding theorem declaration (proof omitted).}
\textit{Source:} \SourceId{DPCut/Lemma13.lean}, line 44.
\begin{lstlisting}[style=lean,firstnumber=44]
theorem lemma13 : OptimizerStabilityStatement
\end{lstlisting}

\paragraph{Comparison note.}
The paper's implicit assumptions $n\geq2$ and $\lambda>0$ are explicit. \texttt{covarianceDistance} uses the Frobenius norm, not the default matrix norm.
\clearpage
\section{Lemma 14: Gaussian covariance privacy}
\SourceId{lem:gaussian-privacy}\hfill\SourceId{STOC2027.tex}:2324--2326
\paragraph{Paper statement.}
\setcounter{theorem}{13}
\begin{lemma}\label{lem:gaussian-privacy}
Let $\del_1 \in(0,1/2)$ be a fixed parameter and $X_0,X_1\in\mathbb{R}^{n\times n}$ be symmetric positive definite matrices. Denote $E=X_0^{-1/2}(X_1-X_0)X_1^{-1/2}$. We use $\pdf_0$ and $\pdf_1$ to denote the probability density functions of $\mathcal{N}(0,X_0)$ and $\mathcal{N}(0,X_1)$ respectively. Suppose $\norm{E}_F<1/2$. Then for \[\eps_1=O\left(k\norm{E}_F^2+\sqrt{k\log\frac{1}{\del_1}}\norm{E}_F+\log\frac{1}{\del_1}\norm{E}_F\right),\] we have \[e^{-\eps_1}\prod_{j=1}^k\pdf_1(Z^{(j)})\leq\prod_{j=1}^k\pdf_0(Z^{(j)})\leq e^{\eps_1}\prod_{j=1}^k\pdf_1(Z^{(j)})\] with probability at least $1-\del_1$ over the joint distribution of the batch   $(Z^{(1)}, \dots, Z^{(k)})$, where each $Z^{(j)} \sim \mathcal{N}(0, X_0)$ is drawn independently.
\end{lemma}
\paragraph{Lean statement.}
\textit{Source:} \SourceId{DPCut/Statements.lean}, line 89.
\begin{lstlisting}[style=lean,firstnumber=89]
def GaussianPrivacyStatement : Prop :=
  ∃ C : ℝ, 0 < C ∧ ∀ (n k : ℕ), 0 < n → 0 < k →
  ∀ (X₀ X₁ : Mat n), X₀.PosDef → X₁.PosDef →
  ∀ (δ₁ : ℝ), 0 < δ₁ → δ₁ < 1 / 2 → covarianceDistance X₀ X₁ < 1 / 2 →
    let d := covarianceDistance X₀ X₁
    let ε₁ := C * (k * d ^ 2 + Real.sqrt (k * Real.log (1 / δ₁)) * d + Real.log (1 / δ₁) * d)
    HighProbability (gaussianBatch X₀ k) δ₁ (fun z => |batchLogRatio X₀ X₁ z| ≤ ε₁)
\end{lstlisting}

\paragraph{Corresponding theorem declaration (proof omitted).}
\textit{Source:} \SourceId{DPCut/Lemma14.lean}, line 135.
\begin{lstlisting}[style=lean,firstnumber=135]
theorem lemma14 : GaussianPrivacyStatement
\end{lstlisting}

\paragraph{Comparison note.}
Lean states the equivalent bound on the absolute log-density ratio. Both statements require $0<\delta_1<1/2$.
\clearpage
\section{Lemma 15: Signed chi-square concentration}
\SourceId{lem:concentration_chi-squares}\hfill\SourceId{STOC2027.tex}:2332--2338
\textit{Background result; not a contribution of this paper.}
\paragraph{Paper statement.}
\setcounter{theorem}{14}
\begin{lemma}[\cite{laurent2000adaptive}]\label{lem:concentration_chi-squares}
For a fixed parameter $t>0$, given independent standard Gaussian variables $\Sigma_i\sim\mathcal{N}(0,1)$ and weight $a_i$ for $i=1,\dots,n$.  Then,
\[
 \Prb\left[\left|\sum_{i=1}^n a_i(\Sigma_i^2-1)\right|
 >2\bigl(\norm a_2\sqrt t+\norm a_\infty t\bigr)\right]\leq2e^{-t}.
\]
\end{lemma}
\paragraph{Lean statement.}
\textit{Source:} \SourceId{DPCut/Statements.lean}, line 98.
\begin{lstlisting}[style=lean,firstnumber=98]
def WeightedChiSquareStatement : Prop :=
  ∀ (n : ℕ) (a : Fin n → ℝ) (t : ℝ), 0 < t →
    prob (standardGaussian n)
      {z | |∑ i, a i * (z i ^ 2 - 1)| >
        2 * (Real.sqrt (∑ i, a i ^ 2) * Real.sqrt t + ‖a‖ * t)} ≤ 2 * Real.exp (-t)
\end{lstlisting}

\paragraph{Corresponding theorem declaration (proof omitted).}
\textit{Source:} \SourceId{DPCut/Lemma15.lean}, line 13.
\begin{lstlisting}[style=lean,firstnumber=13]
theorem lemma15 : WeightedChiSquareStatement
\end{lstlisting}

\clearpage
\section{Lemma 16: Success-charged composition}
\SourceId{lem:pay-what-you-get}\hfill\SourceId{STOC2027.tex}:2406--2434
\paragraph{Paper statement.}
\setcounter{theorem}{15}
\begin{lemma}[Success-Charged Composition]
\label{lem:pay-what-you-get}
Suppose that the following conditions hold for every pair of adjacent inputs $w$ and $w'$. Let $\mathcal{M}_1, \dots, \mathcal{M}_k$ be a sequence of adaptively chosen randomized mechanisms, where each mechanism outputs either a valid result or a designated failure symbol $\bottom$.
For each step $t \in [k]$ and every common transcript prefix $y_{<t} = (y_1, \dots, y_{t-1})$, let $P_t^0$ and $P_t^1$ denote the conditional output distributions of $\mathcal{M}_t(w)$ and $\mathcal{M}_t(w')$, respectively. Suppose there exist auxiliary probability measures $N_t^0$ and $N_t^1$, a privacy loss parameter $\ell_t = \ell_t(y_{<t}) \ge 0$, and global constants $\ell^* < 1/10$ and $\delta' \ge 0$ such that, for any measurable set $A$ of return values and $b \in \{0, 1\}$:
\begin{itemize}
    \item If $P_t^b(A)=0$, then $N_t^b(A)=0;$
    \item $|P_t^b(A) - N_t^b(A)| \le \delta';$
    \item $e^{-\ell_t} N_t^1(A) \le N_t^0(A) \le e^{\ell_t} N_t^1(A),$
\end{itemize}
where $\ell_t(y_{<t}) \le \ell^*$.
Furthermore, suppose there exists a deterministic bound  $L$
such that for any sequence of outputs $y_1, \dots, y_k$ sampled sequentially according to
$y_t\sim N_t^b(\cdot\mid y_{<t}),$ for $t=1,\ldots,k$,
the accumulated privacy loss over successful attempts is bounded by $L$:
\begin{equation*}
    \sum_{t=1}^k \ell_t(y_{<t}) \cdot \ind[y_t \ne \bottom] \le L.
\end{equation*}
Then, for any fixed $\delta'' \in (0,1)$, the adaptive composition $(\mathcal{M}_1, \dots, \mathcal{M}_k)$ is $(\eps, \del)$-differentially private, where
\begin{equation*}
    \eps = 4L + 4\ell^*\log(1/\del''), \qquad \del = \del'' + (1+e^{\eps})k\del'.
\end{equation*}
\end{lemma}
\paragraph{Lean statement.}
\textit{Source:} \SourceId{DPCut/Privacy.lean}, line 76.
\begin{lstlisting}[style=lean,firstnumber=76]
def SuccessChargedStatement : Prop :=
  ∀ (Y : Type) [MeasurableSpace Y] [StandardBorelSpace Y] [Inhabited Y],
  ∀ (bottom : Y) (P₀ P₁ N₀ N₁ : Sequential Y) (k : ℕ)
    (ell : ℕ → List Y → ℝ) (ellStar L δ' δ'' : ℝ),
    (∀ t, IsMarkovKernel (P₀ t) ∧ IsMarkovKernel (P₁ t) ∧
      IsMarkovKernel (N₀ t) ∧ IsMarkovKernel (N₁ t)) →
    (∀ t, Measurable (ell t)) →
    0 ≤ ellStar → ellStar < 1 / 10 → 0 ≤ L → 0 ≤ δ' → 0 < δ'' → δ'' < 1 →
    (∀ t < k, ∀ h, CommonPrefix P₀ P₁ t h →
      0 ≤ ell t h ∧ ell t h ≤ ellStar ∧
      N₀ t h ≪ P₀ t h ∧ N₁ t h ≪ P₁ t h ∧
      TVClose (P₀ t h) (N₀ t h) δ' ∧ TVClose (P₁ t h) (N₁ t h) δ' ∧
      Indistinguishable (N₀ t h) (N₁ t h) (ell t h) 0) →
    (∀ᵐ h ∂sequentialLaw N₀ k, successCharge bottom ell h ≤ L) →
    (∀ᵐ h ∂sequentialLaw N₁ k, successCharge bottom ell h ≤ L) →
    let ε := 4 * L + 4 * ellStar * Real.log (1 / δ'')
    Indistinguishable (sequentialLaw P₀ k) (sequentialLaw P₁ k)
      ε (δ'' + (1 + Real.exp ε) * k * δ')
\end{lstlisting}

\paragraph{Corresponding theorem declaration (proof omitted).}
\textit{Source:} \SourceId{DPCut/Lemma16.lean}, line 20.
\begin{lstlisting}[style=lean,firstnumber=20]
theorem lemma16 : SuccessChargedStatement
\end{lstlisting}

\paragraph{Comparison note.}
Lean states indistinguishability for one neighboring pair. Applying it to every adjacent pair gives DP. Auxiliary laws are measurable Markov kernels, and the successful-charge budget holds almost surely under each auxiliary process.
\clearpage
\section{Lemma 17: Privacy of AggregateWitnesses}
\SourceId{lem:gradient_privacy}\hfill\SourceId{STOC2027.tex}:2441--2446
\paragraph{Paper statement.}
\setcounter{theorem}{16}
\begin{lemma}[Privacy of \Cref{alg:gradient_aggregation}]
\label{lem:gradient_privacy}
Fix any $\eps_{\mathrm{GE}}, \delta_{\mathrm{GE}}$ and $\beta_{\mathrm{GE}}\in(0,1)$, let $\eps_1=\frac{\eps_{\mathrm{GE}}}{320K(1+\log(2/\del_{\mathrm{GE}}))}, \del_1=\frac{\del_{\mathrm{GE}}}{20n^2K}, k=\Theta(\log\frac{n}{\beta_{\chi}\del_1}), \beta_{\chi}=\beta_{\mathrm{GE}}/4n^2K$ and $\lambda=\Theta(\frac{1}{\eps_1}\log\frac{n}{\beta_{\mathrm{GE}}\del_{\mathrm{GE}}}).$ Then each invocation of
\Cref{alg:gradient_aggregation} in \Cref{alg:main} with input $(w_i,h_i,\gamma)$ is
$(\eps_{\mathrm{GE}}, \delta_{\mathrm{GE}})$-differentially private.
\end{lemma}
\paragraph{Lean statement.}
\textit{Source:} \SourceId{DPCut/Statements.lean}, line 145.
\begin{lstlisting}[style=lean,firstnumber=145]
def AggregatePrivacyStatement (c : Constants) : Prop :=
  ∀ (n : ℕ), 2 ≤ n → ∀ (ε δ β γ : ℝ),
  0 < ε → ε < 1 → 0 < δ → δ < 1 → 0 < β → β < 1 → 0 < γ → γ < 1 →
  ∀ h : Weights n, Nonnegative h →
    EdgeDP (fun w => aggregateLaw (oracleConfig n c ε δ β) γ w h) ε δ
\end{lstlisting}

\paragraph{Corresponding theorem declaration (proof omitted).}
\textit{Source:} \SourceId{DPCut/Lemma17.lean}, line 111.
\begin{lstlisting}[style=lean,firstnumber=111]
theorem lemma17 : AggregatePrivacyStatement privacyConstants
\end{lstlisting}

\paragraph{Comparison note.}
The coefficient record \texttt{c} makes the paper's hidden universal constants explicit. The displayed theorem supplies \texttt{privacyConstants}; its definition is reproduced in the notation appendix.
\clearpage
\section{Lemma 18: Privacy of CutRelease}
\SourceId{lem:main-privacy}\hfill\SourceId{STOC2027.tex}:2670--2674
\paragraph{Paper statement.}
\setcounter{theorem}{17}
\begin{lemma}[Privacy of \Cref{alg:main}]\label{lem:main-privacy}
For $0<\varepsilon\leq1$ and $0<\delta<1$, fix $\eps_2=\frac{\eps}{2T_1}$, $\eps_{\mathrm{GE}}=\frac{\eps}{4T_1\sqrt{2T_2\log(2T_1/\delta)}}$, $\delta_{\mathrm{GE}}=\frac{\delta}{2T_1T_2}.$
Then \Cref{alg:main} is
$(\eps, \del)$-differentially private.
\end{lemma}
\paragraph{Lean statement.}
\textit{Source:} \SourceId{DPCut/Statements.lean}, line 161.
\begin{lstlisting}[style=lean,firstnumber=161]
def MainPrivacyStatement (c : Constants) : Prop :=
  ∀ (n : ℕ), 2 ≤ n → ∀ (U : ℝ), 1 ≤ U → ∀ p : Parameters, p.PrivacyValid →
    EdgeDP (cutReleaseLaw (n := n) U p c) p.epsilon p.delta
\end{lstlisting}

\paragraph{Corresponding theorem declaration (proof omitted).}
\textit{Source:} \SourceId{DPCut/Lemma18.lean}, line 75.
\begin{lstlisting}[style=lean,firstnumber=75]
theorem lemma18 : MainPrivacyStatement privacyConstants
\end{lstlisting}

\paragraph{Comparison note.}
\texttt{Parameters.PrivacyValid} gives the broader privacy-only parameter range. The target refers to the same \texttt{cutReleaseLaw} as the utility target.
\clearpage
\section{Lemma 19: The SDP cut certificate}
\SourceId{lem:sdp_error}\hfill\SourceId{STOC2027.tex}:2714--2719
\paragraph{Paper statement.}
\setcounter{theorem}{18}
\begin{lemma}\label{lem:sdp_error}
For every signed vector $D \in \mathbb{R}^{\binom{n}{2}}$ and every cut $S \subseteq [n]$,
\[\langle D, c_S\rangle\leq 2\max\{F(D),0\}+O(\lambda n\log n+\norm{D}_1/n).\]
Moreover, $\langle D,\one\rangle\leq2F(D)$.
\end{lemma}
\paragraph{Lean statement.}
\textit{Source:} \SourceId{DPCut/Statements.lean}, line 43.
\begin{lstlisting}[style=lean,firstnumber=43]
def SDPCertificateStatement : Prop :=
  ∃ C : ℝ, 0 < C ∧ ∀ (n : ℕ), 2 ≤ n → ∀ (lam : ℝ), 0 < lam →
  ∀ (D : Weights n),
    (∀ S, cut D S ≤ 2 * max (optimum lam D) 0 +
      C * (lam * n * Real.log n + l1 D / n)) ∧ mass D ≤ 2 * optimum lam D
\end{lstlisting}

\paragraph{Corresponding theorem declaration (proof omitted).}
\textit{Source:} \SourceId{DPCut/Lemma19.lean}, line 85.
\begin{lstlisting}[style=lean,firstnumber=85]
theorem lemma19 : SDPCertificateStatement
\end{lstlisting}

\clearpage
\section{Lemma 20: Utility of AggregateWitnesses}
\SourceId{lem:gradient_estimation}\hfill\SourceId{STOC2027.tex}:2737--2748
\paragraph{Paper statement.}
\setcounter{theorem}{19}
\begin{lemma}\label{lem:gradient_estimation}
Let $0<\gamma<1$, $K=\Theta(\log(n/\beta_{\mathrm{GE}}))$, $\tau_1=\log(1/\beta_1)/\eps_1$ and $\beta_1=\beta_{\chi}=\beta_{\mathrm{GE}}/4n^2K$. Given nonnegative edge weights $w$ and a nonnegative partial estimate $h$, suppose $w_e\leq U$ for every $e$. Let $(g^+,g^-,r^*)=\mathsf{AggregateWitnesses}(w,h,\gamma)$ (\Cref{alg:gradient_aggregation}).
Then with probability at least $1-\beta_{\mathrm{GE}}$:
\begin{itemize}
    \item $\langle h,g^+\rangle\geq(1+\gamma)\langle w,g^+\rangle;$
    \item $\langle h,g^-\rangle\leq(1-\gamma)\langle w,g^-\rangle;$
    \item For every cut $S$ simultaneously, \[(1-\gamma)\langle w,c_S\odotc r^*\rangle-B\leq \langle h,c_S\odotc r^*\rangle\leq(1+\gamma)\langle w,c_S\odotc r^*\rangle+B,\]
where $B=\Theta\left(\tau_1+\lambda n\log n+nU\right).$
\end{itemize}
\end{lemma}
\paragraph{Lean statement.}
\textit{Source:} \SourceId{DPCut/Statements.lean}, line 152.
\begin{lstlisting}[style=lean,firstnumber=152]
def AggregateUtilityStatement (c : Constants) : Prop :=
  ∀ (n : ℕ), 2 ≤ n → ∀ (U ε δ β γ : ℝ),
  0 ≤ U → 0 < ε → ε < 1 → 0 < δ → δ < 1 → 0 < β → β < 1 → 0 < γ → γ < 1 →
  ∀ (w h : Weights n), Bounded U w → Nonnegative h →
    let o := oracleConfig n c ε δ β
    let B := c.certificate * (o.tau1 + o.lam * n * Real.log n + n * U)
    HighProbability (aggregateLaw o γ w h) β (fun a => a.Partition ∧ a.Accurate w h γ B)
\end{lstlisting}

\paragraph{Corresponding theorem declaration (proof omitted).}
\textit{Source:} \SourceId{DPCut/Lemma20.lean}, line 15.
\begin{lstlisting}[style=lean,firstnumber=15]
theorem lemma20 : AggregateUtilityStatement utilityConstants
\end{lstlisting}

\paragraph{Comparison note.}
Lean also includes the deterministic partition invariant $g^++g^-+r=\mathbf1$ in \texttt{Aggregate.Partition}. \texttt{Aggregate.Accurate} expands to the three displayed accuracy requirements. The supplied certificate coefficient is $16$.
\clearpage
\section{Lemma 21: Number of inner iterations}
\SourceId{lem:inner_round}\hfill\SourceId{STOC2027.tex}:2816--2824
\paragraph{Paper statement.}
\setcounter{theorem}{20}
\begin{lemma}\label{lem:inner_round}
Let $0 < \gamma \leq 1/4$, $\eta=\gamma/[8(1+\gamma)]$, $\tau_2=\log(1/\beta_2)/\eps_2$,
$B=O(\tau_1+\lambda n\log n+nU)$ as in \Cref{alg:main}, and
$M_{\mathrm{stop}}=C(nU\log\frac{nU}{\delta\beta\gamma\eps}+\frac{n}{\gamma\varepsilon}(\log\frac{nU}{\delta\beta\gamma\eps})^6)\geq10(B+\tau_2/\gamma)$, for a sufficiently large universal constant $C$.
For any fixed outer iteration $i$, with probability at least $1-T_2\beta_{\mathrm{GE}}-\beta_2$, if $\widetilde M_i>M_{\rm stop}$, then within at most $T_2=1+\left\lceil64(1+\gamma)^2\log n/\gamma^2\right\rceil=O(\log n/\gamma^2)$ inner iterations, $\mathsf{AggregateWitnesses}(w_i,h_i,\gamma)$ returns $(g^+_i,g^-_i,r_i)$ satisfying
\[
\langle h_i,r_i\rangle\geq\max\{\langle h_i,g^+_i\rangle,\langle h_i,g^-_i\rangle\}.
\]
\end{lemma}
\paragraph{Lean statement.}
\textit{Source:} \SourceId{DPCut/Utility.lean}, line 31.
\begin{lstlisting}[style=lean,firstnumber=31]
def InnerRoundStatement (c : Constants) : Prop :=
  ∀ (n : ℕ), 2 ≤ n → ∀ (U : ℝ), 1 ≤ U → ∀ p : Parameters, p.IterationValid →
  ∀ (w a : Weights n), Bounded U w → IsMask a →
    let cfg := runConfig n U p c
    cfg.stop ≥ 10 * (cfg.B + cfg.tau2 / p.gamma) →
    HighProbability (innerRoundLaw p cfg w a) (cfg.T2 * cfg.betaGE + cfg.beta2)
      (fun result => cfg.stop < result.1 → result.2.isSome)
\end{lstlisting}

\paragraph{Corresponding theorem declaration (proof omitted).}
\textit{Source:} \SourceId{DPCut/Lemma21.lean}, line 117.
\begin{lstlisting}[style=lean,firstnumber=117]
theorem lemma21 : InnerRoundStatement utilityConstants
\end{lstlisting}

\paragraph{Comparison note.}
Lean makes the algorithm's surrounding bounded-input setting explicit. \texttt{Parameters.IterationValid} permits $0<\delta,\beta<1$, without the main theorem's restrictions $\delta,\beta<1/10$. The proof establishes the oracle and noise prerequisites over this broader range. For each fixed preceding history, success probability is joint over the new mass estimate and the inner loop. No nonzero-mask hypothesis is needed: on the good mass-estimation event, continuing implies positive true mass. \texttt{result.2.isSome} means that \texttt{innerLoop} returned a candidate and mask; that branch requires \texttt{Aggregate.Accepts}.
\clearpage
\section{Lemma 22: Number of outer iterations}
\SourceId{lem:outer_round}\hfill\SourceId{STOC2027.tex}:2944--2946
\paragraph{Paper statement.}
\setcounter{theorem}{21}
\begin{lemma}\label{lem:outer_round}
With probability at least $1-T_1T_2\beta_{\mathrm{GE}}-T_1\beta_2$, \Cref{alg:main} terminates with $a_i=0$ or $\widetilde{M}_i\leq M_{\mathsf{stop}}$ in at most $T_1 = O(\log(nU))$ outer rounds.
\end{lemma}
\paragraph{Lean statement.}
\textit{Source:} \SourceId{DPCut/Utility.lean}, line 54.
\begin{lstlisting}[style=lean,firstnumber=54]
def OuterRoundStatement (c : Constants) : Prop :=
  ∀ (n : ℕ), 2 ≤ n → ∀ (U : ℝ), 1 ≤ U → ∀ p : Parameters, p.IterationValid →
  ∀ w : Weights n, Bounded U w →
    let cfg := runConfig n U p c
    HighProbability (outerStopsLaw p cfg w cfg.T1 1)
      (cfg.T1 * cfg.T2 * cfg.betaGE + cfg.T1 * cfg.beta2) (fun stopped => stopped = true)
\end{lstlisting}

\paragraph{Corresponding theorem declaration (proof omitted).}
\textit{Source:} \SourceId{DPCut/Lemma22.lean}, line 37.
\begin{lstlisting}[style=lean,firstnumber=37]
theorem lemma22 : OuterRoundStatement outerUtilityConstants
\end{lstlisting}

\paragraph{Comparison note.}
Lean makes the algorithm's surrounding bounded-input setting explicit and uses \texttt{Parameters.IterationValid}, allowing $0<\delta,\beta<1$. The stopping-threshold calibration is proved over this broader range, not assumed. \texttt{outerStopsLaw} records whether the actual stopping test is reached within the allotted horizon, rather than merely whether the loop runs out of iterations.
\clearpage
\section{Lemma 23: Utility of CutRelease}
\SourceId{lem:utility}\hfill\SourceId{STOC2027.tex}:2982--2998
\paragraph{Paper statement.}
\setcounter{theorem}{22}
\begin{lemma}\label{lem:utility}
Let $n\geq2$, $U\geq1$, $0<\gamma\leq1/4$, $0<\eps\leq1$, and $0<\delta,\beta<1/10$.
Choose all parameters as in \Cref{alg:main}, in particular
$\eps_{\mathrm{GE}}=\eps/(4T_1\sqrt{2T_2\log(2T_1/\delta)})$,
$\delta_{\mathrm{GE}}=\delta/(2T_1T_2)$, and $M_{\mathrm{stop}}=C(nU\log\frac{nU}{\delta\beta\gamma\eps}+\frac{n}{\gamma\varepsilon}(\log\frac{nU}{\delta\beta\gamma\eps})^6)
\geq10(B+\tau_2/\gamma)$, for a sufficiently large universal constant $C$.
Given a nonnegative weighted $n$-vertex graph $G$ with edge weights bounded by $U$, \Cref{alg:main} outputs a
nonnegative weighted graph $\widetilde{G}$ satisfying,
with probability at least $1-\beta$, simultaneously
for every $S\subseteq V$,
\begin{align*}
 \left|w_G(S,V\setminus S)-w_{\widetilde{G}}(S,V\setminus S)\right|
 \leq \gamma w_G(S,V\setminus S)  + O\left(nU\log\frac{nU}{\delta\beta\gamma\eps}+\frac{n}{\gamma\varepsilon}
 \left(\log\frac{nU}{\delta\beta\gamma\eps}\right)^7\right).
\end{align*}
\end{lemma}
\paragraph{Lean statement.}
\textit{Source:} \SourceId{DPCut/Statements.lean}, line 166.
\begin{lstlisting}[style=lean,firstnumber=166]
def BoundedUtilityStatement (c : Constants) : Prop :=
  ∃ C : ℝ, 0 < C ∧ ∀ (n : ℕ), 2 ≤ n → ∀ (U : ℝ), 1 ≤ U →
  ∀ p : Parameters, p.Valid → ∀ w : Weights n, Bounded U w →
    HighProbability (cutReleaseLaw U p c w) p.beta
      (fun h => Nonnegative h ∧ AllCuts w h p.gamma (C * boundedError n U p))
\end{lstlisting}

\paragraph{Corresponding theorem declaration (proof omitted).}
\textit{Source:} \SourceId{DPCut/Lemma23.lean}, line 30.
\begin{lstlisting}[style=lean,firstnumber=30]
theorem lemma23 : BoundedUtilityStatement outerUtilityConstants
\end{lstlisting}

\paragraph{Comparison note.}
\texttt{boundedError} contains the two additive terms, including exponent $7$. The multiplicative constant is quantified before all inputs and parameters. Lean also records nonnegative output on the accuracy event.
\clearpage
\section{Lemma 24: Chi-square tail bound}
\SourceId{lem:chi_square_tail}\hfill\SourceId{STOC2027.tex}:3384--3389
\textit{Background result; not a contribution of this paper.}
\paragraph{Paper statement.}
\setcounter{theorem}{23}
\begin{lemma}[Chi-squared Tail Bound~\cite{laurent2000adaptive}]\label{lem:chi_square_tail}
Let $Y_1, \dots, Y_k$ be independent standard normal random variables, i.e., $Y_i \sim \mathcal{N}(0, 1)$. The random variable $Q = \sum_{i=1}^k Y_i^2$ follows a Chi-squared distribution with $k$ degrees of freedom, denoted as $Q \sim \chi_k^2$. For any such $Q$, there exists a universal constant $c > 0$ such that
\[
 \Prb\left[\frac{Q}{k} \notin \left[\frac{1}{2}, 2\right]\right] \leq 2e^{-ck}.
\]
\end{lemma}
\paragraph{Lean statement.}
\textit{Source:} \SourceId{DPCut/Statements.lean}, line 105.
\begin{lstlisting}[style=lean,firstnumber=105]
def ChiSquareTailStatement : Prop :=
  ∃ c : ℝ, 0 < c ∧ ∀ k : ℕ, 0 < k →
    prob (standardGaussian k)
      {z | (∑ i, z i ^ 2) / k ∉ Set.Icc (1 / 2 : ℝ) 2} ≤ 2 * Real.exp (-c * k)
\end{lstlisting}

\clearpage
\section{Lemma 25: Gaussian hypercontractivity}
\SourceId{lem:hypercontractivity}\hfill\SourceId{STOC2027.tex}:3393--3399
\textit{Background result; not a contribution of this paper.}
\paragraph{Paper statement.}
\setcounter{theorem}{24}
\begin{lemma}[Gaussian Hypercontractivity for Degree-Two Polynomials~\cite{Janson97}]\label{lem:hypercontractivity}
Let $W$ be a degree-two polynomial in independent standard Gaussian random variables. Then its higher moments can be bounded by its second moment. Specifically, for the $4$-th moment, we have:
\[
    \|W\|_4 \le 3\|W\|_2,
\]
where $\|W\|_p = (\E[|W|^p])^{1/p}$ denotes the $L_p$ norm of the random variable.
\end{lemma}
\paragraph{Lean statement.}
\textit{Source:} \SourceId{DPCut/Statements.lean}, line 111.
\begin{lstlisting}[style=lean,firstnumber=111]
def HypercontractivityStatement : Prop :=
  ∀ (m : ℕ) (P : MvPolynomial (Fin m) ℝ), P.totalDegree ≤ 2 →
    eLpNorm (fun z => MvPolynomial.eval z P) 4 (standardGaussian m) ≤
      3 * eLpNorm (fun z => MvPolynomial.eval z P) 2 (standardGaussian m)
\end{lstlisting}

\clearpage
\section{Lemma 26: Interpolation of Lp norms}
\SourceId{lem:holder_interpolation}\hfill\SourceId{STOC2027.tex}:3401--3407
\textit{Background result; not a contribution of this paper.}
\paragraph{Paper statement.}
\setcounter{theorem}{25}
\begin{lemma}[Log-Convexity of $L_p$ Norms]\label{lem:holder_interpolation}
For any random variable $W$ and $1 \le p < r < q \le \infty$, the $L_p$ norms satisfy the interpolation inequality $\|W\|_r \le \|W\|_p^\theta \|W\|_q^{1-\theta}$, where $\theta \in (0, 1)$ is determined by $\frac{1}{r} = \frac{\theta}{p} + \frac{1-\theta}{q}$.
In particular, choosing $p=1, r=2$, and $q=4$ yields $\theta = 1/3$, which gives:
\[
    \|W\|_2 \le \|W\|_1^{1/3} \|W\|_4^{2/3}.
\]
\end{lemma}
\paragraph{Lean statement.}
\textit{Source:} \SourceId{DPCut/Statements.lean}, line 123.
\begin{lstlisting}[style=lean,firstnumber=123]
def GeneralHolderInterpolationStatement : Prop :=
  ∀ (Ω : Type) [MeasurableSpace Ω] (μ : Measure Ω) (W : Ω → ℝ),
  ∀ (p r q : ℝ≥0∞) (θ : ℝ), AEStronglyMeasurable W μ →
    1 ≤ p → p < r → r < q → 0 < θ → θ < 1 →
    r⁻¹ = ENNReal.ofReal θ * p⁻¹ + ENNReal.ofReal (1 - θ) * q⁻¹ →
    eLpNorm W r μ ≤ (eLpNorm W p μ) ^ θ * (eLpNorm W q μ) ^ (1 - θ)
\end{lstlisting}

\textit{Source:} \SourceId{DPCut/Statements.lean}, line 117.
\begin{lstlisting}[style=lean,firstnumber=117]
def HolderInterpolationStatement : Prop :=
  ∀ (Ω : Type) [MeasurableSpace Ω] (μ : Measure Ω) (W : Ω → ℝ),
    AEStronglyMeasurable W μ →
    eLpNorm W 2 μ ≤ (eLpNorm W 1 μ) ^ (1 / 3 : ℝ) * (eLpNorm W 4 μ) ^ (2 / 3 : ℝ)
\end{lstlisting}

\paragraph{Comparison note.}
The first target covers the full exponent range, including $q=\infty$. The second is the $(1,2,4)$ instance used in the witness argument. Lean uses the extended $L_p$ seminorm and makes measurability explicit.
\clearpage
\section{Proposition 27: Unbiased fractional-query estimator}
\SourceId{prop:unbiased_estimator}\hfill\SourceId{STOC2027.tex}:3411--3417
\paragraph{Paper statement.}
\setcounter{theorem}{26}
\begin{proposition}[Unbiased Estimator of Fractional Query]\label{prop:unbiased_estimator}
Let $X \in \mathcal{X}$ be a valid covariance matrix and let $Z \sim \mathcal{N}(0, X)$. For any pair of vertices $e = \{u,v\}$, we have
\[
    \E\left[\frac{(Z_u - Z_v)^2}{4}\right] = q_e(X).
\]
Consequently, for $\overline{g}(u,v) = \frac{1}{k} \sum_{j=1}^k \frac{(Z^{(j)}_u - Z^{(j)}_v)^2}{4}$ where $Z^{(j)}$ are drawn independently from $\mathcal{N}(0, X)$, it holds that $\E[\overline{g}(u,v)] = q_e(X)$.
\end{proposition}
\paragraph{Lean statement.}
\textit{Source:} \SourceId{DPCut/Statements.lean}, line 57.
\begin{lstlisting}[style=lean,firstnumber=57]
def UnbiasedWitnessStatement : Prop :=
  ∀ (n : ℕ), 2 ≤ n → ∀ X ∈ feasible n, ∀ (k : ℕ), 0 < k → ∀ e : Edge n,
    (∫ z, (z e.val.1 - z e.val.2) ^ 2 / 4 ∂gaussian X) = fractionalQuery X e ∧
    (∫ z, rawWitness z e ∂gaussianBatch X k) = fractionalQuery X e
\end{lstlisting}

\clearpage
\section{Lemma 28: KL Pythagorean inequality}
\SourceId{lem:pythagorean}\hfill\SourceId{STOC2027.tex}:3594--3607
\textit{Background result; not a contribution of this paper.}
\paragraph{Paper statement.}
\setcounter{theorem}{27}
\begin{lemma}[Generalized Pythagorean Inequality for KL Divergence~\cite{Csiszar75}]\label{lem:pythagorean}
Let $\mathcal{K}$ be a closed convex set of probability distributions, and let $\overline{p}$ be an arbitrary strictly positive probability distribution. If $p'$ is the projection of $\overline{p}$ onto $\mathcal{K}$, defined as:
\[
    p' = \argmin_{z \in \mathcal{K}} \KL(z \| \overline{p}),
\]
then for any target distribution $p^* \in \mathcal{K}$, it holds that:
\[
    \KL(p^* \| \overline{p}) \ge \KL(p^* \| p') + \KL(p' \| \overline{p}).
\]
In particular, since $\KL(p' \| \overline{p}) \ge 0$, this implies the distance-reduction property:
\[
    \KL(p^* \| p') \le \KL(p^* \| \overline{p}).
\]
\end{lemma}
\paragraph{Lean statement.}
\textit{Source:} \SourceId{DPCut/Applications.lean}, line 94.
\begin{lstlisting}[style=lean,firstnumber=94]
def KLPythagoreanStatement : Prop :=
  ∀ (m : ℕ) (K : Set (Fin m → ℝ)), IsClosed K → Convex ℝ K →
    (∀ p ∈ K, IsDistribution p) →
    ∀ (pbar p' pstar : Fin m → ℝ), IsDistribution pbar → (∀ i, 0 < pbar i) →
      p' ∈ K → (∀ q ∈ K, finiteKL p' pbar ≤ finiteKL q pbar) → pstar ∈ K →
      finiteKL pstar p' + finiteKL p' pbar ≤ finiteKL pstar pbar
\end{lstlisting}

\paragraph{Comparison note.}
The paper states this cited auxiliary result without restricting the underlying space to a finite set. Lean proves the finite-domain specialization indexed by \texttt{Fin m} that the paper's applications require; formalizing the full generality of this background result is outside the validation scope. Membership and minimality of the projection are hypotheses; the nonnegative extended KL divergence also implies the final distance-reduction inequality.
\clearpage
\section{Lemma 29: Capped unit-bounded cut release}
\SourceId{lem:capped-unit-release}\hfill\SourceId{STOC2027.tex}:3609--3620
\paragraph{Paper statement.}
\setcounter{theorem}{28}
\begin{lemma}[Capped unit-bounded cut release]\label{lem:capped-unit-release}
Let $D\subseteq\binom V2$ be public and let $x\in[0,1]^D$. For $0<\gamma\leq1/4$, $0<\varepsilon\leq1$, and $0<\delta,\beta_{\rm cap}<1/10$, there is a polynomial-time $(\varepsilon, \delta)$-DP mechanism under adjacency defined in Definition~\ref{def:DP} on $x$ that outputs $q\in[0,1]^D$ such that, except with probability $\beta_{\rm cap}$, simultaneously for every $S\subseteq V$,
\begin{equation*}
 \left|q(\delta_D(S))-x(\delta_D(S))\right| \leq \gamma x(\delta_D(S))+A_{\rm cap},
\end{equation*}
where
\[
 A_{\rm cap} \leq O\left(\frac{n}{\gamma\varepsilon}
 \left(\log\frac{n}{\delta\beta_{\rm cap}\gamma\varepsilon}\right)^7\right).
\]
\end{lemma}
\paragraph{Lean statement.}
\textit{Source:} \SourceId{DPCut/Applications.lean}, line 80.
\begin{lstlisting}[style=lean,firstnumber=80]
def CappedReleaseStatement : Prop :=
  ∃ C : ℝ, 0 < C ∧ ∀ (n : ℕ), 2 ≤ n → ∀ (D : Finset (Edge n)) (p : Parameters), p.Valid →
  ∃ M : Kernel (Weights n) (Weights n), IsMarkovKernel M ∧
    DP (fun w w' => CappedOn D w ∧ CappedOn D w' ∧ Adjacent w w') M p.epsilon p.delta ∧
    ∀ w, CappedOn D w → (∀ᵐ h ∂M w, CappedOn D h) ∧
      HighProbability (M w) p.beta (fun h => AllCuts w h p.gamma
        (C * mainError n p))
\end{lstlisting}

\paragraph{Corresponding theorem declaration (proof omitted).}
\textit{Source:} \SourceId{DPCut/Lemma29.lean}, line 69.
\begin{lstlisting}[style=lean,firstnumber=69]
theorem lemma29 : CappedReleaseStatement
\end{lstlisting}

\paragraph{Combined implementation theorem (proof omitted).}
\textit{Source:} \SourceId{DPCut/CappedLSWRuntime.lean}, line 115. \textit{Namespace:} \SourceId{DPCut.CappedLSW}.
\begin{lstlisting}[style=lean,firstnumber=115]
theorem capped_release_of_lsw (hLSW : LSWStatement) :
    ∃ Ctime : ℝ, 0 < Ctime ∧ ∃ degree : ℕ,
    ∀ n (hn : 2 ≤ n) (p : Parameters) (hp : p.Valid) (D : Finset (Edge n)),
      let F := cuttingPlaneFamily hLSW verifiedLinearAlgebra
      DP (fun w w' => CappedOn D w ∧ CappedOn D w' ∧ Adjacent w w')
        (kernel F hn p hp D) p.epsilon p.delta ∧
      ∀ w : Weights n,
        (run F hn p hp D w).cost ≤ ENNReal.ofReal (Ctime * runtimeScale n p ^ degree) ∧
        (∀ᵐ h ∂(run F hn p hp D w).law, CappedOn D h) ∧
        (CappedOn D w → HighProbability (run F hn p hp D w).law p.beta
          (fun h => AllCuts w h p.gamma ((2:ℝ)^138 * mainError n p)))
\end{lstlisting}

\paragraph{Comparison note.}
This Lean target contains the statistical guarantee and capped output support. As in the paper's proof, \texttt{cappedRunConfig} keeps $\gamma$ unchanged and multiplies the main stopping threshold by ten, ensuring $M_{\mathrm{stop}}\geq100(B+\tau_2/\gamma)$. The error is \texttt{C * mainError n p}, with fixed logarithmic exponent seven and a universal coefficient; \texttt{lemma29} chooses $C=2^{126}$. The combined theorem \texttt{CappedLSW.capped\_release\_of\_lsw} proves capped support, privacy, utility with coefficient $2^{138}$, and polynomial assigned cost in the full parameter scale under the formal LSW parameter. This does not certify operational running time.
\clearpage
\section{Theorem 30: Private correlation-clustering reduction}
\SourceId{thm:cc-black-box-reduction}\hfill\SourceId{STOC2027.tex}:3714--3724
\paragraph{Paper statement.}
\setcounter{theorem}{29}
\begin{theorem}[Black-box private correlation-clustering reduction]
\label{thm:cc-black-box-reduction}
Let $\rho\geq1$ and $p_0>0$ be constants, and let $\mathcal A$ be a polynomial-time randomized algorithm such that, on every unweighted complete-graph min-disagree instance $Q$,
\[
 \Prb\left[\operatorname{cc}_Q(\mathcal A(Q)) \leq \rho\operatorname{OPT}_Q\right] \geq p_0.
\]
Then, for every fixed $\xi>0$, $\varepsilon>0$, and $\delta=n^{-10}$, there is a polynomial-time $(\varepsilon, \delta)$-DP algorithm that, on every complete-graph min-disagree instance $G$, with high probability, outputs a partition $\widehat\Pi$ satisfying
\[
 \operatorname{cc}_G(\widehat\Pi) \leq (\rho+\xi)\operatorname{OPT}_G+\widetilde O_{\rho,\xi,\varepsilon}(n).
\]
\end{theorem}
\paragraph{Lean statement.}
\textit{Source:} \SourceId{DPCut/Applications.lean}, line 160.
\begin{lstlisting}[style=lean,firstnumber=160]
def CCBlackBoxStatement : Prop :=
  ∀ (ρ ξ ε : ℝ), 1 ≤ ρ → 0 < ξ → 0 < ε → ∃ C Cfail : ℝ, ∃ exponent : ℕ,
    0 < C ∧ 0 < Cfail ∧ ∀ p₀ : ℝ, 0 < p₀ →
  ∀ algorithm : (n : ℕ) → Kernel (Weights n) (Partition n),
    (∀ n, IsMarkovKernel (algorithm n)) →
    (∀ (n : ℕ), 2 ≤ n → ∀ w, IsUnweighted w →
      HighProbability (algorithm n w) (1 - p₀) (fun π => ccCost w π ≤ ρ * ccOptimum w)) →
    PrivateCCGuarantee C Cfail ε (ρ + ξ) exponent id
\end{lstlisting}

\paragraph{Corresponding theorem declaration (proof omitted).}
\textit{Source:} \SourceId{DPCut/CCApplication.lean}, line 81.
\begin{lstlisting}[style=lean,firstnumber=81]
theorem theorem30 (baseline : CCBaselineFamily) : CCBlackBoxStatement
\end{lstlisting}

\paragraph{Combined implementation theorem (proof omitted).}
\textit{Source:} \SourceId{DPCut/ImplementedApplications.lean}, line 124. \textit{Namespace:} \SourceId{DPCut.LSWApplications}.
\begin{lstlisting}[style=lean,firstnumber=124]
theorem theorem30 (hLSW : LSWStatement) (baseline : CCBaselineFamily)
    {ρ ξ ε : ℝ} (hρ : 1 ≤ ρ) (hξ : 0 < ξ) (hε : 0 < ε) :
    ∃ C Cfail : ℝ, ∃ exponent : ℕ, 0 < C ∧ 0 < Cfail ∧
    ∀ s : ℝ, 0 < s → ∀ alg : CCApproximationContract ρ s,
      ImplementedCCGuarantee hLSW (baseline ε hε).toPolynomialGraphAlgorithm
        alg.toPolynomialGraphAlgorithm C Cfail ε (ρ+ξ) exponent id
\end{lstlisting}

\paragraph{Comparison note.}
The original Lean target records the statistical guarantee. The combined implementation theorem also provides a polynomial assigned-cost bound for the specified approximate program, using the formal LSW parameter and cited external algorithm contracts. These interfaces allow arbitrary kernels with cost zero; actual running times of the external algorithms and full reduction are not certified. It makes $\widetilde O(n)$ explicit as $C n(\log n)^d$ and high probability as failure at most $C_{\rm fail}/n^2$, with constants chosen before $n$. The source algorithm is a family of Markov kernels with success probability at least $p_0$. The constants are chosen after $\rho,\xi,\varepsilon$ and before $p_0$ and the source algorithm, matching the paper's uniform additive bound. The theorem takes the cited private baseline as an explicit external input.
\clearpage
\section{Corollary 31: Private fixed-k correlation clustering}
\SourceId{cor:cc-fixed-k}\hfill\SourceId{STOC2027.tex}:3800--3811
\paragraph{Paper statement.}
\setcounter{theorem}{30}
\begin{corollary}[Private fixed-$k$ correlation clustering]
\label{cor:cc-fixed-k}
For every fixed $k\geq2$ and $\xi>0$, constant $\varepsilon>0$, and $\delta=n^{-10}$, there is a
polynomial-time $(\varepsilon, \delta)$-edge-DP algorithm for complete unweighted
Min-Disagree with at most $k$ clusters which, with high probability, outputs a partition
$\widehat\Pi$ satisfying
\[
 \operatorname{cc}_G(\widehat\Pi)
 \leq (1+\xi)\operatorname{OPT}_{G,k}
   +\widetilde O_{\xi,\varepsilon}(n).
\]
\end{corollary}
\paragraph{Lean statement.}
\textit{Source:} \SourceId{DPCut/Applications.lean}, line 170.
\begin{lstlisting}[style=lean,firstnumber=170]
def FixedKCorrelationClusteringStatement : Prop :=
  ∀ (ξ ε : ℝ), 0 < ξ → 0 < ε → ∃ C Cfail : ℝ, ∃ exponent : ℕ,
    0 < C ∧ 0 < Cfail ∧ ∀ k : ℕ, 2 ≤ k →
      PrivateCCGuarantee C Cfail ε (1 + ξ) exponent (fun _ => k)
\end{lstlisting}

\paragraph{Corresponding theorem declaration (proof omitted).}
\textit{Source:} \SourceId{DPCut/Theorem30.lean}, line 10.
\begin{lstlisting}[style=lean,firstnumber=10]
theorem corollary31 (baseline : CCBaselineFamily)
    (ptas : ∀ k : ℕ, 2 ≤ k → ∀ slack : ℝ, 0 < slack → FixedKCCContract k slack (1 / 2)) :
    FixedKCorrelationClusteringStatement
\end{lstlisting}

\paragraph{Combined implementation theorem (proof omitted).}
\textit{Source:} \SourceId{DPCut/ImplementedApplications.lean}, line 154. \textit{Namespace:} \SourceId{DPCut.LSWApplications}.
\begin{lstlisting}[style=lean,firstnumber=154]
theorem corollary31 (hLSW : LSWStatement) (baseline : CCBaselineFamily)
    (ptas : ∀ k : ℕ, 2 ≤ k → ∀ slack : ℝ, 0 < slack → FixedKCCContract k slack (1/2))
    {ξ ε : ℝ} (hξ : 0 < ξ) (hε : 0 < ε) :
    ∃ C Cfail : ℝ, ∃ exponent : ℕ, 0 < C ∧ 0 < Cfail ∧
      ∀ k : ℕ, ∀ hk : 2 ≤ k, ImplementedCCGuarantee hLSW (baseline ε hε).toPolynomialGraphAlgorithm
        (ptas k hk (ξ/2) (half_pos hξ)).toPolynomialGraphAlgorithm C Cfail ε (1+ξ) exponent (fun _ => k)
\end{lstlisting}

\paragraph{Comparison note.}
The original Lean target records the statistical guarantee. The combined implementation theorem also provides a polynomial assigned-cost bound for the specified approximate program, using the formal LSW parameter and cited external algorithm contracts. These interfaces allow arbitrary kernels with cost zero; actual running times of the external algorithms and full reduction are not certified. It makes $\widetilde O(n)$ explicit as $C n(\log n)^d$ and high probability as failure at most $C_{\rm fail}/n^2$, with constants chosen before $n$. The error coefficient, failure coefficient, and logarithmic exponent are chosen after $\xi,\varepsilon$ and before $k$, matching the paper's additive bound independent of $k$. The theorem takes the cited private baseline and fixed-$k$ PTAS family as explicit external inputs; running-time constants may still depend on $k$.
\clearpage
\section{Theorem 32: Private Max-k-Cut}
\SourceId{thm:private-max-k-cut}\hfill\SourceId{STOC2027.tex}:3845--3870
\paragraph{Paper statement.}
\setcounter{theorem}{31}
\begin{theorem}[Private Max-$k$-Cut]
\label{thm:private-max-k-cut}
Suppose a polynomial-time nonprivate algorithm, on every nonnegative weighted graph $Q$, outputs
an at-most-$k$ partition $\Pi$ and, for some $0<\rho_k\leq1$, satisfies
\[
 \E[\operatorname{cut}_Q(\Pi)]
 \geq\rho_k\operatorname{OPT}^{\operatorname{Max-}k\operatorname{-Cut}}_Q.
\]
For every fixed $\xi>0$, constant $\varepsilon>0$, and
$\delta=n^{-10}$, there is a polynomial-time $(\varepsilon, \delta)$-edge-DP algorithm which, with
high probability, outputs an at-most-$k$ partition $\widehat\Pi$ satisfying
\[
 \operatorname{cut}_G(\widehat\Pi)
 \geq (\rho_k-\xi)\operatorname{OPT}^{\operatorname{Max-}k\operatorname{-Cut}}_G
   -\widetilde O_{\xi,\varepsilon}(n).
\]
The additive term is independent of $k$.  In particular, for Max-Cut, the algorithm of
Goemans and Williamson~\cite{GoemansWilliamson95} gives
\[
 \operatorname{cut}_G(\widehat\Pi)
 \geq (0.878-\xi)\operatorname{OPT}^{\operatorname{Max-Cut}}_G
   -\widetilde O_{\xi,\varepsilon}(n).
\]
For $k\geq3$, one may use the standard SDP ratio $\rho_k$ of Frieze and
Jerrum~\cite{FriezeJerrum97}.
\end{theorem}
\paragraph{Lean statement.}
\textit{Source:} \SourceId{DPCut/Applications.lean}, line 183.
\begin{lstlisting}[style=lean,firstnumber=183]
def PrivateMaxKStatement : Prop :=
  ∀ (ξ ε : ℝ), 0 < ξ → 0 < ε → ∃ C Cfail : ℝ, ∃ exponent : ℕ,
    0 < C ∧ 0 < Cfail ∧ ∀ (k : ℕ), 2 ≤ k → ∀ ρ : ℝ, 0 < ρ → ρ ≤ 1 →
    ∀ algorithm : (n : ℕ) → Kernel (Weights n) (Partition n),
      (∀ n, IsMarkovKernel (algorithm n)) →
      (∀ (n : ℕ), 2 ≤ n → ∀ w, Nonnegative w →
        (∀ᵐ π ∂algorithm n w, clusterCount π ≤ k) ∧
        ρ * maxKOptimum w k ≤ ∫ π, partitionCut w π ∂algorithm n w) →
      PrivateMaxKGuarantee C Cfail ε (ρ - ξ) exponent k
\end{lstlisting}

\textit{Source:} \SourceId{DPCut/Applications.lean}, line 124.
\begin{lstlisting}[style=lean,firstnumber=124]
def MaxKTransferStatement : Prop :=
  ∀ (n k : ℕ), 2 ≤ n → 2 ≤ k → ∀ (w h : Weights n) (η A ρ : ℝ) (π : Partition n),
    Nonnegative w → Nonnegative h → 0 ≤ η → η < 1 → 0 ≤ A → 0 ≤ ρ → ρ ≤ 1 →
    AllCuts w h η A → clusterCount π ≤ k → ρ * maxKOptimum h k ≤ partitionCut h π →
    ρ * (1 - η) / (1 + η) * maxKOptimum w k - 2 * (ρ + 1) / (1 + η) * A ≤ partitionCut w π
\end{lstlisting}

\paragraph{Combined implementation theorem (proof omitted).}
\textit{Source:} \SourceId{DPCut/ImplementedApplications.lean}, line 171. \textit{Namespace:} \SourceId{DPCut.LSWApplications}.
\begin{lstlisting}[style=lean,firstnumber=171]
theorem theorem32 (hLSW : LSWStatement) {ξ ε : ℝ} (hξ : 0 < ξ) (hε : 0 < ε) :
    ∃ C Cfail : ℝ, ∃ exponent : ℕ, 0 < C ∧ 0 < Cfail ∧
    ∀ k : ℕ, 2 ≤ k → ∀ ρ : ℝ, 0 < ρ → ρ ≤ 1 →
    ∀ alg : MaxKCutContract k ρ,
      ImplementedMaxKGuarantee hLSW alg.toPolynomialGraphAlgorithm C Cfail ε (ρ-ξ) exponent k
\end{lstlisting}

\paragraph{Comparison note.}
The original Lean target records the statistical guarantee. The combined implementation theorem also provides a polynomial assigned-cost bound for the specified approximate program, using the formal LSW parameter and cited external algorithm contracts. These interfaces allow arbitrary kernels with cost zero; actual running times of the external algorithms and full reduction are not certified. It makes $\widetilde O(n)$ explicit as $C n(\log n)^d$ and high probability as failure at most $C_{\rm fail}/n^2$, with constants chosen before $n$. Constants precede $k$ in the quantifiers, recording independence of the additive term from $k$. The second target is the deterministic transfer inequality used by the reduction; the Max-Cut specialization also appears in Corollary~\ref{cor:maxcut-intro}.
\clearpage
\section{Remark 33: Approximate-DP Max-Cut lower bound}
\SourceId{rem:maxcut-tightness}\hfill\SourceId{STOC2027.tex}:3914--3953
\paragraph{Paper statement.}
\setcounter{theorem}{32}
\begin{remark}[Approximate-DP Max-Cut lower bound]\label{rem:maxcut-tightness}
A direct-sum adaptation of the privacy comparison in~\cite{dharangutte2026price}
shows that the linear additive dependence on $n$ is necessary even for unweighted
Max-Cut under approximate differential privacy. Specifically, let $n\geq8$, $0\leq\varepsilon\leq1$,
$0\leq\delta,\beta\leq1/100$, and $0.868\leq\rho\leq1$.
If an $(\varepsilon,\delta)$-edge-DP algorithm $\mathcal M$ returns a cut of value
at least $\rho\operatorname{OPT}_G-B$, where $B\geq0$, with probability at least
$1-\beta$ on every unweighted $n$-vertex graph, then
$B\geq\lfloor n/8\rfloor/100=\Omega(n)$.
To see this, partition $8m$ vertices into $m=\lfloor n/8\rfloor$ public blocks
of size eight, and put one independently uniform random edge in each block;
any remaining vertices are isolated. Every such graph has optimum $m$.
Fix a block and the edges in all other blocks. Let $G_0$ leave this block empty,
and let $G_e=G_0+e$ for each of its $28$ possible edges $e$.
Let $E_e$ be the event that the output cut does not cut $e$.
Differential privacy gives
\[
 \Prb[\mathcal M(G_e)\in E_e]
 \geq e^{-\varepsilon}\bigl(\Prb[\mathcal M(G_0)\in E_e]-\delta\bigr).
\]
Every cut separates at most $16$ of the $28$ possible edges in the block, so
averaging over $e$ bounds its expected contribution under $\mathcal M(G_e)$ by
\[
 q:=1-e^{-\varepsilon}(3/7-\delta)
 \leq1-e^{-1}(3/7-1/100)<0.847.
\]
This holds for every choice of the other blocks. Averaging over all input edges
and summing over blocks therefore gives some graph $G$ with
$\E[\operatorname{cut}_G(\mathcal M(G))]\leq qm$.
The assumed utility guarantee and nonnegativity of cut values give
$\E[\operatorname{cut}_G(\mathcal M(G))]\geq(1-\beta)\rho m-B$.
Consequently,
\[
 B\geq\bigl((1-\beta)\rho-q\bigr)m
 >\bigl(0.99\cdot0.868-0.847\bigr)m>m/100.
\]
In particular, this applies to the ratio $0.878-\xi$ for $0<\xi\leq1/100$,
establishing the claimed tightness up to polylogarithmic factors in that range.
\end{remark}
\paragraph{Lean statement.}
\textit{Source:} \SourceId{DPCut/Applications.lean}, line 131.
\begin{lstlisting}[style=lean,firstnumber=131]
def MaxCutLowerBoundStatement : Prop :=
  ∀ (n : ℕ), 8 ≤ n → ∀ (ε δ β ρ B : ℝ),
    0 ≤ ε → ε ≤ 1 → 0 ≤ δ → δ ≤ 1 / 100 → 0 ≤ β → β ≤ 1 / 100 →
    (868 / 1000 : ℝ) ≤ ρ → ρ ≤ 1 → 0 ≤ B →
  ∀ M : Kernel (Weights n) (Finset (Fin n)), IsMarkovKernel M →
    DP (fun w w' => IsUnweighted w ∧ IsUnweighted w' ∧ Adjacent w w') M ε δ →
    (∀ w, IsUnweighted w → HighProbability (M w) β
      (fun S => ρ * sSup (Set.range (cut w)) - B ≤ cut w S)) →
    ((n / 8 : ℕ) : ℝ) / 100 ≤ B
\end{lstlisting}

\paragraph{Corresponding theorem declaration (proof omitted).}
\textit{Source:} \SourceId{DPCut/Remark33.lean}, line 71.
\begin{lstlisting}[style=lean,firstnumber=71]
theorem remark33 : MaxCutLowerBoundStatement
\end{lstlisting}

\paragraph{Comparison note.}
The Lean theorem proves the remark's lower bound $B\geq\lfloor n/8\rfloor/100$ with the same parameter ranges, including $\varepsilon=0$, $\delta=0$, and $\beta=0$. The remark's explanatory argument is retained in the paper excerpt. No external algorithmic hypothesis is needed.
\clearpage
\section{Theorem 34: Private bicriteria cut minimization}
\SourceId{thm:private-constrained-cut}\hfill\SourceId{STOC2027.tex}:3982--4012
\paragraph{Paper statement.}
\setcounter{theorem}{33}
\begin{theorem}[Private bicriteria cut minimization]
\label{thm:private-constrained-cut}
Fix $0<\eta\leq1/4$, $0<\varepsilon\leq1$, $0<\delta,\beta<1/10$,
$\rho\geq1$, and $0\leq\tau<1$.  Let
$\emptyset\ne\mathcal F,\mathcal F'\subseteq2^V$ be public families and
define
\[
 \operatorname{OPT}_{G,\mathcal F}
 :=\min_{S\in\mathcal F}w_G(S,V\setminus S).
\]
Suppose a polynomial-time nonprivate algorithm, on every nonnegative weighted graph, returns
$\widehat S\in\mathcal F'$ with probability at least $1-\tau$ and, on that event,
\[
 w(\widehat S,V\setminus\widehat S)
 \leq \rho\min_{S\in\mathcal F}w(S,V\setminus S).
\]
Then there is a polynomial-time $(\varepsilon, \delta)$-edge-DP algorithm which, with probability
at least $1-\beta-\tau$, returns $\widehat S\in\mathcal F'$ satisfying
\[
 w_G(\widehat S,V\setminus\widehat S)
 \leq \rho\frac{1+\eta}{1-\eta}\operatorname{OPT}_{G,\mathcal F}
   +\frac{\rho+1}{1-\eta}A_\eta.
\]
Here $A_\eta$ is the quantity in \eqref{eq:application-additive-error} for the displayed
$(\varepsilon, \delta,\beta)$.
Consequently, for every fixed $\xi>0$ and $\rho\geq1$, choosing
$\eta=\min\{1/4,\xi/(2\rho+\xi)\}$ gives
$(\rho+\xi)\operatorname{OPT}_{G,\mathcal F}
+\widetilde O_{\xi,\varepsilon}(\rho^{8}n)$ for constant $\varepsilon$ and
inverse-polynomial $\delta,\beta$.
\end{theorem}
\paragraph{Lean statement.}
\textit{Source:} \SourceId{DPCut/Applications.lean}, line 109.
\begin{lstlisting}[style=lean,firstnumber=109]
def ConstrainedCutStatement : Prop :=
  ∀ (n : ℕ), 2 ≤ n → ∀ (F F' : Set (Finset (Fin n))), F.Nonempty → F'.Nonempty →
  ∀ (η ε δ β ρ τ A : ℝ), 0 < η → η ≤ 1 / 4 → 0 < ε → ε ≤ 1 →
    0 < δ → δ < 1 / 10 → 0 < β → β < 1 / 10 → 1 ≤ ρ → 0 ≤ τ → τ < 1 → 0 ≤ A →
  ∀ (M : Kernel (Weights n) (Weights n)) (K : Kernel (Weights n) (Finset (Fin n))),
    IsMarkovKernel M → IsMarkovKernel K → GraphValued M → EdgeDP M ε δ →
    (∀ w, Nonnegative w → HighProbability (M w) β (fun h => AllCuts w h η A)) →
    (∀ h, Nonnegative h → HighProbability (K h) τ
      (fun S => S ∈ F' ∧ cut h S ≤ ρ * constrainedOptimum h F)) →
    EdgeDP (fun w => (M w).bind K) ε δ ∧ ∀ w, Nonnegative w →
      HighProbability ((M w).bind K) (β + τ)
        (fun S => S ∈ F' ∧ cut w S ≤ ρ * (1 + η) / (1 - η) * constrainedOptimum w F +
          (ρ + 1) / (1 - η) * A)
\end{lstlisting}

\paragraph{Combined implementation theorem (proof omitted).}
\textit{Source:} \SourceId{DPCut/ImplementedApplications.lean}, line 279. \textit{Namespace:} \SourceId{DPCut.LSWApplications}.
\begin{lstlisting}[style=lean,firstnumber=279]
theorem constrained_cut_general (hLSW : LSWStatement)
    (alg : PolynomialGraphAlgorithm (fun n => Finset (Fin n))) :
    ∃ Ctime : ℝ, 0 < Ctime ∧ ∃ degree : ℕ,
    ∀ n (hn : 2 ≤ n) (p : Parameters) (hp : p.Valid)
      (S S' : Set (Finset (Fin n))), S.Nonempty →
    ∀ ρ τ : ℝ, 0 ≤ ρ → 0 ≤ τ →
    (∀ h, Nonnegative h → HighProbability (alg.kernel n h) τ
      (fun T => T ∈ S' ∧ cut h T ≤ ρ*constrainedOptimum h S)) →
    let M := postprocessKernel (family hLSW) hn p hp (alg.kernel n)
    EdgeDP M p.epsilon p.delta ∧
    RealizedWithin M (ENNReal.ofReal (Ctime*runtimeScale n p^degree))
      (costedPostprocess (family hLSW) hn p hp (alg.run n)) ∧
    ∀ w, Nonnegative w → HighProbability (M w) (p.beta+τ)
      (fun T => T ∈ S' ∧ cut w T ≤ ρ*(1+p.gamma)/(1-p.gamma)*constrainedOptimum w S +
        (ρ+1)/(1-p.gamma)*((2:ℝ)^178*mainError n p))
\end{lstlisting}

\paragraph{Combined implementation theorem (proof omitted).}
\textit{Source:} \SourceId{DPCut/ImplementedApplications.lean}, line 313. \textit{Namespace:} \SourceId{DPCut.LSWApplications}.
\begin{lstlisting}[style=lean,firstnumber=313]
theorem theorem34 (hLSW : LSWStatement)
    (alg : PolynomialGraphAlgorithm (fun n => Finset (Fin n)))
    {η ε : ℝ} (hη : 0 < η) (hη4 : η ≤ 1/4) (hε : 0 < ε) :
    ∃ CA Ctime : ℝ, 0 < CA ∧ 0 < Ctime ∧ ∃ degree : ℕ,
    ∀ n (hn : 2 ≤ n) (S S' : Set (Finset (Fin n))), S.Nonempty →
    ∀ ρ τ : ℝ, 0 ≤ ρ → 0 ≤ τ →
    (∀ h, Nonnegative h → HighProbability (alg.kernel n h) τ
      (fun T => T ∈ S' ∧ cut h T ≤ ρ*constrainedOptimum h S)) →
    let p := applicationParameters n η ε
    let M := postprocessKernel (family hLSW) hn p (applicationParameters_valid hn hη hη4 hε) (alg.kernel n)
    EdgeDP M ε ((n:ℝ)⁻¹^10) ∧
    RealizedWithin M (ENNReal.ofReal (Ctime*(1+(n:ℝ))^degree))
      (costedPostprocess (family hLSW) hn p (applicationParameters_valid hn hη hη4 hε) (alg.run n)) ∧
    ∀ w, Nonnegative w → HighProbability (M w) ((n:ℝ)⁻¹^10+τ)
      (fun T => T ∈ S' ∧ cut w T ≤ ρ*(1+η)/(1-η)*constrainedOptimum w S +
        (ρ+1)/(1-η)*(CA*n*Real.log n^7))
\end{lstlisting}

\paragraph{Comparison note.}
The Lean statement is conditional on a release mechanism with additive error $A$ and on a nonprivate approximation kernel. It preserves the exact approximation factors and failure bound $\beta+\tau$. Instantiating the release with Theorem~\ref{thm:main} gives the paper's $A_\eta$. The combined general theorem instantiates the approximate LSW program and includes a uniform polynomial assigned-cost bound, subject to the cost-model limitation above. The application specialization additionally calibrates the error at $\delta=\beta=n^{-10}$.
\clearpage
\section{Lemma 35: Bicriteria balanced recursion}
\SourceId{lem:hc-bicriteria-recursion}\hfill\SourceId{STOC2027.tex}:4108--4120
\paragraph{Paper statement.}
\setcounter{theorem}{34}
\begin{lemma}[Bicriteria balanced recursion]\label{lem:hc-bicriteria-recursion}
Fix constants $0<b\leq1/3$ and $0<b'\leq1/3$, and let $a\geq1$.  Suppose an algorithm, on every
nonnegative weighted graph with vertex set $A$ of size at least $1/b$, returns a $b'$-balanced
cut $S$ satisfying
\[
 \phi_{G[A]}(S)\leq a\min\{\phi_{G[A]}(R):R\text{ is $b$-balanced}\}.
\]
Recursively applying this algorithm, and splitting smaller clusters as evenly as possible,
produces a hierarchical clustering tree $T$ with
\[
 \operatorname{hcost}_G(T)\leq O_b(a)\operatorname{OPT}^{\rm HC}_G.
\]
\end{lemma}
\paragraph{Lean statement.}
\textit{Source:} \SourceId{DPCut/Applications.lean}, line 65.
\begin{lstlisting}[style=lean,firstnumber=65]
def HCRecursionStatement : Prop :=
  ∀ b : ℝ, 0 < b → b ≤ 1 / 3 → ∃ C : ℝ, 0 < C ∧
    ∀ (b' a : ℝ), 0 < b' → b' ≤ 1 / 3 → 1 ≤ a →
    ∀ (n : ℕ), 2 ≤ n → ∀ (w : Weights n) (T : Hierarchy n), Nonnegative w →
      BalancedHierarchy b' T → RecursiveApproximation w b a T → hcCost w T ≤ C * a * hcOptimum w
\end{lstlisting}

\paragraph{Corresponding theorem declaration (proof omitted).}
\textit{Source:} \SourceId{DPCut/Lemma35.lean}, line 88.
\begin{lstlisting}[style=lean,firstnumber=88]
theorem lemma35 : HCRecursionStatement
\end{lstlisting}

\paragraph{Comparison note.}
Lean expresses the recursive algorithm's hypothesis as properties of its output tree: \texttt{BalancedHierarchy} and \texttt{RecursiveApproximation}. The approximation constant may depend on $b$, and is quantified before $b'$, $a$, and the graph.
\clearpage
\section{Theorem 36: Private hierarchical clustering}
\SourceId{thm:private-hierarchical-clustering}\hfill\SourceId{STOC2027.tex}:4171--4184
\paragraph{Paper statement.}
\setcounter{theorem}{35}
\begin{theorem}[Private hierarchical clustering]
\label{thm:private-hierarchical-clustering}
For every $0<\eta\leq1/4$, constant $0<\varepsilon\leq1$, and $\delta=n^{-10}$, there is a
polynomial-time $(\varepsilon, \delta)$-edge-DP algorithm that, with high probability, outputs a
$b'$-balanced hierarchical clustering tree $\widehat T$, for a universal constant $b'>0$,
satisfying
\[
 \operatorname{hcost}_G(\widehat T)
 \leq O(\sqrt{\log n})\operatorname{OPT}^{\rm HC}_G
  +O_\eta(A_\eta n\log^{3/2}n).
\]
In particular, for fixed $\eta$ the additive term is
$\widetilde O_{\eta,\varepsilon}(n^2)$.
\end{theorem}
\paragraph{Lean statement.}
\textit{Source:} \SourceId{DPCut/Applications.lean}, line 201.
\begin{lstlisting}[style=lean,firstnumber=201]
def PrivateHierarchyStatement : Prop :=
  ∃ b' Cratio CA Cfail : ℝ, 0 < b' ∧ b' ≤ 1 / 3 ∧ 0 < Cratio ∧ 0 < CA ∧ 0 < Cfail ∧
    ∀ η : ℝ, 0 < η → η ≤ 1 / 4 → ∃ Cadd : ℝ, 0 < Cadd ∧
    ∀ ε : ℝ, 0 < ε → ε ≤ 1 → ∀ (n : ℕ), 2 ≤ n →
    ∃ M : Kernel (Weights n) (Hierarchy n), IsMarkovKernel M ∧ EdgeDP M ε ((n : ℝ)⁻¹ ^ 10) ∧
      ∀ w, Nonnegative w → HighProbability (M w) (Cfail / (n : ℝ) ^ 2)
        (fun T => BalancedHierarchy b' T ∧ hcCost w T ≤
          Cratio * Real.sqrt (Real.log n) * hcOptimum w +
          Cadd * (CA * mainError n ⟨η, ε, (n : ℝ)⁻¹ ^ 10, (n : ℝ)⁻¹ ^ 10⟩) *
            n * (Real.log n) ^ (3 / 2 : ℝ))
\end{lstlisting}

\textit{Source:} \SourceId{DPCut/Applications.lean}, line 72.
\begin{lstlisting}[style=lean,firstnumber=72]
def HCTransferStatement : Prop :=
  ∀ (n : ℕ), 2 ≤ n → ∀ (w h : Weights n) (T : Hierarchy n) (η A : ℝ),
    Nonnegative w → Nonnegative h → IsHierarchy T → 0 ≤ η → 0 ≤ A → AllCuts w h η A →
    |hcCost h T - hcCost w T| ≤ η * hcCost w T +
      A / 2 * (n + ∑ cluster ∈ T.filter (fun cluster => 2 ≤ cluster.card), (cluster.card : ℝ))
\end{lstlisting}

\paragraph{Combined implementation theorem (proof omitted).}
\textit{Source:} \SourceId{DPCut/ImplementedApplications.lean}, line 356. \textit{Namespace:} \SourceId{DPCut.LSWApplications}.
\begin{lstlisting}[style=lean,firstnumber=356]
theorem theorem36 (hLSW : LSWStatement) (o : InducedARVContract) :
    ∃ b' Cratio CA Cfail : ℝ, 0 < b' ∧ b' ≤ 1/3 ∧ 0 < Cratio ∧ 0 < CA ∧ 0 < Cfail ∧
    ∀ η : ℝ, ∀ (hη : 0 < η) (hη4 : η ≤ 1/4), ∃ Cadd : ℝ, 0 < Cadd ∧
    ∀ ε : ℝ, ∀ hε : 0 < ε, ε ≤ 1 → ∃ Ctime : ℝ, 0 < Ctime ∧ ∃ degree : ℕ,
    ∀ n (hn : 2 ≤ n), ∃ M : Kernel (Weights n) (Hierarchy n), IsMarkovKernel M ∧
      RealizedWithin M (ENNReal.ofReal (Ctime*(1+(n:ℝ))^degree))
        (costedPrivateHierarchy hLSW o hn (applicationParameters n η ε)
          (applicationParameters_valid hn hη hη4 hε)) ∧
      EdgeDP M ε ((n:ℝ)⁻¹^10) ∧
      ∀ w, Nonnegative w → HighProbability (M w) (Cfail/(n:ℝ)^2)
        (fun T => BalancedHierarchy b' T ∧ hcCost w T ≤
          Cratio*Real.sqrt (Real.log n)*hcOptimum w +
          Cadd*(CA*mainError n ⟨η, ε, (n:ℝ)⁻¹^10, (n:ℝ)⁻¹^10⟩)*n*(Real.log n)^(3/2:ℝ))
\end{lstlisting}

\paragraph{Comparison note.}
The original Lean target records the statistical guarantee, with failure at most $C_{\rm fail}/n^2$. The combined theorem includes a polynomial realization using LSW and the weighted ARV contract. It evaluates the additive error at $\delta=\beta=n^{-10}$ and assumes $0<\varepsilon\leq1$, as in the paper. The additive coefficient may depend on $\eta$; the balance and ratio constants are universal. The second target is the deterministic hierarchy-transfer inequality used in the analysis.
\clearpage
\section{Lemma 37: Graph discrepancy}
\SourceId{lem:ekkl-discrepancy}\hfill\SourceId{STOC2027.tex}:4377--4388
\textit{Background result; not a contribution of this paper.}
\paragraph{Paper statement.}
\setcounter{theorem}{36}
\begin{lemma}[{\cite[Lemma~5.2]{EKKL20}}]
\label{lem:ekkl-discrepancy}
There is an absolute constant $0<\kappa\leq1$ such that the following holds.
Let $n\geq2$, $1\leq d\leq n/2$, and $\sigma\in[0,1]$.
Let $G,G'$ be simple unweighted graphs on the same $n$ vertices, with every
vertex having degree in $[d/2,2d]$ in each graph. If
$|E(G)\mathbin{\triangle}E(G')|\geq\sigma nd$, then
\[
 \max_{\substack{S,T\subseteq V\\S\cap T=\emptyset}}
       |w_G(S,T)-w_{G'}(S,T)|\geq\kappa\sigma n\sqrt d.
\]
\end{lemma}
\paragraph{Lean statement.}
\textit{Source:} \SourceId{DPCut/Statements.lean}, line 224.
\begin{lstlisting}[style=lean,firstnumber=224]
def EKKLDiscrepancyStatement : Prop :=
  ∃ κ : ℝ, 0 < κ ∧ κ ≤ 1 ∧ ∀ (n : ℕ), 2 ≤ n →
  ∀ (d σ : ℝ), 1 ≤ d → d ≤ n / 2 → 0 ≤ σ → σ ≤ 1 →
  ∀ (w w' : Weights n), IsUnweighted w → IsUnweighted w' →
    (∀ v, d / 2 ≤ vertexDegree w v ∧ vertexDegree w v ≤ 2 * d) →
    (∀ v, d / 2 ≤ vertexDegree w' v ∧ vertexDegree w' v ≤ 2 * d) →
    σ * n * d ≤ l1 (w - w') →
    ∃ S T : Finset (Fin n), Disjoint S T ∧
      κ * σ * n * Real.sqrt d ≤ |between w S T - between w' S T|
\end{lstlisting}

\paragraph{Comparison note.}
Simple graphs are encoded by binary edge weights. The maximum-over-disjoint-sets conclusion is expressed by the existence of disjoint finite sets attaining the required lower bound.
\clearpage
\section{Theorem 38: Lower bound for relative cut release}
\SourceId{thm:lower_bound}\hfill\SourceId{STOC2027.tex}:4398--4418
\paragraph{Paper statement.}
\setcounter{theorem}{37}
\begin{theorem}[Lower Bound for Relative Cut Release]
\label{thm:lower_bound}
Let $n\geq33$, $\varepsilon\in(0,1)$, $0\leq\delta\leq\varepsilon/100$,
$\beta\in (0,1/40]$, and $\gamma \in [0, 1)$.
Suppose that an $(\varepsilon,\delta)$-differentially private mechanism
$\mathcal M$, on every nonnegative weighted graph $G$ on a fixed set $V$
of $n$ vertices, outputs a nonnegative weighted graph $H$ such that,
with probability at least $1-\beta$, simultaneously for all $S\subseteq V$,
\[
 \left|w_G(S, V\setminus S) - w_H(S, V\setminus S)\right| \leq \gamma w_G(S, V\setminus S) + A.
\]
Then the additive error must satisfy
\[
 A=\Omega\left(
   (1-\gamma)\cdot\min\left\{\frac{n^{3/2}}{\varepsilon},
                    \frac{n}{\varepsilon\gamma}\right\}
 \right).
\]
where the implicit constant is absolute and $1/\gamma$ is interpreted as
$+\infty$ when $\gamma=0$.
\end{theorem}
\paragraph{Lean statement.}
\textit{Source:} \SourceId{DPCut/Statements.lean}, line 212.
\begin{lstlisting}[style=lean,firstnumber=212]
def LowerBoundStatement : Prop :=
  ∃ c : ℝ, 0 < c ∧ ∀ (n : ℕ), 33 ≤ n → ∀ (ε δ β γ A : ℝ),
  0 < ε → ε < 1 → 0 ≤ δ → δ ≤ ε / 100 → 0 < β → β ≤ 1 / 40 →
  0 ≤ γ → γ < 1 → 0 ≤ A →
  ∀ M : Kernel (Weights n) (Weights n), IsMarkovKernel M → GraphValued M → EdgeDP M ε δ →
    (∀ w : Weights n, Nonnegative w → HighProbability (M w) β (fun h => AllCuts w h γ A)) →
    c * lowerScale n ε γ ≤ A
\end{lstlisting}

\paragraph{Corresponding theorem declaration (proof omitted).}
\textit{Source:} \SourceId{DPCut/Theorem38.lean}, line 65.
\begin{lstlisting}[style=lean,firstnumber=65]
theorem theorem38 : LowerBoundStatement
\end{lstlisting}

\paragraph{Comparison note.}
\texttt{lowerScale} explicitly treats $\gamma=0$ by using $n^{3/2}/\varepsilon$ for the minimum. The universal lower-bound constant is quantified before all parameters; the displayed proof chooses $c=2^{-30}$.
\clearpage
\section{Notation appendix: definitions used by the Lean statements}
\subsection{Combined application implementation guarantees}
\textit{Source:} \SourceId{DPCut/ImplementedApplications.lean}, line 19.
\begin{lstlisting}[style=lean,firstnumber=19]
def RealizedWithin {n : ℕ} {Y : Type} [MeasurableSpace Y]
    (M : Kernel (Weights n) Y) (bound : ℝ≥0∞) (run : Weights n → CostedLaw Y) : Prop :=
  ∀ w, (run w).law = M w ∧ (run w).cost ≤ bound
\end{lstlisting}

\textit{Source:} \SourceId{DPCut/ImplementedApplications.lean}, line 24.
\begin{lstlisting}[style=lean,firstnumber=24]
def ImplementedCCGuarantee (hLSW : LSWStatement) (base alg : PolynomialGraphAlgorithm Partition)
    (C Cfail ε ratio : ℝ) (exponent : ℕ) (clusters : ℕ → ℕ) : Prop :=
  ∃ Ctime : ℝ, 0 < Ctime ∧ ∃ degree : ℕ,
  ∀ n (hn : 2 ≤ n), ∃ (p : Parameters) (hp : p.Valid) (t : ℕ) (seed : Partition n)
    (M : Kernel (Weights n) (Partition n)), IsMarkovKernel M ∧
    RealizedWithin M (ENNReal.ofReal (Ctime * (1+(n:ℝ))^degree))
      (LSWApplications.costedCC (LSWApplications.family hLSW) hn p hp (base.run n)
        (costedCCSolver (alg.run n) (clusters n) t seed)) ∧
    DP (fun w w' => IsUnweighted w ∧ IsUnweighted w' ∧ Adjacent w w') M ε ((n:ℝ)⁻¹^10) ∧
    ∀ w, IsUnweighted w → HighProbability (M w) (Cfail/(n:ℝ)^2)
      (fun π => clusterCount π ≤ clusters n ∧
        ccCost w π ≤ ratio * ccOptimumK w (clusters n) + C*n*Real.log n^exponent)
\end{lstlisting}

\textit{Source:} \SourceId{DPCut/ImplementedApplications.lean}, line 38.
\begin{lstlisting}[style=lean,firstnumber=38]
def ImplementedMaxKGuarantee (hLSW : LSWStatement) (alg : PolynomialGraphAlgorithm Partition)
    (C Cfail ε ratio : ℝ) (exponent k : ℕ) : Prop :=
  ∃ Ctime : ℝ, 0 < Ctime ∧ ∃ degree : ℕ,
  ∀ n (hn : 2 ≤ n), ∃ (p : Parameters) (hp : p.Valid) (t : ℕ) (seed : Partition n)
    (M : Kernel (Weights n) (Partition n)), IsMarkovKernel M ∧
    RealizedWithin M (ENNReal.ofReal (Ctime * (1+(n:ℝ))^degree))
      (LSWApplications.costedPostprocess (LSWApplications.family hLSW) hn p hp
        (costedMaxKSolver (alg.run n) t seed)) ∧
    EdgeDP M ε ((n:ℝ)⁻¹^10) ∧
    ∀ w, Nonnegative w → HighProbability (M w) (Cfail/(n:ℝ)^2)
      (fun π => clusterCount π ≤ k ∧
        ratio * maxKOptimum w k - C*n*Real.log n^exponent ≤ partitionCut w π)
\end{lstlisting}

\subsection{Graph notation and probability}
\textit{Source:} \SourceId{DPCut/Definitions.lean}, line 36.
\begin{lstlisting}[style=lean,firstnumber=36]
abbrev Edge (n : ℕ) := {p : Fin n × Fin n // p.1 < p.2}
\end{lstlisting}

\textit{Source:} \SourceId{DPCut/Definitions.lean}, line 37.
\begin{lstlisting}[style=lean,firstnumber=37]
abbrev Weights (n : ℕ) := Edge n → ℝ
\end{lstlisting}

\textit{Source:} \SourceId{DPCut/Definitions.lean}, line 38.
\begin{lstlisting}[style=lean,firstnumber=38]
abbrev Mat (n : ℕ) := Matrix (Fin n) (Fin n) ℝ
\end{lstlisting}

\textit{Source:} \SourceId{DPCut/Definitions.lean}, line 40.
\begin{lstlisting}[style=lean,firstnumber=40]
def Nonnegative {n : ℕ} (w : Weights n) : Prop := ∀ e, 0 ≤ w e
\end{lstlisting}

\textit{Source:} \SourceId{DPCut/Definitions.lean}, line 41.
\begin{lstlisting}[style=lean,firstnumber=41]
def Bounded {n : ℕ} (U : ℝ) (w : Weights n) : Prop := ∀ e, 0 ≤ w e ∧ w e ≤ U
\end{lstlisting}

\textit{Source:} \SourceId{DPCut/Definitions.lean}, line 42.
\begin{lstlisting}[style=lean,firstnumber=42]
def IsMask {n : ℕ} (r : Weights n) : Prop := Bounded 1 r
\end{lstlisting}

\textit{Source:} \SourceId{DPCut/Definitions.lean}, line 43.
\begin{lstlisting}[style=lean,firstnumber=43]
def l1 {n : ℕ} (w : Weights n) : ℝ := ∑ e, |w e|
\end{lstlisting}

\textit{Source:} \SourceId{DPCut/Definitions.lean}, line 44.
\begin{lstlisting}[style=lean,firstnumber=44]
def mass {n : ℕ} (w : Weights n) : ℝ := ∑ e, w e
\end{lstlisting}

\textit{Source:} \SourceId{DPCut/Definitions.lean}, line 45.
\begin{lstlisting}[style=lean,firstnumber=45]
def pairing {n : ℕ} (w f : Weights n) : ℝ := ∑ e, w e * f e
\end{lstlisting}

\textit{Source:} \SourceId{DPCut/Definitions.lean}, line 46.
\begin{lstlisting}[style=lean,firstnumber=46]
def hadamard {n : ℕ} (w f : Weights n) : Weights n := fun e => w e * f e
\end{lstlisting}

\textit{Source:} \SourceId{DPCut/Definitions.lean}, line 47.
\begin{lstlisting}[style=lean,firstnumber=47]
def Adjacent {n : ℕ} (w w' : Weights n) : Prop := l1 (w - w') ≤ 1
\end{lstlisting}

\textit{Source:} \SourceId{DPCut/Definitions.lean}, line 49.
\begin{lstlisting}[style=lean,firstnumber=49]
def cutQuery {n : ℕ} (S : Finset (Fin n)) : Weights n :=
  fun e => if (e.val.1 ∈ S) ↔ (e.val.2 ∈ S) then 0 else 1
\end{lstlisting}

\textit{Source:} \SourceId{DPCut/Definitions.lean}, line 51.
\begin{lstlisting}[style=lean,firstnumber=51]
def cut {n : ℕ} (w : Weights n) (S : Finset (Fin n)) : ℝ := pairing w (cutQuery S)
\end{lstlisting}

\textit{Source:} \SourceId{DPCut/Definitions.lean}, line 52.
\begin{lstlisting}[style=lean,firstnumber=52]
def between {n : ℕ} (w : Weights n) (S T : Finset (Fin n)) : ℝ :=
  ∑ e, if (e.val.1 ∈ S ∧ e.val.2 ∈ T) ∨ (e.val.2 ∈ S ∧ e.val.1 ∈ T)
    then w e else 0
\end{lstlisting}

\textit{Source:} \SourceId{DPCut/Definitions.lean}, line 55.
\begin{lstlisting}[style=lean,firstnumber=55]
def AllCuts {n : ℕ} (w h : Weights n) (γ A : ℝ) : Prop :=
  ∀ S : Finset (Fin n), |cut w S - cut h S| ≤ γ * cut w S + A
\end{lstlisting}

\textit{Source:} \SourceId{DPCut/Definitions.lean}, line 59.
\begin{lstlisting}[style=lean,firstnumber=59]
def prob {α : Type*} [MeasurableSpace α] (μ : Measure α) (s : Set α) : ℝ := (μ s).toReal
\end{lstlisting}

\textit{Source:} \SourceId{DPCut/Definitions.lean}, line 60.
\begin{lstlisting}[style=lean,firstnumber=60]
def HighProbability {α : Type*} [MeasurableSpace α]
    (μ : Measure α) (β : ℝ) (P : α → Prop) : Prop :=
  MeasurableSet {x | P x} ∧ 1 - β ≤ prob μ {x | P x}
\end{lstlisting}

\textit{Source:} \SourceId{DPCut/Definitions.lean}, line 63.
\begin{lstlisting}[style=lean,firstnumber=63]
def DP {α β : Type*} [MeasurableSpace β] (adj : α → α → Prop)
    (M : α → Measure β) (ε δ : ℝ) : Prop :=
  ∀ x x', adj x x' → ∀ s, MeasurableSet s →
    M x s ≤ ENNReal.ofReal (Real.exp ε) * M x' s + ENNReal.ofReal δ
\end{lstlisting}

\textit{Source:} \SourceId{DPCut/Definitions.lean}, line 67.
\begin{lstlisting}[style=lean,firstnumber=67]
def EdgeDP {n : ℕ} {α : Type*} [MeasurableSpace α]
    (M : Weights n → Measure α) (ε δ : ℝ) : Prop :=
  DP (fun w w' => Nonnegative w ∧ Nonnegative w' ∧ Adjacent w w') M ε δ
\end{lstlisting}

\textit{Source:} \SourceId{DPCut/Definitions.lean}, line 70.
\begin{lstlisting}[style=lean,firstnumber=70]
def TVClose {α : Type*} [MeasurableSpace α] (μ ν : Measure α) (η : ℝ) : Prop :=
  ∀ s, MeasurableSet s → |prob μ s - prob ν s| ≤ η
\end{lstlisting}

\textit{Source:} \SourceId{DPCut/Definitions.lean}, line 72.
\begin{lstlisting}[style=lean,firstnumber=72]
def GraphValued {n : ℕ} (M : Weights n → Measure (Weights n)) : Prop :=
  ∀ w, Nonnegative w → ∀ᵐ h ∂M w, Nonnegative h
\end{lstlisting}

\subsection{Parameters and error expressions}
\textit{Source:} \SourceId{DPCut/Definitions.lean}, line 75.
\begin{lstlisting}[style=lean,firstnumber=75]
structure Parameters where
  gamma : ℝ
  epsilon : ℝ
  delta : ℝ
  beta : ℝ
\end{lstlisting}

\textit{Source:} \SourceId{DPCut/Definitions.lean}, line 81.
\begin{lstlisting}[style=lean,firstnumber=81]
def Parameters.Valid (p : Parameters) : Prop :=
  0 < p.gamma ∧ p.gamma ≤ 1 / 4 ∧ 0 < p.epsilon ∧ p.epsilon ≤ 1 ∧
  0 < p.delta ∧ p.delta < 1 / 10 ∧ 0 < p.beta ∧ p.beta < 1 / 10
\end{lstlisting}

\textit{Source:} \SourceId{DPCut/Definitions.lean}, line 84.
\begin{lstlisting}[style=lean,firstnumber=84]
def Parameters.PrivacyValid (p : Parameters) : Prop :=
  0 < p.gamma ∧ p.gamma < 1 ∧ 0 < p.epsilon ∧ p.epsilon ≤ 1 ∧
  0 < p.delta ∧ p.delta < 1 ∧ 0 < p.beta ∧ p.beta < 1
\end{lstlisting}

\textit{Source:} \SourceId{DPCut/Definitions.lean}, line 87.
\begin{lstlisting}[style=lean,firstnumber=87]
def logFactor (n : ℕ) (U : ℝ) (p : Parameters) : ℝ :=
  Real.log (n * U / (p.delta * p.beta * p.gamma * p.epsilon))
\end{lstlisting}

\textit{Source:} \SourceId{DPCut/Definitions.lean}, line 89.
\begin{lstlisting}[style=lean,firstnumber=89]
def mainError (n : ℕ) (p : Parameters) : ℝ :=
  n / (p.gamma * p.epsilon) * logFactor n 1 p ^ 7
\end{lstlisting}

\textit{Source:} \SourceId{DPCut/Definitions.lean}, line 91.
\begin{lstlisting}[style=lean,firstnumber=91]
def boundedError (n : ℕ) (U : ℝ) (p : Parameters) : ℝ :=
  n * U * logFactor n U p + n / (p.gamma * p.epsilon) * logFactor n U p ^ 7
\end{lstlisting}

\textit{Source:} \SourceId{DPCut/Definitions.lean}, line 201.
\begin{lstlisting}[style=lean,firstnumber=201]
def lowerScale (n : ℕ) (ε γ : ℝ) : ℝ :=
  (1 - γ) * if γ = 0 then n * Real.sqrt n / ε
    else min (n * Real.sqrt n / ε) (n / (ε * γ))
\end{lstlisting}

\textit{Source:} \SourceId{DPCut/Definitions.lean}, line 193.
\begin{lstlisting}[style=lean,firstnumber=193]
def innerHorizon (n : ℕ) (γ : ℝ) : ℕ :=
  1 + ⌈64 * (1 + γ) ^ 2 * Real.log n / γ ^ 2⌉₊
\end{lstlisting}

\textit{Source:} \SourceId{DPCut/Definitions.lean}, line 195.
\begin{lstlisting}[style=lean,firstnumber=195]
def outerHorizon (n : ℕ) (U C : ℝ) : ℕ := ⌈C * Real.log (n * U)⌉₊
\end{lstlisting}

\textit{Source:} \SourceId{DPCut/Definitions.lean}, line 196.
\begin{lstlisting}[style=lean,firstnumber=196]
def stepSize (γ : ℝ) : ℝ := γ / (8 * (1 + γ))
\end{lstlisting}

\textit{Source:} \SourceId{DPCut/Definitions.lean}, line 197.
\begin{lstlisting}[style=lean,firstnumber=197]
def stoppingThreshold (n : ℕ) (U C : ℝ) (p : Parameters) : ℝ :=
  C * (n * U * logFactor n U p + n / (p.gamma * p.epsilon) * logFactor n U p ^ 6)
\end{lstlisting}

\subsection{Parameter range for standalone iteration bounds}
\textit{Source:} \SourceId{DPCut/Utility.lean}, line 13.
\begin{lstlisting}[style=lean,firstnumber=13]
def Parameters.IterationValid (p : Parameters) : Prop :=
  0 < p.gamma ∧ p.gamma ≤ 1 / 4 ∧ 0 < p.epsilon ∧ p.epsilon ≤ 1 ∧
  0 < p.delta ∧ p.delta < 1 ∧ 0 < p.beta ∧ p.beta < 1
\end{lstlisting}

\subsection{SDP geometry and covariance distance}
\textit{Source:} \SourceId{DPCut/Definitions.lean}, line 96.
\begin{lstlisting}[style=lean,firstnumber=96]
def feasible (n : ℕ) : Set (Mat n) :=
  {X | X.IsHermitian ∧ (∀ i, X i i = 1) ∧ (X - (1 / (n : ℝ)) • 1).PosSemidef}
\end{lstlisting}

\textit{Source:} \SourceId{DPCut/Definitions.lean}, line 100.
\begin{lstlisting}[style=lean,firstnumber=100]
def fractionalQuery {n : ℕ} (X : Mat n) : Weights n := fun e => (1 - X e.val.1 e.val.2) / 2
\end{lstlisting}

\textit{Source:} \SourceId{DPCut/Definitions.lean}, line 101.
\begin{lstlisting}[style=lean,firstnumber=101]
def edgeVector {n : ℕ} (e : Edge n) : Fin n → ℝ :=
  fun i => (if i = e.val.1 then 1 else 0) - (if i = e.val.2 then 1 else 0)
\end{lstlisting}

\textit{Source:} \SourceId{DPCut/Definitions.lean}, line 104.
\begin{lstlisting}[style=lean,firstnumber=104]
def objectiveMatrix {n : ℕ} (D : Weights n) : Mat n :=
  ∑ e, (D e / 4) • Matrix.vecMulVec (edgeVector e) (edgeVector e)
\end{lstlisting}

\textit{Source:} \SourceId{DPCut/Definitions.lean}, line 106.
\begin{lstlisting}[style=lean,firstnumber=106]
def matrixPairing {n : ℕ} (A X : Mat n) : ℝ := ∑ i, ∑ j, A i j * X i j
\end{lstlisting}

\textit{Source:} \SourceId{DPCut/Definitions.lean}, line 107.
\begin{lstlisting}[style=lean,firstnumber=107]
def objective {n : ℕ} (lam : ℝ) (D : Weights n) (X : Mat n) : ℝ :=
  pairing D (fractionalQuery X) + lam * Real.log X.det
\end{lstlisting}

\textit{Source:} \SourceId{DPCut/Definitions.lean}, line 109.
\begin{lstlisting}[style=lean,firstnumber=109]
def IsOptimizer {n : ℕ} (lam : ℝ) (D : Weights n) (X : Mat n) : Prop :=
  X ∈ feasible n ∧ ∀ Y ∈ feasible n, objective lam D Y ≤ objective lam D X
\end{lstlisting}

\textit{Source:} \SourceId{DPCut/Definitions.lean}, line 113.
\begin{lstlisting}[style=lean,firstnumber=113]
def optimum {n : ℕ} (lam : ℝ) (D : Weights n) : ℝ :=
  sSup (objective lam D '' feasible n)
\end{lstlisting}

\textit{Source:} \SourceId{DPCut/Definitions.lean}, line 118.
\begin{lstlisting}[style=lean,firstnumber=118]
def optimizer {n : ℕ} (lam : ℝ) (D : Weights n) : Mat n :=
  if h : ∃ X, IsOptimizer lam D X then Classical.choose h else 1
\end{lstlisting}

\textit{Source:} \SourceId{DPCut/Definitions.lean}, line 125.
\begin{lstlisting}[style=lean,firstnumber=125]
def frobenius {n : ℕ} (X : Mat n) : ℝ :=
  @norm (Mat n) Matrix.frobeniusNormedAddCommGroup.toNorm X
\end{lstlisting}

\textit{Source:} \SourceId{DPCut/Definitions.lean}, line 127.
\begin{lstlisting}[style=lean,firstnumber=127]
def relativeDifference {n : ℕ} (X Y : Mat n) : Mat n :=
  (CFC.sqrt X)⁻¹ * (Y - X) * (CFC.sqrt Y)⁻¹
\end{lstlisting}

\textit{Source:} \SourceId{DPCut/Definitions.lean}, line 129.
\begin{lstlisting}[style=lean,firstnumber=129]
def covarianceDistance {n : ℕ} (X Y : Mat n) : ℝ := frobenius (relativeDifference X Y)
\end{lstlisting}

\subsection{Matrix optimizer predicate}
\textit{Source:} \SourceId{DPCut/Statements.lean}, line 26.
\begin{lstlisting}[style=lean,firstnumber=26]
def IsMatrixOptimizer {n : ℕ} (lam : ℝ) (A X : Mat n) : Prop :=
  X ∈ feasible n ∧ ∀ Y ∈ feasible n,
    matrixPairing A Y + lam * Real.log Y.det ≤ matrixPairing A X + lam * Real.log X.det
\end{lstlisting}

\subsection{Gaussian laws and witnesses}
\textit{Source:} \SourceId{DPCut/Definitions.lean}, line 132.
\begin{lstlisting}[style=lean,firstnumber=132]
def standardGaussian (n : ℕ) : Measure (Fin n → ℝ) :=
  Measure.pi (fun _ => gaussianReal 0 1)
\end{lstlisting}

\textit{Source:} \SourceId{DPCut/Definitions.lean}, line 134.
\begin{lstlisting}[style=lean,firstnumber=134]
def gaussian {n : ℕ} (X : Mat n) : Measure (Fin n → ℝ) :=
  (standardGaussian n).map (fun z => (CFC.sqrt X).mulVec z)
\end{lstlisting}

\textit{Source:} \SourceId{DPCut/Definitions.lean}, line 136.
\begin{lstlisting}[style=lean,firstnumber=136]
def gaussianBatch {n : ℕ} (X : Mat n) (k : ℕ) : Measure (Fin k → Fin n → ℝ) :=
  Measure.pi (fun _ : Fin k => gaussian X)
\end{lstlisting}

\textit{Source:} \SourceId{DPCut/Definitions.lean}, line 138.
\begin{lstlisting}[style=lean,firstnumber=138]
def gaussianLogDensity {n : ℕ} (X : Mat n) (z : Fin n → ℝ) : ℝ :=
  -(n : ℝ) / 2 * Real.log (2 * Real.pi) - Real.log X.det / 2 -
    (∑ i, z i * (X⁻¹.mulVec z) i) / 2
\end{lstlisting}

\textit{Source:} \SourceId{DPCut/Definitions.lean}, line 141.
\begin{lstlisting}[style=lean,firstnumber=141]
def batchLogRatio {n k : ℕ} (X Y : Mat n) (z : Fin k → Fin n → ℝ) : ℝ :=
  ∑ j, (gaussianLogDensity X (z j) - gaussianLogDensity Y (z j))
\end{lstlisting}

\textit{Source:} \SourceId{DPCut/Definitions.lean}, line 143.
\begin{lstlisting}[style=lean,firstnumber=143]
def rawWitness {n k : ℕ} (z : Fin k → Fin n → ℝ) : Weights n :=
  fun e => (∑ j, (z j e.val.1 - z j e.val.2) ^ 2 / 4) / k
\end{lstlisting}

\textit{Source:} \SourceId{DPCut/Definitions.lean}, line 145.
\begin{lstlisting}[style=lean,firstnumber=145]
def clippedWitness {n k : ℕ} (z : Fin k → Fin n → ℝ) : Weights n :=
  fun e => min (rawWitness z e / 2) 1
\end{lstlisting}

\textit{Source:} \SourceId{DPCut/Definitions.lean}, line 147.
\begin{lstlisting}[style=lean,firstnumber=147]
def ChiGood {n k : ℕ} (X : Mat n) (z : Fin k → Fin n → ℝ) : Prop :=
  ∀ e, fractionalQuery X e / 4 ≤ clippedWitness z e ∧
    clippedWitness z e ≤ fractionalQuery X e
\end{lstlisting}

\subsection{Preprocessing and aggregate guarantees}
\textit{Source:} \SourceId{DPCut/Definitions.lean}, line 155.
\begin{lstlisting}[style=lean,firstnumber=155]
def laplaceDensity (b x : ℝ) : ℝ := (2 * b)⁻¹ * Real.exp (-|x| / b)
\end{lstlisting}

\textit{Source:} \SourceId{DPCut/Definitions.lean}, line 156.
\begin{lstlisting}[style=lean,firstnumber=156]
def laplace (b : ℝ) : Measure ℝ := volume.withDensity (fun x => ENNReal.ofReal (laplaceDensity b x))
\end{lstlisting}

\textit{Source:} \SourceId{DPCut/Definitions.lean}, line 157.
\begin{lstlisting}[style=lean,firstnumber=157]
def clipWeights {n : ℕ} (U : ℝ) (w : Weights n) : Weights n := fun e => min (w e) U
\end{lstlisting}

\textit{Source:} \SourceId{DPCut/Definitions.lean}, line 158.
\begin{lstlisting}[style=lean,firstnumber=158]
def excess {n : ℕ} (U : ℝ) (w : Weights n) : Weights n := fun e => max (w e - U) 0
\end{lstlisting}

\textit{Source:} \SourceId{DPCut/Definitions.lean}, line 159.
\begin{lstlisting}[style=lean,firstnumber=159]
def clipThreshold (n : ℕ) (ε₀ β₀ γ : ℝ) : ℝ := 8 / (ε₀ * γ) * Real.log (n / β₀)
\end{lstlisting}

\textit{Source:} \SourceId{DPCut/Definitions.lean}, line 160.
\begin{lstlisting}[style=lean,firstnumber=160]
def noiseShift (n : ℕ) (ε₀ β₀ : ℝ) : ℝ := 2 / ε₀ * Real.log (n / β₀)
\end{lstlisting}

\textit{Source:} \SourceId{DPCut/Definitions.lean}, line 161.
\begin{lstlisting}[style=lean,firstnumber=161]
def noisyExcess {n : ℕ} (U V : ℝ) (w noise : Weights n) : Weights n :=
  fun e => max 0 (excess U w e + noise e - V)
\end{lstlisting}

\textit{Source:} \SourceId{DPCut/Definitions.lean}, line 163.
\begin{lstlisting}[style=lean,firstnumber=163]
def preprocessLaw {n : ℕ} (ε₀ β₀ γ : ℝ) (w : Weights n) : Measure (Weights n) :=
  (Measure.pi (fun _ : Edge n => laplace (1 / ε₀))).map
    (noisyExcess (clipThreshold n ε₀ β₀ γ) (noiseShift n ε₀ β₀) w)
\end{lstlisting}

\textit{Source:} \SourceId{DPCut/Definitions.lean}, line 167.
\begin{lstlisting}[style=lean,firstnumber=167]
structure Aggregate (n : ℕ) where
  plus : Weights n
  minus : Weights n
  residual : Weights n
  deriving Inhabited
\end{lstlisting}

\textit{Source:} \SourceId{DPCut/Definitions.lean}, line 173.
\begin{lstlisting}[style=lean,firstnumber=173]
def Aggregate.Partition {n : ℕ} (a : Aggregate n) : Prop :=
  IsMask a.plus ∧ IsMask a.minus ∧ IsMask a.residual ∧
  ∀ e, a.plus e + a.minus e + a.residual e = 1
\end{lstlisting}

\textit{Source:} \SourceId{DPCut/Definitions.lean}, line 176.
\begin{lstlisting}[style=lean,firstnumber=176]
def Aggregate.Accurate {n : ℕ} (a : Aggregate n) (w h : Weights n) (γ B : ℝ) : Prop :=
  pairing h a.plus ≥ (1 + γ) * pairing w a.plus ∧
  pairing h a.minus ≤ (1 - γ) * pairing w a.minus ∧
  AllCuts (hadamard w a.residual) (hadamard h a.residual) γ B
\end{lstlisting}

\textit{Source:} \SourceId{DPCut/Definitions.lean}, line 180.
\begin{lstlisting}[style=lean,firstnumber=180]
def Aggregate.Accepts {n : ℕ} (a : Aggregate n) (h : Weights n) : Prop :=
  max (pairing h a.plus) (pairing h a.minus) ≤ pairing h a.residual
\end{lstlisting}

\textit{Source:} \SourceId{DPCut/Definitions.lean}, line 182.
\begin{lstlisting}[style=lean,firstnumber=182]
def mwUpdate {n : ℕ} (M θ η : ℝ) (h g : Weights n) : Weights n :=
  fun e => M * (h e * Real.exp (-θ * η * g e)) /
    (∑ f, h f * Real.exp (-θ * η * g f))
\end{lstlisting}

\textit{Source:} \SourceId{DPCut/Definitions.lean}, line 185.
\begin{lstlisting}[style=lean,firstnumber=185]
def initialEstimate {n : ℕ} (M : ℝ) (a : Weights n) : Weights n := fun e => M * a e / mass a
\end{lstlisting}

\textit{Source:} \SourceId{DPCut/Definitions.lean}, line 186.
\begin{lstlisting}[style=lean,firstnumber=186]
def residualUpdate {n : ℕ} (a r : Weights n) : Weights n := fun e => a e * (1 - r e)
\end{lstlisting}

\subsection{Privacy-composition notation}
\textit{Source:} \SourceId{DPCut/Privacy.lean}, line 12.
\begin{lstlisting}[style=lean,firstnumber=12]
def Indistinguishable {α : Type*} [MeasurableSpace α]
    (μ ν : Measure α) (ε δ : ℝ) : Prop :=
  (∀ s, MeasurableSet s → μ s ≤ ENNReal.ofReal (Real.exp ε) * ν s + ENNReal.ofReal δ) ∧
  (∀ s, MeasurableSet s → ν s ≤ ENNReal.ofReal (Real.exp ε) * μ s + ENNReal.ofReal δ)
\end{lstlisting}

\textit{Source:} \SourceId{DPCut/Privacy.lean}, line 24.
\begin{lstlisting}[style=lean,firstnumber=24]
def laplaceMechanism {X : Type*} (k : ℕ) (f : X → Fin k → ℝ) (Δ ε : ℝ)
    (x : X) : Measure (Fin k → ℝ) :=
  if Δ = 0 then Measure.dirac (f x)
  else (Measure.pi (fun _ : Fin k => laplace (Δ / ε))).map (fun z => f x + z)
\end{lstlisting}

\textit{Source:} \SourceId{DPCut/Privacy.lean}, line 37.
\begin{lstlisting}[style=lean,firstnumber=37]
abbrev Sequential (Y : Type*) [MeasurableSpace Y] := ℕ → Kernel (List Y) Y
\end{lstlisting}

\textit{Source:} \SourceId{DPCut/Privacy.lean}, line 38.
\begin{lstlisting}[style=lean,firstnumber=38]
def sequentialLaw {Y : Type*} [MeasurableSpace Y] (P : Sequential Y) : ℕ → Measure (List Y)
  | 0 => Measure.dirac []
  | t + 1 => (sequentialLaw P t).bind fun h => (P t h).map (fun y => h ++ [y])
\end{lstlisting}

\textit{Source:} \SourceId{DPCut/Privacy.lean}, line 51.
\begin{lstlisting}[style=lean,firstnumber=51]
def advancedEpsilon (k : ℕ) (ε₀ δ' : ℝ) : ℝ :=
  ε₀ * Real.sqrt (2 * k * Real.log (1 / δ')) + k * ε₀ * (Real.exp ε₀ - 1)
\end{lstlisting}

\textit{Source:} \SourceId{DPCut/Privacy.lean}, line 64.
\begin{lstlisting}[style=lean,firstnumber=64]
def CommonPrefix {Y : Type*} [MeasurableSpace Y] (P₀ P₁ : Sequential Y)
    (t : ℕ) (h : List Y) : Prop :=
  let μ₀ := sequentialLaw P₀ t
  let μ₁ := sequentialLaw P₁ t
  0 < μ₀.rnDeriv (μ₀ + μ₁) h ∧ 0 < μ₁.rnDeriv (μ₀ + μ₁) h
\end{lstlisting}

\textit{Source:} \SourceId{DPCut/Privacy.lean}, line 70.
\begin{lstlisting}[style=lean,firstnumber=70]
def successCharge {Y : Type*} (bottom : Y) (ell : ℕ → List Y → ℝ) (h : List Y) : ℝ :=
  ∑ t ∈ Finset.range h.length, if h[t]?.getD bottom = bottom then 0 else ell t (h.take t)
\end{lstlisting}

\subsection{Algorithm parameter records}
\textit{Source:} \SourceId{DPCut/Algorithms.lean}, line 16.
\begin{lstlisting}[style=lean,firstnumber=16]
structure Constants where
  outer : ℝ
  trials : ℝ
  batch : ℝ
  regularization : ℝ
  certificate : ℝ
  stopping : ℝ
\end{lstlisting}

\textit{Source:} \SourceId{DPCut/Algorithms.lean}, line 24.
\begin{lstlisting}[style=lean,firstnumber=24]
def Constants.Positive (c : Constants) : Prop :=
  0 < c.outer ∧ 0 < c.trials ∧ 0 < c.batch ∧ 0 < c.regularization ∧
  0 < c.certificate ∧ 0 < c.stopping
\end{lstlisting}

\textit{Source:} \SourceId{DPCut/Algorithms.lean}, line 28.
\begin{lstlisting}[style=lean,firstnumber=28]
structure OracleConfig where
  k : ℕ
  K : ℕ
  lam : ℝ
  eps1 : ℝ
  delta1 : ℝ
  beta1 : ℝ
  betaChi : ℝ
  tau1 : ℝ
\end{lstlisting}

\textit{Source:} \SourceId{DPCut/Algorithms.lean}, line 38.
\begin{lstlisting}[style=lean,firstnumber=38]
def oracleConfig (n : ℕ) (c : Constants) (ε δ β : ℝ) : OracleConfig :=
  let K := ⌈c.trials * Real.log (n / β)⌉₊
  let ε₁ := ε / (320 * K * (1 + Real.log (2 / δ)))
  let δ₁ := δ / (20 * (n : ℝ) ^ 2 * K)
  let β₁ := β / (4 * (n : ℝ) ^ 2 * K)
  { k := ⌈c.batch * Real.log (n / (β₁ * δ₁))⌉₊, K := K,
    lam := c.regularization / ε₁ * Real.log (n / (β * δ)),
    eps1 := ε₁, delta1 := δ₁, beta1 := β₁, betaChi := β₁,
    tau1 := Real.log (1 / β₁) / ε₁ }
\end{lstlisting}

\textit{Source:} \SourceId{DPCut/Algorithms.lean}, line 48.
\begin{lstlisting}[style=lean,firstnumber=48]
structure RunConfig where
  T1 : ℕ
  T2 : ℕ
  eps2 : ℝ
  beta2 : ℝ
  tau2 : ℝ
  massCap : ℝ
  epsGE : ℝ
  deltaGE : ℝ
  betaGE : ℝ
  oracle : OracleConfig
  B : ℝ
  stop : ℝ
  eta : ℝ
\end{lstlisting}

\textit{Source:} \SourceId{DPCut/Algorithms.lean}, line 63.
\begin{lstlisting}[style=lean,firstnumber=63]
def runConfig (n : ℕ) (U : ℝ) (p : Parameters) (c : Constants) : RunConfig :=
  let T₁ := outerHorizon n U c.outer
  let T₂ := innerHorizon n p.gamma
  let ε₂ := p.epsilon / (2 * T₁)
  let β₂ := p.beta / (2 * T₁)
  let εG := p.epsilon / (4 * T₁ * Real.sqrt (2 * T₂ * Real.log (2 * T₁ / p.delta)))
  let δG := p.delta / (2 * T₁ * T₂)
  let βG := p.beta / (2 * T₁ * T₂)
  let o := oracleConfig n c εG δG βG
  { T1 := T₁, T2 := T₂, eps2 := ε₂, beta2 := β₂, tau2 := Real.log (1 / β₂) / ε₂,
    massCap := Fintype.card (Edge n) * U + 2 * (Real.log (1 / β₂) / ε₂),
    epsGE := εG, deltaGE := δG, betaGE := βG, oracle := o,
    B := c.certificate * (o.tau1 + o.lam * n * Real.log n + n * U),
    stop := stoppingThreshold n U c.stopping p, eta := stepSize p.gamma }
\end{lstlisting}

\textit{Source:} \SourceId{DPCut/Algorithms.lean}, line 80.
\begin{lstlisting}[style=lean,firstnumber=80]
def clipNoisyMass (cap M z τ : ℝ) : ℝ :=
  min (max 0 cap) (max 0 (M + z + τ))
\end{lstlisting}

\subsection{Concrete privacy constants}
\textit{Source:} \SourceId{DPCut/PublicOracleCalibration.lean}, line 13.
\begin{lstlisting}[style=lean,firstnumber=13]
def privacyConstants : Constants := ⟨1, 8192, 512, 1073741824, 1, 1⟩
\end{lstlisting}

\subsection{Concrete oracle-utility constants}
\textit{Source:} \SourceId{DPCut/UtilityCalibration.lean}, line 10.
\begin{lstlisting}[style=lean,firstnumber=10]
def utilityConstants : Constants := { privacyConstants with certificate := 16 }
\end{lstlisting}

\subsection{Concrete outer-loop constants}
\textit{Source:} \SourceId{DPCut/OuterUtilityCalibration.lean}, line 8.
\begin{lstlisting}[style=lean,firstnumber=8]
def outerUtilityConstants : Constants :=
  { utilityConstants with outer := 80, stopping := 2 ^ 120 }
\end{lstlisting}

\subsection{Algorithm laws referred to by the targets}
\textit{Source:} \SourceId{DPCut/Algorithms.lean}, line 83.
\begin{lstlisting}[style=lean,firstnumber=83]
def activeMaximum {n : ℕ} (r g : Weights n) : ℝ :=
  sSup (insert 0 (g '' {e | 0 < r e}))
\end{lstlisting}

\textit{Source:} \SourceId{DPCut/Algorithms.lean}, line 85.
\begin{lstlisting}[style=lean,firstnumber=85]
def normalizeWitness {n : ℕ} (r g : Weights n) : Weights n :=
  fun e => r e * g e / activeMaximum r g
\end{lstlisting}

\textit{Source:} \SourceId{DPCut/Algorithms.lean}, line 89.
\begin{lstlisting}[style=lean,firstnumber=89]
def findTrials {n : ℕ} (o : OracleConfig) (D r : Weights n) : ℕ → Measure (Option (Weights n))
  | 0 => Measure.dirac none
  | t + 1 => (gaussianBatch (optimizer o.lam D) o.k).bind fun z =>
      (laplace (1 / o.eps1)).bind fun noise =>
        let g := clippedWitness z
        if o.tau1 ≤ pairing D g + noise ∧ 0 < activeMaximum r g
        then Measure.dirac (some (normalizeWitness r g))
        else findTrials o D r t
\end{lstlisting}

\textit{Source:} \SourceId{DPCut/Algorithms.lean}, line 98.
\begin{lstlisting}[style=lean,firstnumber=98]
def findWitness {n : ℕ} (o : OracleConfig) (D r : Weights n) : Measure (Option (Weights n)) :=
  if r = 0 then Measure.dirac none else findTrials o D r o.K
\end{lstlisting}

\textit{Source:} \SourceId{DPCut/Algorithms.lean}, line 101.
\begin{lstlisting}[style=lean,firstnumber=101]
def plusDiscrepancy {n : ℕ} (γ : ℝ) (w h r : Weights n) : Weights n :=
  fun e => r e * (h e - (1 + γ) * w e)
\end{lstlisting}

\textit{Source:} \SourceId{DPCut/Algorithms.lean}, line 103.
\begin{lstlisting}[style=lean,firstnumber=103]
def minusDiscrepancy {n : ℕ} (γ : ℝ) (w h r : Weights n) : Weights n :=
  fun e => r e * ((1 - γ) * w e - h e)
\end{lstlisting}

\textit{Source:} \SourceId{DPCut/Algorithms.lean}, line 108.
\begin{lstlisting}[style=lean,firstnumber=108]
def aggregateLoop {n : ℕ} (o : OracleConfig) (γ : ℝ) (w h : Weights n) :
    ℕ → Aggregate n → Measure (Aggregate n)
  | 0, a => Measure.dirac a
  | t + 1, a =>
      if a.residual = 0 then Measure.dirac a else
      (findWitness o (plusDiscrepancy γ w h a.residual) a.residual).bind fun gp =>
      match gp with
      | some g => aggregateLoop o γ w h t
          { a with plus := a.plus + g, residual := a.residual - g }
      | none =>
        (findWitness o (minusDiscrepancy γ w h a.residual) a.residual).bind fun gm =>
        match gm with
        | some g => aggregateLoop o γ w h t
            { a with minus := a.minus + g, residual := a.residual - g }
        | none => Measure.dirac a
\end{lstlisting}

\textit{Source:} \SourceId{DPCut/Algorithms.lean}, line 124.
\begin{lstlisting}[style=lean,firstnumber=124]
def aggregateLaw {n : ℕ} (o : OracleConfig) (γ : ℝ) (w h : Weights n) : Measure (Aggregate n) :=
  aggregateLoop o γ w h (Fintype.card (Edge n)) { plus := 0, minus := 0, residual := 1 }
\end{lstlisting}

\textit{Source:} \SourceId{DPCut/Algorithms.lean}, line 128.
\begin{lstlisting}[style=lean,firstnumber=128]
def innerLoop {n : ℕ} (p : Parameters) (c : RunConfig) (w : Weights n) (M : ℝ) :
    ℕ → Weights n → Measure (Option (Weights n × Weights n))
  | 0, _ => Measure.dirac none
  | t + 1, h => (aggregateLaw c.oracle p.gamma w h).bind fun a =>
      if a.Accepts h then Measure.dirac (some (h, a.residual)) else
      let over := pairing h a.minus ≤ pairing h a.plus
      let θ : ℝ := if over then 1 else -1
      let g := if over then a.plus else a.minus
      innerLoop p c w M t (mwUpdate M θ c.eta h g)
\end{lstlisting}

\textit{Source:} \SourceId{DPCut/Algorithms.lean}, line 139.
\begin{lstlisting}[style=lean,firstnumber=139]
def outerLoop {n : ℕ} (p : Parameters) (c : RunConfig) (w : Weights n) :
    ℕ → Weights n → Weights n → Measure (Weights n)
  | 0, _, output => Measure.dirac output
  | t + 1, a, output =>
      let wi := hadamard a w
      (laplace (1 / c.eps2)).bind fun noise =>
      let M := clipNoisyMass c.massCap (mass wi) noise c.tau2
      if a = 0 ∨ M ≤ c.stop then Measure.dirac output else
      (innerLoop p c wi M c.T2 (initialEstimate M a)).bind fun result =>
      match result with
      | none => outerLoop p c w t a output
      | some (h, r) => outerLoop p c w t (residualUpdate a r) (output + hadamard h r)
\end{lstlisting}

\textit{Source:} \SourceId{DPCut/Algorithms.lean}, line 152.
\begin{lstlisting}[style=lean,firstnumber=152]
def cutReleaseLaw {n : ℕ} (U : ℝ) (p : Parameters) (c : Constants)
    (w : Weights n) : Measure (Weights n) :=
  let cfg := runConfig n U p c
  outerLoop p cfg w cfg.T1 1 0
\end{lstlisting}

\subsection{Inner success and outer stopping laws}
\textit{Source:} \SourceId{DPCut/Utility.lean}, line 22.
\begin{lstlisting}[style=lean,firstnumber=22]
def innerRoundLaw {n : ℕ} (p : Parameters) (cfg : RunConfig) (w a : Weights n) :
    Measure (ℝ × Option (Weights n × Weights n)) :=
  let wi := hadamard a w
  (laplace (1 / cfg.eps2)).bind fun noise =>
    let M := clipNoisyMass cfg.massCap (mass wi) noise cfg.tau2
    (innerLoop p cfg wi M cfg.T2 (initialEstimate M a)).map (fun result => (M, result))
\end{lstlisting}

\textit{Source:} \SourceId{DPCut/Utility.lean}, line 40.
\begin{lstlisting}[style=lean,firstnumber=40]
def outerStopsLaw {n : ℕ} (p : Parameters) (cfg : RunConfig) (w : Weights n) :
    ℕ → Weights n → Measure Bool
  | 0, _ => Measure.dirac false
  | t + 1, a =>
    let wi := hadamard a w
    (laplace (1 / cfg.eps2)).bind fun noise =>
      let M := clipNoisyMass cfg.massCap (mass wi) noise cfg.tau2
      if a = 0 ∨ M ≤ cfg.stop then Measure.dirac true else
      (innerLoop p cfg wi M cfg.T2 (initialEstimate M a)).bind fun result =>
        match result with
        | none => outerStopsLaw p cfg w t a
        | some (_, r) => outerStopsLaw p cfg w t (residualUpdate a r)
\end{lstlisting}

\subsection{Computational-model interface}
\textit{Source:} \SourceId{DPCut/Statements.lean}, line 190.
\begin{lstlisting}[style=lean,firstnumber=190]
structure ComputationalModel where
  Program : Type
  realizes : ∀ n : ℕ, Program → Parameters → Kernel (Weights n) (Weights n) → Prop
  cost : ∀ n : ℕ, Program → Parameters → Weights n → ℝ≥0∞
\end{lstlisting}

\textit{Source:} \SourceId{DPCut/Statements.lean}, line 195.
\begin{lstlisting}[style=lean,firstnumber=195]
def runtimeScale (n : ℕ) (p : Parameters) : ℝ :=
  1 + n + 1 / p.gamma + 1 / p.epsilon + Real.log (1 / p.delta) + Real.log (1 / p.beta)
\end{lstlisting}

\subsection{Costed laws and real-model block charges}
\textit{Source:} \SourceId{DPCut/CostedProgram.lean}, line 31.
\begin{lstlisting}[style=lean,firstnumber=31]
structure CostedLaw (α : Type) [MeasurableSpace α] where
  law : Measure α
  cost : ℝ≥0∞
\end{lstlisting}

\textit{Source:} \SourceId{DPCut/CostedProgram.lean}, line 63.
\begin{lstlisting}[style=lean,firstnumber=63]
def vectorBlockCost (n : ℕ) : ℝ≥0∞ := 100 * (n + 1) ^ 2
\end{lstlisting}

\textit{Source:} \SourceId{DPCut/CostedProgram.lean}, line 64.
\begin{lstlisting}[style=lean,firstnumber=64]
def gaussianBatchCost (n k : ℕ) : ℝ≥0∞ := vectorBlockCost n * (n + 1) * (k + 1)
\end{lstlisting}

\subsection{Reported completion times and timeout laws}
\textit{Source:} \SourceId{DPCut/TimedAlgorithm.lean}, line 31.
\begin{lstlisting}[style=lean,firstnumber=31]
def timeoutResult (fuel : ℕ) (z : ℕ × α) : Option α :=
  if z.1 ≤ fuel then some z.2 else none
\end{lstlisting}

\textit{Source:} \SourceId{DPCut/TimedAlgorithm.lean}, line 40.
\begin{lstlisting}[style=lean,firstnumber=40]
structure TimedAlgorithm (Θ α : Type) [MeasurableSpace Θ] [MeasurableSpace α] where
  completed : Kernel Θ (ℕ × α)
  completed_markov : IsMarkovKernel completed
  run : ℕ → Θ → CostedLaw (Option α)
  run_law : ∀ fuel θ, (run fuel θ).law =
    (completed θ).map (timeoutResult fuel)
  run_cost : ∀ fuel θ, (run fuel θ).cost ≤ fuel + 1
\end{lstlisting}

\textit{Source:} \SourceId{DPCut/TimedAlgorithm.lean}, line 54. \textit{Namespace:} \SourceId{DPCut.TimedAlgorithm}.
\begin{lstlisting}[style=lean,firstnumber=54]
def expectedCost (θ : Θ) : ℝ≥0∞ := ∫⁻ z, (z.1 : ℝ≥0∞) ∂A.completed θ
\end{lstlisting}

\subsection{Abstract LSW specification: problem, oracle, and guarantee}
\textit{Source:} \SourceId{DPCut/LSWStatement.lean}, line 24. \textit{Namespace:} \SourceId{DPCut.LSW}.
\begin{lstlisting}[style=lean,firstnumber=24]
abbrev Space (ι : Type) [Fintype ι] := EuclideanSpace ℝ ι
\end{lstlisting}

\textit{Source:} \SourceId{DPCut/LSWStatement.lean}, line 28. \textit{Namespace:} \SourceId{DPCut.LSW}.
\begin{lstlisting}[style=lean,firstnumber=28]
structure Problem (ι : Type) [Fintype ι] (Θ : Type) [MeasurableSpace Θ] where
  body : Set (Space ι)
  value : Θ → Space ι → ℝ
  measurable_value : Measurable (Function.uncurry value)
  convex : Convex ℝ body
  compact : IsCompact body
  concave : ∀ θ, ConcaveOn ℝ body (value θ)
  innerRadius : ℝ
  outerRadius : ℝ
  inner_pos : 0 < innerRadius
  radii_order : innerRadius ≤ outerRadius
  inner_ball : closedBall 0 innerRadius ⊆ body
  outer_ball : body ⊆ closedBall 0 outerRadius
  rangeBound : ℝ
  range_one_le : 1 ≤ rangeBound
  range_le : ∀ θ x, x ∈ body → ∀ y, y ∈ body → value θ x - value θ y ≤ rangeBound
\end{lstlisting}

\textit{Source:} \SourceId{DPCut/LSWStatement.lean}, line 47. \textit{Namespace:} \SourceId{DPCut.LSW}.
\begin{lstlisting}[style=lean,firstnumber=47]
abbrev Reply (ι : Type) [Fintype ι] := Space ι ⊕ (ℝ × Space ι)
\end{lstlisting}

\textit{Source:} \SourceId{DPCut/LSWStatement.lean}, line 49. \textit{Namespace:} \SourceId{DPCut.LSW}.
\begin{lstlisting}[style=lean,firstnumber=49]
def CorrectReply {ι : Type} [Fintype ι] {Θ : Type} [MeasurableSpace Θ] (P : Problem ι Θ)
    (θ : Θ) (x : Space ι) : Reply ι → Prop
  | .inl g => x ∉ P.body ∧ g ≠ 0 ∧ ∀ y ∈ P.body, ⟪g, y - x⟫ ≤ 0
  | .inr (v, g) => x ∈ P.body ∧ v = P.value θ x ∧
      ∀ y ∈ P.body, P.value θ y ≤ v + ⟪g, y - x⟫
\end{lstlisting}

\textit{Source:} \SourceId{DPCut/LSWStatement.lean}, line 57. \textit{Namespace:} \SourceId{DPCut.LSW}.
\begin{lstlisting}[style=lean,firstnumber=57]
structure Oracle {ι : Type} [Fintype ι] {Θ : Type} [MeasurableSpace Θ] (P : Problem ι Θ) where
  answer : Θ × Space ι → Arithmetic (Reply ι)
  measurable_answer : Measurable (fun q => (answer q).value)
  correct : ∀ θ x, CorrectReply P θ x (answer (θ, x)).value
  charge : ℕ
  cost_le : ∀ q, (answer q).steps ≤ charge
\end{lstlisting}

\textit{Source:} \SourceId{DPCut/LSWStatement.lean}, line 64. \textit{Namespace:} \SourceId{DPCut.LSW}.
\begin{lstlisting}[style=lean,firstnumber=64]
def Accurate {ι : Type} [Fintype ι] {Θ : Type} [MeasurableSpace Θ] (P : Problem ι Θ)
    (ε : ℝ) (θ : Θ) (x : Space ι) : Prop :=
  x ∈ P.body ∧ ∀ y ∈ P.body, P.value θ y - P.value θ x ≤ ε
\end{lstlisting}

\textit{Source:} \SourceId{DPCut/LSWStatement.lean}, line 70. \textit{Namespace:} \SourceId{DPCut.LSW}.
\begin{lstlisting}[style=lean,firstnumber=70]
def scale {ι : Type} [Fintype ι] {Θ : Type} [MeasurableSpace Θ] (P : Problem ι Θ)
    (O : Oracle P) (ε : ℝ) : ℝ :=
  1 + Fintype.card ι + O.charge + Real.log (P.outerRadius / P.innerRadius) +
    Real.log (P.rangeBound / ε)
\end{lstlisting}

\textit{Source:} \SourceId{DPCut/LSWStatement.lean}, line 77. \textit{Namespace:} \SourceId{DPCut.LSW}.
\begin{lstlisting}[style=lean,firstnumber=77]
def Statement : Prop :=
  ∃ C : ℝ, 0 < C ∧ ∃ degree : ℕ,
    ∀ (ι : Type) (_ : Fintype ι), 1 ≤ Fintype.card ι → ∀ (Θ : Type) (_ : MeasurableSpace Θ)
      (P : Problem ι Θ) (O : Oracle P) (ε : ℝ), 0 < ε → ε ≤ 1 / 2 →
      ∃ A : TimedAlgorithm Θ (Space ι),
        (∀ θ, ∀ᵐ z ∂A.completed θ, Accurate P ε θ z.2) ∧
        ∀ θ, A.expectedCost θ ≤ ENNReal.ofReal (C * scale P O ε ^ degree)
\end{lstlisting}

\textit{Source:} \SourceId{DPCut/LSWStatement.lean}, line 88. \textit{Namespace:} \SourceId{_root_}.
\begin{lstlisting}[style=lean,firstnumber=88]
def DPCut.LSWStatement : Prop := DPCut.LSW.Statement
\end{lstlisting}

\subsection{Linear-algebra interface implemented internally}
\textit{Source:} \SourceId{DPCut/SDPLinearOracle.lean}, line 16.
\begin{lstlisting}[style=lean,firstnumber=16]
structure LinearAlgebraPrimitives where
  psdWitness : ∀ n, Mat n → Option (Fin n → ℝ)
  measurable_psdWitness : ∀ n, Measurable (psdWitness n)
  none_iff : ∀ n (M : Mat n), M.IsHermitian →
    (psdWitness n M = none ↔ M.PosSemidef)
  witness_negative : ∀ n (M : Mat n) v, psdWitness n M = some v →
    v ⬝ᵥ (M *ᵥ v) < 0
\end{lstlisting}

\textit{Source:} \SourceId{DPCut/SDPLinearOracle.lean}, line 27.
\begin{lstlisting}[style=lean,firstnumber=27]
def linearAlgebraPrimitiveCost (n : ℕ) : ℕ := 100 * (n + 1) ^ 3
\end{lstlisting}

\subsection{Value/counter record (not an instruction language)}
\textit{Source:} \SourceId{DPCut/Arithmetic.lean}, line 8.
\begin{lstlisting}[style=lean,firstnumber=8]
structure Arithmetic (α : Type) where
  value : α
  steps : Nat
\end{lstlisting}

\subsection{Constructed linear-algebra implementation certificate}
\textit{Source:} \SourceId{DPCut/VerifiedLinearAlgebra.lean}, line 13.
\begin{lstlisting}[style=lean,firstnumber=13]
structure LinearAlgebraImplementation where
  primitives : LinearAlgebraPrimitives
  psdProgram : ∀ n, Mat n → Arithmetic (Option (Fin n → ℝ))
  psd_value : ∀ n M, (psdProgram n M).value = primitives.psdWitness n M
  psd_cost : ∀ n M, (psdProgram n M).steps ≤ linearAlgebraPrimitiveCost n
  oracleProgram : ∀ n, ℝ → ℝ → Weights n × SDPCoordinates.Space n →
    Arithmetic (LSW.Reply (Fin n × Fin n))
  oracle_value : ∀ n, 2 ≤ n → ∀ lam B q,
    (oracleProgram n lam B q).value = SDPCoordinates.answer primitives lam B q
  oracle_cost : ∀ n lam B q, (oracleProgram n lam B q).steps ≤ SDPCoordinates.oracleCharge n
  decodeProgram : ∀ n, Option (SDPCoordinates.Space n) → Arithmetic (Mat n)
  decode_value : ∀ n (hn : 2 ≤ n) o,
    (decodeProgram n o).value = (SDPCoordinates.decode hn o).val
  decode_cost : ∀ n o, (decodeProgram n o).steps ≤ SDPCoordinates.oracleCharge n
  gaussianProgram : ∀ n, Mat n → (k : ℕ) → CostedLaw (Fin k → Fin n → ℝ)
  gaussian_law : ∀ n X, X.PosDef → ∀ k, (gaussianProgram n X k).law = gaussianBatch X k
  gaussian_cost : ∀ n X k, (gaussianProgram n X k).cost ≤ gaussianBatchCost n k
\end{lstlisting}

\textit{Source:} \SourceId{DPCut/VerifiedLinearAlgebra.lean}, line 33.
\begin{lstlisting}[style=lean,firstnumber=33]
def verifiedLinearAlgebra : LinearAlgebraImplementation where
  primitives := MatrixElimination.primitives
  psdProgram _ := MatrixElimination.Counted.totalPSD
  psd_value _ := MatrixElimination.Counted.totalPSD_value
  psd_cost _ := MatrixElimination.Counted.totalPSD_steps
  oracleProgram _ := SDPCoordinates.Counted.answer
  oracle_value _ := SDPCoordinates.Counted.answer_value
  oracle_cost _ := SDPCoordinates.Counted.answer_steps
  decodeProgram _ := SDPCoordinates.Counted.decode
  decode_value _ := SDPCoordinates.Counted.decode_value
  decode_cost _ := SDPCoordinates.Counted.decode_steps
  gaussianProgram _ := verifiedGaussianBatch
  gaussian_law _ := verifiedGaussianBatch_law
  gaussian_cost _ := verifiedGaussianBatch_cost
\end{lstlisting}

\subsection{Solver family constructed from LSW}
\textit{Source:} \SourceId{DPCut/LSWSolverFamily.lean}, line 13.
\begin{lstlisting}[style=lean,firstnumber=13]
structure CuttingPlaneFamily where
  implementation : LinearAlgebraImplementation
  coefficient : ℝ
  coefficient_pos : 0 < coefficient
  degree : ℕ
  algorithm : ∀ (n : ℕ) (_hn : 2 ≤ n) (lam B ε : ℝ) (_hlam : 0 ≤ lam) (_hB : 0 ≤ B),
    0 < ε → ε ≤ 1/2 → TimedAlgorithm (Weights n) (SDPCoordinates.Space n)
  accurate : ∀ n hn lam B ε hlam hB hε hεhalf D,
    ∀ᵐ z ∂(algorithm n hn lam B ε hlam hB hε hεhalf).completed D,
      LSW.Accurate (SDPCoordinates.problem hn hlam hB) ε D z.2
  expected : ∀ n hn lam B ε hlam hB hε hεhalf D,
    (algorithm n hn lam B ε hlam hB hε hεhalf).expectedCost D ≤
      ENNReal.ofReal (coefficient * sdpLSWScale n lam B ε ^ degree)
\end{lstlisting}

\textit{Source:} \SourceId{DPCut/LSWSolverFamily.lean}, line 35.
\begin{lstlisting}[style=lean,firstnumber=35]
def cuttingPlaneFamily (hLSW : LSWStatement) (I : LinearAlgebraImplementation) : CuttingPlaneFamily :=
  Classical.choose (cuttingPlaneFamily_exists hLSW I)
\end{lstlisting}

\subsection{Public objective bound}
\textit{Source:} \SourceId{DPCut/SDPRuntimeData.lean}, line 13.
\begin{lstlisting}[style=lean,firstnumber=13]
def publicObjectiveBound (n : ℕ) (p : Parameters) : ℝ :=
  (fullRunConfig n p).massCap + (1+(boundedParameters p).gamma) * Fintype.card (Edge n) * preprocessingThreshold n p
\end{lstlisting}

\subsection{Approximation budgets}
\textit{Source:} \SourceId{DPCut/ApproximationTransfer.lean}, line 12.
\begin{lstlisting}[style=lean,firstnumber=12]
def precisionParameters (p : Parameters) : Parameters :=
  { p with delta := p.delta / 2, beta := p.beta / 2 }
\end{lstlisting}

\textit{Source:} \SourceId{DPCut/ApproximationTransfer.lean}, line 21.
\begin{lstlisting}[style=lean,firstnumber=21]
def implementationTV (p : Parameters) : ℝ :=
  min (p.beta / 2) (p.delta / (2 * (1 + Real.exp p.epsilon)))
\end{lstlisting}

\subsection{Public optimizer precision}
\textit{Source:} \SourceId{DPCut/ApproxSDPPrecision.lean}, line 11.
\begin{lstlisting}[style=lean,firstnumber=11]
def approximationBudget (n : ℕ) (p : Parameters) : ℝ :=
  let c := fullRunConfig n (precisionParameters p)
  4 * Fintype.card (Edge n) * c.T1 * c.T2 * c.oracle.K * (c.oracle.k + 1)
\end{lstlisting}

\textit{Source:} \SourceId{DPCut/ApproxSDPPrecision.lean}, line 15.
\begin{lstlisting}[style=lean,firstnumber=15]
def approximationRadius (n : ℕ) (p : Parameters) : ℝ :=
  p.beta * p.delta / (16 * (1 + approximationBudget n p))
\end{lstlisting}

\subsection{Nonnegative-output post-processing}
\textit{Source:} \SourceId{DPCut/ApproxSDPRelease.lean}, line 13.
\begin{lstlisting}[style=lean,firstnumber=13]
def positiveWeights {n : ℕ} (w : Weights n) : Weights n := fun e => max 0 (w e)
\end{lstlisting}

\subsection{Current LSW release program and real-arithmetic model}
\textit{Source:} \SourceId{DPCut/Theorem1LSW.lean}, line 44. \textit{Namespace:} \SourceId{DPCut.LSWRelease}.
\begin{lstlisting}[style=lean,firstnumber=44]
inductive Program where
  | fullRelease
\end{lstlisting}

\textit{Source:} \SourceId{DPCut/Theorem1LSW.lean}, line 47. \textit{Namespace:} \SourceId{DPCut.LSWRelease}.
\begin{lstlisting}[style=lean,firstnumber=47]
def Program.run {n : ℕ} (F : CuttingPlaneFamily) (_ : Program) (p : Parameters) (w : Weights n) :
    CostedLaw (Weights n) :=
  if h : 2 ≤ n ∧ p.Valid then LSWRelease.run F h.1 p h.2 w else .pure 0
\end{lstlisting}

\textit{Source:} \SourceId{DPCut/Theorem1LSW.lean}, line 60. \textit{Namespace:} \SourceId{DPCut.LSWRelease}.
\begin{lstlisting}[style=lean,firstnumber=60]
def model (F : CuttingPlaneFamily) : ComputationalModel where
  Program := Program
  realizes n program p M := ∀ w : Weights n, M w = (program.run F p w).law
  cost _ program p w := (program.run F p w).cost
\end{lstlisting}

\textit{Source:} \SourceId{DPCut/Theorem1LSW.lean}, line 85.
\begin{lstlisting}[style=lean,firstnumber=85]
def realArithmeticLSWImplementationModel (hLSW : LSWStatement) (I : LinearAlgebraImplementation) : ComputationalModel :=
  LSWRelease.model (cuttingPlaneFamily hLSW I)
\end{lstlisting}

\textit{Source:} \SourceId{DPCut/Theorem1LSW.lean}, line 95.
\begin{lstlisting}[style=lean,firstnumber=95]
def realArithmeticLSWModel (hLSW : LSWStatement) : ComputationalModel :=
  realArithmeticLSWImplementationModel hLSW verifiedLinearAlgebra
\end{lstlisting}

\subsection{Legacy exact-SDP contract (historical interface)}
\textit{Source:} \SourceId{DPCut/CostedProgram.lean}, line 21.
\begin{lstlisting}[style=lean,firstnumber=21]
structure ExactSDPOracleCost where
  charge : ℕ → ℝ≥0∞
  coefficient : ℝ
  degree : ℕ
  coefficient_pos : 0 < coefficient
  charge_le : ∀ n, charge n ≤ ENNReal.ofReal (coefficient * (1 + (n : ℝ)) ^ degree)
\end{lstlisting}

\subsection{Finite KL and capped support}
\textit{Source:} \SourceId{DPCut/Definitions.lean}, line 187.
\begin{lstlisting}[style=lean,firstnumber=187]
def discreteMeasure {ι : Type*} [Fintype ι] [MeasurableSpace ι] (p : ι → ℝ) : Measure ι :=
  ∑ i, ENNReal.ofReal (p i) • Measure.dirac i
\end{lstlisting}

\subsection{Finite KL and capped support (continued)}
\textit{Source:} \SourceId{DPCut/Applications.lean}, line 36.
\begin{lstlisting}[style=lean,firstnumber=36]
def CappedOn {n : ℕ} (D : Finset (Edge n)) (w : Weights n) : Prop :=
  Bounded 1 w ∧ ∀ e, e ∉ D → w e = 0
\end{lstlisting}

\textit{Source:} \SourceId{DPCut/Applications.lean}, line 88.
\begin{lstlisting}[style=lean,firstnumber=88]
def IsDistribution {m : ℕ} (p : Fin m → ℝ) : Prop := (∀ i, 0 ≤ p i) ∧ (∑ i, p i) = 1
\end{lstlisting}

\textit{Source:} \SourceId{DPCut/Applications.lean}, line 89.
\begin{lstlisting}[style=lean,firstnumber=89]
def finiteKL {m : ℕ} (p q : Fin m → ℝ) : ℝ≥0∞ :=
  InformationTheory.klDiv (discreteMeasure p) (discreteMeasure q)
\end{lstlisting}

\subsection{Partitions and application objectives}
\textit{Source:} \SourceId{DPCut/Applications.lean}, line 16.
\begin{lstlisting}[style=lean,firstnumber=16]
abbrev Partition (n : ℕ) := Fin n → Fin n
\end{lstlisting}

\textit{Source:} \SourceId{DPCut/Applications.lean}, line 17.
\begin{lstlisting}[style=lean,firstnumber=17]
def partitionQuery {n : ℕ} (π : Partition n) : Weights n :=
  fun e => if π e.val.1 = π e.val.2 then 0 else 1
\end{lstlisting}

\textit{Source:} \SourceId{DPCut/Applications.lean}, line 19.
\begin{lstlisting}[style=lean,firstnumber=19]
def partitionCut {n : ℕ} (w : Weights n) (π : Partition n) : ℝ := pairing w (partitionQuery π)
\end{lstlisting}

\textit{Source:} \SourceId{DPCut/Applications.lean}, line 20.
\begin{lstlisting}[style=lean,firstnumber=20]
def clusterCount {n : ℕ} (π : Partition n) : ℕ := (Finset.univ.image π).card
\end{lstlisting}

\textit{Source:} \SourceId{DPCut/Applications.lean}, line 21.
\begin{lstlisting}[style=lean,firstnumber=21]
def ccCost {n : ℕ} (positive : Weights n) (π : Partition n) : ℝ :=
  ∑ e, if π e.val.1 = π e.val.2 then 1 - positive e else positive e
\end{lstlisting}

\textit{Source:} \SourceId{DPCut/Applications.lean}, line 23.
\begin{lstlisting}[style=lean,firstnumber=23]
def ccOptimum {n : ℕ} (positive : Weights n) : ℝ := sInf (Set.range (ccCost positive))
\end{lstlisting}

\textit{Source:} \SourceId{DPCut/Applications.lean}, line 24.
\begin{lstlisting}[style=lean,firstnumber=24]
def ccOptimumK {n : ℕ} (positive : Weights n) (k : ℕ) : ℝ :=
  sInf (ccCost positive '' {π | clusterCount π ≤ k})
\end{lstlisting}

\textit{Source:} \SourceId{DPCut/Applications.lean}, line 26.
\begin{lstlisting}[style=lean,firstnumber=26]
def maxKOptimum {n : ℕ} (w : Weights n) (k : ℕ) : ℝ :=
  sSup (partitionCut w '' {π | clusterCount π ≤ k})
\end{lstlisting}

\textit{Source:} \SourceId{DPCut/Applications.lean}, line 28.
\begin{lstlisting}[style=lean,firstnumber=28]
def constrainedOptimum {n : ℕ} (w : Weights n) (F : Set (Finset (Fin n))) : ℝ :=
  sInf (cut w '' F)
\end{lstlisting}

\subsection{Statistical application guarantees}
\textit{Source:} \SourceId{DPCut/Applications.lean}, line 146.
\begin{lstlisting}[style=lean,firstnumber=146]
def PrivateCCGuarantee (C Cfail ε ratio : ℝ) (exponent : ℕ) (clusters : ℕ → ℕ) : Prop :=
  ∀ (n : ℕ), 2 ≤ n → ∃ M : Kernel (Weights n) (Partition n), IsMarkovKernel M ∧
    DP (fun w w' => IsUnweighted w ∧ IsUnweighted w' ∧ Adjacent w w') M ε ((n : ℝ)⁻¹ ^ 10) ∧
    ∀ w, IsUnweighted w → HighProbability (M w) (Cfail / (n : ℝ) ^ 2)
      (fun π => clusterCount π ≤ clusters n ∧
        ccCost w π ≤ ratio * ccOptimumK w (clusters n) + C * n * Real.log n ^ exponent)
\end{lstlisting}

\textit{Source:} \SourceId{DPCut/Applications.lean}, line 175.
\begin{lstlisting}[style=lean,firstnumber=175]
def PrivateMaxKGuarantee (C Cfail ε ratio : ℝ) (exponent k : ℕ) : Prop :=
  ∀ (n : ℕ), 2 ≤ n → ∃ M : Kernel (Weights n) (Partition n), IsMarkovKernel M ∧
    EdgeDP M ε ((n : ℝ)⁻¹ ^ 10) ∧ ∀ w, Nonnegative w →
      HighProbability (M w) (Cfail / (n : ℝ) ^ 2) (fun π => clusterCount π ≤ k ∧
        ratio * maxKOptimum w k - C * n * Real.log n ^ exponent ≤ partitionCut w π)
\end{lstlisting}

\subsection{External correlation-clustering algorithm contracts}
\textit{Source:} \SourceId{DPCut/ExternalApplicationContracts.lean}, line 21.
\begin{lstlisting}[style=lean,firstnumber=21]
structure PolynomialGraphAlgorithm (Output : ℕ → Type) [∀ n, MeasurableSpace (Output n)] where
  kernel : (n : ℕ) → Kernel (Weights n) (Output n)
  markov : ∀ n, IsMarkovKernel (kernel n)
  run : (n : ℕ) → Weights n → CostedLaw (Output n)
  realizes : ∀ n w, (run n w).law = kernel n w
  coefficient : ℝ
  degree : ℕ
  coefficient_pos : 0 < coefficient
  cost_le : ∀ n w, (run n w).cost ≤ ENNReal.ofReal (coefficient * (1 + (n : ℝ)) ^ degree)
\end{lstlisting}

\textit{Source:} \SourceId{DPCut/ExternalApplicationContracts.lean}, line 52.
\begin{lstlisting}[style=lean,firstnumber=52]
structure PrivateCCBaselineContract (ε : ℝ) (δ : ℕ → ℝ) extends PolynomialGraphAlgorithm Partition where
  ratio : ℝ
  additiveCoefficient : ℝ
  logDegree : ℕ
  ratio_ge_one : 1 ≤ ratio
  additiveCoefficient_pos : 0 < additiveCoefficient
  privacy : ∀ n, DP (fun w w' => IsUnweighted w ∧ IsUnweighted w' ∧ Adjacent w w')
    (kernel n) ε (δ n)
  approximation : ∀ n, 2 ≤ n → ∀ w, IsUnweighted w →
    HighProbability (kernel n w) ((n : ℝ)⁻¹ ^ 2) (fun π =>
      ccCost w π ≤ ratio * ccOptimum w + additiveCoefficient * n * Real.log n ^ logDegree)
\end{lstlisting}

\textit{Source:} \SourceId{DPCut/ExternalApplicationContracts.lean}, line 32.
\begin{lstlisting}[style=lean,firstnumber=32]
structure CCApproximationContract (ρ p₀ : ℝ) extends PolynomialGraphAlgorithm Partition where
  approximation : ∀ n, 2 ≤ n → ∀ w, IsUnweighted w →
    HighProbability (kernel n w) (1 - p₀) (fun π => ccCost w π ≤ ρ * ccOptimum w)
\end{lstlisting}

\textit{Source:} \SourceId{DPCut/ExternalApplicationContracts.lean}, line 37.
\begin{lstlisting}[style=lean,firstnumber=37]
structure FixedKCCContract (k : ℕ) (slack p₀ : ℝ) extends PolynomialGraphAlgorithm Partition where
  approximation : ∀ n, 2 ≤ n → ∀ w, IsUnweighted w →
    HighProbability (kernel n w) (1 - p₀) (fun π => clusterCount π ≤ k ∧
      ccCost w π ≤ (1 + slack) * ccOptimumK w k)
\end{lstlisting}

\subsection{Private baseline family for correlation clustering}
\textit{Source:} \SourceId{DPCut/CCApplication.lean}, line 12.
\begin{lstlisting}[style=lean,firstnumber=12]
abbrev CCBaselineFamily := ∀ ε : ℝ, 0 < ε →
  PrivateCCBaselineContract (ε / 2) (fun n => (n : ℝ)⁻¹ ^ 10 / 2)
\end{lstlisting}

\subsection{Hierarchical-clustering predicates}
\textit{Source:} \SourceId{DPCut/Applications.lean}, line 40.
\begin{lstlisting}[style=lean,firstnumber=40]
abbrev Hierarchy (n : ℕ) := Finset (Finset (Fin n))
\end{lstlisting}

\textit{Source:} \SourceId{DPCut/Applications.lean}, line 41.
\begin{lstlisting}[style=lean,firstnumber=41]
def IsSplit {n : ℕ} (T : Hierarchy n) (A B C : Finset (Fin n)) : Prop :=
  A ∈ T ∧ B ∈ T ∧ C ∈ T ∧ B.Nonempty ∧ C.Nonempty ∧ Disjoint B C ∧ A = B ∪ C
\end{lstlisting}

\textit{Source:} \SourceId{DPCut/Applications.lean}, line 43.
\begin{lstlisting}[style=lean,firstnumber=43]
def IsHierarchy {n : ℕ} (T : Hierarchy n) : Prop :=
  Finset.univ ∈ T ∧ ∅ ∉ T ∧ (∀ v, {v} ∈ T) ∧
  (∀ A ∈ T, ∀ B ∈ T, Disjoint A B ∨ A ⊆ B ∨ B ⊆ A) ∧
  ∀ A ∈ T, 2 ≤ A.card → ∃ B C, IsSplit T A B C
\end{lstlisting}

\textit{Source:} \SourceId{DPCut/Applications.lean}, line 47.
\begin{lstlisting}[style=lean,firstnumber=47]
def BalancedHierarchy {n : ℕ} (b : ℝ) (T : Hierarchy n) : Prop :=
  IsHierarchy T ∧ ∀ A B C, IsSplit T A B C →
    b * A.card ≤ B.card ∧ b * A.card ≤ C.card
\end{lstlisting}

\textit{Source:} \SourceId{DPCut/Applications.lean}, line 50.
\begin{lstlisting}[style=lean,firstnumber=50]
def lcaSize {n : ℕ} (T : Hierarchy n) (e : Edge n) : ℕ :=
  sInf ((fun A : Finset (Fin n) => A.card) '' {A | A ∈ T ∧ e.val.1 ∈ A ∧ e.val.2 ∈ A})
\end{lstlisting}

\textit{Source:} \SourceId{DPCut/Applications.lean}, line 52.
\begin{lstlisting}[style=lean,firstnumber=52]
def hcCost {n : ℕ} (w : Weights n) (T : Hierarchy n) : ℝ := ∑ e, w e * lcaSize T e
\end{lstlisting}

\textit{Source:} \SourceId{DPCut/Applications.lean}, line 53.
\begin{lstlisting}[style=lean,firstnumber=53]
def hcOptimum {n : ℕ} (w : Weights n) : ℝ := sInf (hcCost w '' {T | IsHierarchy T})
\end{lstlisting}

\textit{Source:} \SourceId{DPCut/Applications.lean}, line 54.
\begin{lstlisting}[style=lean,firstnumber=54]
def inducedExpansion {n : ℕ} (w : Weights n) (A S : Finset (Fin n)) : ℝ :=
  between w S (A \ S) / min (S.card : ℝ) ((A \ S).card : ℝ)
\end{lstlisting}

\textit{Source:} \SourceId{DPCut/Applications.lean}, line 56.
\begin{lstlisting}[style=lean,firstnumber=56]
def BalancedWithin {n : ℕ} (b : ℝ) (A S : Finset (Fin n)) : Prop :=
  S ⊆ A ∧ b * A.card ≤ S.card ∧ b * A.card ≤ (A \ S).card
\end{lstlisting}

\textit{Source:} \SourceId{DPCut/Applications.lean}, line 58.
\begin{lstlisting}[style=lean,firstnumber=58]
def RecursiveApproximation {n : ℕ} (w : Weights n) (b a : ℝ) (T : Hierarchy n) : Prop :=
  ∀ A B C, IsSplit T A B C →
    if 1 / b ≤ (A.card : ℝ) then
      inducedExpansion w A B ≤ a * sInf (inducedExpansion w A '' {R | BalancedWithin b A R})
    else B.card ≤ C.card + 1 ∧ C.card ≤ B.card + 1
\end{lstlisting}

\subsection{Unweighted graphs and degrees}
\textit{Source:} \SourceId{DPCut/Statements.lean}, line 221.
\begin{lstlisting}[style=lean,firstnumber=221]
def IsUnweighted {n : ℕ} (w : Weights n) : Prop := ∀ e, w e = 0 ∨ w e = 1
\end{lstlisting}

\textit{Source:} \SourceId{DPCut/Statements.lean}, line 222.
\begin{lstlisting}[style=lean,firstnumber=222]
def vertexDegree {n : ℕ} (w : Weights n) (v : Fin n) : ℝ :=
  ∑ e, if e.val.1 = v ∨ e.val.2 = v then w e else 0
\end{lstlisting}

\end{document}
